
%% $Id$

\documentclass[12pt,a4paper,fleqn]{report}
\usepackage{latexsym,graphicx}
\usepackage[refpage]{nomencl}
\usepackage{../iman,../extra,../isar,../proof}
\usepackage{Thy/document/isabelle,Thy/document/isabellesym}
\usepackage{style}
\usepackage{../pdfsetup}


\hyphenation{Isabelle}
\hyphenation{Isar}

\isadroptag{theory}
\title{\includegraphics[scale=0.5]{isabelle_isar}
  \\[4ex] The Isabelle/Isar Implementation}
\author{\emph{Makarius Wenzel}}

\makeglossary
\makeindex


\begin{document}

\maketitle 

\begin{abstract}
  We describe the key concepts underlying the Isabelle/Isar
  implementation, including ML references for the most important
  elements.  The aim is to give some insight into the overall system
  architecture, and provide clues on implementing user extensions.
\end{abstract}

\vspace*{2.5cm}
\begin{quote}
  
  {\small\em Isabelle was not designed; it evolved.  Not everyone
    likes this idea.  Specification experts rightly abhor
    trial-and-error programming.  They suggest that no one should
    write a program without first writing a complete formal
    specification. But university departments are not software houses.
    Programs like Isabelle are not products: when they have served
    their purpose, they are discarded.}
  
  Lawrence C. Paulson, ``Isabelle: The Next 700 Theorem Provers''

  \vspace*{1cm}
  
  {\small\em As I did 20 years ago, I still fervently believe that the
    only way to make software secure, reliable, and fast is to make it
    small.  Fight features.}
  
  Andrew S. Tanenbaum

\end{quote}

\thispagestyle{empty}\clearpage

\pagenumbering{roman} \tableofcontents \clearfirst


\chapter{Introduction}

\section{Quick start}

FIXME examples, ProofGeneral setup

\section{Examples}

\section{How to write Isar proofs anyway?}


%%% Local Variables: 
%%% mode: latex
%%% TeX-master: "isar-ref"
%%% End: 

%
\begin{isabellebody}%
\def\isabellecontext{prelim}%
%
\isadelimtheory
\isanewline
\isanewline
\isanewline
%
\endisadelimtheory
%
\isatagtheory
\isacommand{theory}\isamarkupfalse%
\ prelim\ \isakeyword{imports}\ base\ \isakeyword{begin}%
\endisatagtheory
{\isafoldtheory}%
%
\isadelimtheory
%
\endisadelimtheory
%
\isamarkupchapter{Preliminaries%
}
\isamarkuptrue%
%
\isamarkupsection{Contexts \label{sec:context}%
}
\isamarkuptrue%
%
\begin{isamarkuptext}%
A logical context represents the background that is required for
  formulating statements and composing proofs.  It acts as a medium to
  produce formal content, depending on earlier material (declarations,
  results etc.).

  For example, derivations within the Isabelle/Pure logic can be
  described as a judgment \isa{{\isasymGamma}\ {\isasymturnstile}\isactrlsub {\isasymTheta}\ {\isasymphi}}, which means that a
  proposition \isa{{\isasymphi}} is derivable from hypotheses \isa{{\isasymGamma}}
  within the theory \isa{{\isasymTheta}}.  There are logical reasons for
  keeping \isa{{\isasymTheta}} and \isa{{\isasymGamma}} separate: theories can be
  liberal about supporting type constructors and schematic
  polymorphism of constants and axioms, while the inner calculus of
  \isa{{\isasymGamma}\ {\isasymturnstile}\ {\isasymphi}} is strictly limited to Simple Type Theory (with
  fixed type variables in the assumptions).

  \medskip Contexts and derivations are linked by the following key
  principles:

  \begin{itemize}

  \item Transfer: monotonicity of derivations admits results to be
  transferred into a \emph{larger} context, i.e.\ \isa{{\isasymGamma}\ {\isasymturnstile}\isactrlsub {\isasymTheta}\ {\isasymphi}} implies \isa{{\isasymGamma}{\isacharprime}\ {\isasymturnstile}\isactrlsub {\isasymTheta}\isactrlsub {\isacharprime}\ {\isasymphi}} for contexts \isa{{\isasymTheta}{\isacharprime}\ {\isasymsupseteq}\ {\isasymTheta}} and \isa{{\isasymGamma}{\isacharprime}\ {\isasymsupseteq}\ {\isasymGamma}}.

  \item Export: discharge of hypotheses admits results to be exported
  into a \emph{smaller} context, i.e.\ \isa{{\isasymGamma}{\isacharprime}\ {\isasymturnstile}\isactrlsub {\isasymTheta}\ {\isasymphi}}
  implies \isa{{\isasymGamma}\ {\isasymturnstile}\isactrlsub {\isasymTheta}\ {\isasymDelta}\ {\isasymLongrightarrow}\ {\isasymphi}} where \isa{{\isasymGamma}{\isacharprime}\ {\isasymsupseteq}\ {\isasymGamma}} and
  \isa{{\isasymDelta}\ {\isacharequal}\ {\isasymGamma}{\isacharprime}\ {\isacharminus}\ {\isasymGamma}}.  Note that \isa{{\isasymTheta}} remains unchanged here,
  only the \isa{{\isasymGamma}} part is affected.

  \end{itemize}

  \medskip By modeling the main characteristics of the primitive
  \isa{{\isasymTheta}} and \isa{{\isasymGamma}} above, and abstracting over any
  particular logical content, we arrive at the fundamental notions of
  \emph{theory context} and \emph{proof context} in Isabelle/Isar.
  These implement a certain policy to manage arbitrary \emph{context
  data}.  There is a strongly-typed mechanism to declare new kinds of
  data at compile time.

  The internal bootstrap process of Isabelle/Pure eventually reaches a
  stage where certain data slots provide the logical content of \isa{{\isasymTheta}} and \isa{{\isasymGamma}} sketched above, but this does not stop there!
  Various additional data slots support all kinds of mechanisms that
  are not necessarily part of the core logic.

  For example, there would be data for canonical introduction and
  elimination rules for arbitrary operators (depending on the
  object-logic and application), which enables users to perform
  standard proof steps implicitly (cf.\ the \isa{rule} method
  \cite{isabelle-isar-ref}).

  \medskip Thus Isabelle/Isar is able to bring forth more and more
  concepts successively.  In particular, an object-logic like
  Isabelle/HOL continues the Isabelle/Pure setup by adding specific
  components for automated reasoning (classical reasoner, tableau
  prover, structured induction etc.) and derived specification
  mechanisms (inductive predicates, recursive functions etc.).  All of
  this is ultimately based on the generic data management by theory
  and proof contexts introduced here.%
\end{isamarkuptext}%
\isamarkuptrue%
%
\isamarkupsubsection{Theory context \label{sec:context-theory}%
}
\isamarkuptrue%
%
\begin{isamarkuptext}%
\glossary{Theory}{FIXME}

  A \emph{theory} is a data container with explicit named and unique
  identifier.  Theories are related by a (nominal) sub-theory
  relation, which corresponds to the dependency graph of the original
  construction; each theory is derived from a certain sub-graph of
  ancestor theories.

  The \isa{merge} operation produces the least upper bound of two
  theories, which actually degenerates into absorption of one theory
  into the other (due to the nominal sub-theory relation).

  The \isa{begin} operation starts a new theory by importing
  several parent theories and entering a special \isa{draft} mode,
  which is sustained until the final \isa{end} operation.  A draft
  theory acts like a linear type, where updates invalidate earlier
  versions.  An invalidated draft is called ``stale''.

  The \isa{checkpoint} operation produces an intermediate stepping
  stone that will survive the next update: both the original and the
  changed theory remain valid and are related by the sub-theory
  relation.  Checkpointing essentially recovers purely functional
  theory values, at the expense of some extra internal bookkeeping.

  The \isa{copy} operation produces an auxiliary version that has
  the same data content, but is unrelated to the original: updates of
  the copy do not affect the original, neither does the sub-theory
  relation hold.

  \medskip The example in \figref{fig:ex-theory} below shows a theory
  graph derived from \isa{Pure}, with theory \isa{Length}
  importing \isa{Nat} and \isa{List}.  The body of \isa{Length} consists of a sequence of updates, working mostly on
  drafts.  Intermediate checkpoints may occur as well, due to the
  history mechanism provided by the Isar top-level, cf.\
  \secref{sec:isar-toplevel}.

  \begin{figure}[htb]
  \begin{center}
  \begin{tabular}{rcccl}
        &            & \isa{Pure} \\
        &            & \isa{{\isasymdown}} \\
        &            & \isa{FOL} \\
        & $\swarrow$ &              & $\searrow$ & \\
  $Nat$ &            &              &            & \isa{List} \\
        & $\searrow$ &              & $\swarrow$ \\
        &            & \isa{Length} \\
        &            & \multicolumn{3}{l}{~~$\isarkeyword{imports}$} \\
        &            & \multicolumn{3}{l}{~~$\isarkeyword{begin}$} \\
        &            & $\vdots$~~ \\
        &            & \isa{{\isasymbullet}}~~ \\
        &            & $\vdots$~~ \\
        &            & \isa{{\isasymbullet}}~~ \\
        &            & $\vdots$~~ \\
        &            & \multicolumn{3}{l}{~~$\isarkeyword{end}$} \\
  \end{tabular}
  \caption{A theory definition depending on ancestors}\label{fig:ex-theory}
  \end{center}
  \end{figure}

  \medskip There is a separate notion of \emph{theory reference} for
  maintaining a live link to an evolving theory context: updates on
  drafts are propagated automatically.  The dynamic stops after an
  explicit \isa{end} only.

  Derived entities may store a theory reference in order to indicate
  the context they belong to.  This implicitly assumes monotonic
  reasoning, because the referenced context may become larger without
  further notice.%
\end{isamarkuptext}%
\isamarkuptrue%
%
\isadelimmlref
%
\endisadelimmlref
%
\isatagmlref
%
\begin{isamarkuptext}%
\begin{mldecls}
  \indexmltype{theory}\verb|type theory| \\
  \indexml{Theory.subthy}\verb|Theory.subthy: theory * theory -> bool| \\
  \indexml{Theory.merge}\verb|Theory.merge: theory * theory -> theory| \\
  \indexml{Theory.checkpoint}\verb|Theory.checkpoint: theory -> theory| \\
  \indexml{Theory.copy}\verb|Theory.copy: theory -> theory| \\[1ex]
  \indexmltype{theory-ref}\verb|type theory_ref| \\
  \indexml{Theory.self-ref}\verb|Theory.self_ref: theory -> theory_ref| \\
  \indexml{Theory.deref}\verb|Theory.deref: theory_ref -> theory| \\
  \end{mldecls}

  \begin{description}

  \item \verb|theory| represents theory contexts.  This is
  essentially a linear type!  Most operations destroy the original
  version, which then becomes ``stale''.

  \item \verb|Theory.subthy|~\isa{{\isacharparenleft}thy\isactrlsub {\isadigit{1}}{\isacharcomma}\ thy\isactrlsub {\isadigit{2}}{\isacharparenright}}
  compares theories according to the inherent graph structure of the
  construction.  This sub-theory relation is a nominal approximation
  of inclusion (\isa{{\isasymsubseteq}}) of the corresponding content.

  \item \verb|Theory.merge|~\isa{{\isacharparenleft}thy\isactrlsub {\isadigit{1}}{\isacharcomma}\ thy\isactrlsub {\isadigit{2}}{\isacharparenright}}
  absorbs one theory into the other.  This fails for unrelated
  theories!

  \item \verb|Theory.checkpoint|~\isa{thy} produces a safe
  stepping stone in the linear development of \isa{thy}.  The next
  update will result in two related, valid theories.

  \item \verb|Theory.copy|~\isa{thy} produces a variant of \isa{thy} that holds a copy of the same data.  The result is not
  related to the original; the original is unchanched.

  \item \verb|theory_ref| represents a sliding reference to an
  always valid theory; updates on the original are propagated
  automatically.

  \item \verb|Theory.self_ref|~\isa{thy} and \verb|Theory.deref|~\isa{thy{\isacharunderscore}ref} convert between \verb|theory| and \verb|theory_ref|.  As the referenced theory
  evolves monotonically over time, later invocations of \verb|Theory.deref| may refer to a larger context.

  \end{description}%
\end{isamarkuptext}%
\isamarkuptrue%
%
\endisatagmlref
{\isafoldmlref}%
%
\isadelimmlref
%
\endisadelimmlref
%
\isamarkupsubsection{Proof context \label{sec:context-proof}%
}
\isamarkuptrue%
%
\begin{isamarkuptext}%
\glossary{Proof context}{The static context of a structured proof,
  acts like a local ``theory'' of the current portion of Isar proof
  text, generalizes the idea of local hypotheses \isa{{\isasymGamma}} in
  judgments \isa{{\isasymGamma}\ {\isasymturnstile}\ {\isasymphi}} of natural deduction calculi.  There is a
  generic notion of introducing and discharging hypotheses.
  Arbritrary auxiliary context data may be adjoined.}

  A proof context is a container for pure data with a back-reference
  to the theory it belongs to.  The \isa{init} operation creates a
  proof context from a given theory.  Modifications to draft theories
  are propagated to the proof context as usual, but there is also an
  explicit \isa{transfer} operation to force resynchronization
  with more substantial updates to the underlying theory.  The actual
  context data does not require any special bookkeeping, thanks to the
  lack of destructive features.

  Entities derived in a proof context need to record inherent logical
  requirements explicitly, since there is no separate context
  identification as for theories.  For example, hypotheses used in
  primitive derivations (cf.\ \secref{sec:thms}) are recorded
  separately within the sequent \isa{{\isasymGamma}\ {\isasymturnstile}\ {\isasymphi}}, just to make double
  sure.  Results could still leak into an alien proof context do to
  programming errors, but Isabelle/Isar includes some extra validity
  checks in critical positions, notably at the end of sub-proof.

  Proof contexts may be manipulated arbitrarily, although the common
  discipline is to follow block structure as a mental model: a given
  context is extended consecutively, and results are exported back
  into the original context.  Note that the Isar proof states model
  block-structured reasoning explicitly, using a stack of proof
  contexts internally, cf.\ \secref{sec:isar-proof-state}.%
\end{isamarkuptext}%
\isamarkuptrue%
%
\isadelimmlref
%
\endisadelimmlref
%
\isatagmlref
%
\begin{isamarkuptext}%
\begin{mldecls}
  \indexmltype{Proof.context}\verb|type Proof.context| \\
  \indexml{ProofContext.init}\verb|ProofContext.init: theory -> Proof.context| \\
  \indexml{ProofContext.theory-of}\verb|ProofContext.theory_of: Proof.context -> theory| \\
  \indexml{ProofContext.transfer}\verb|ProofContext.transfer: theory -> Proof.context -> Proof.context| \\
  \end{mldecls}

  \begin{description}

  \item \verb|Proof.context| represents proof contexts.  Elements
  of this type are essentially pure values, with a sliding reference
  to the background theory.

  \item \verb|ProofContext.init|~\isa{thy} produces a proof context
  derived from \isa{thy}, initializing all data.

  \item \verb|ProofContext.theory_of|~\isa{ctxt} selects the
  background theory from \isa{ctxt}, dereferencing its internal
  \verb|theory_ref|.

  \item \verb|ProofContext.transfer|~\isa{thy\ ctxt} promotes the
  background theory of \isa{ctxt} to the super theory \isa{thy}.

  \end{description}%
\end{isamarkuptext}%
\isamarkuptrue%
%
\endisatagmlref
{\isafoldmlref}%
%
\isadelimmlref
%
\endisadelimmlref
%
\isamarkupsubsection{Generic contexts \label{sec:generic-context}%
}
\isamarkuptrue%
%
\begin{isamarkuptext}%
A generic context is the disjoint sum of either a theory or proof
  context.  Occasionally, this enables uniform treatment of generic
  context data, typically extra-logical information.  Operations on
  generic contexts include the usual injections, partial selections,
  and combinators for lifting operations on either component of the
  disjoint sum.

  Moreover, there are total operations \isa{theory{\isacharunderscore}of} and \isa{proof{\isacharunderscore}of} to convert a generic context into either kind: a theory
  can always be selected from the sum, while a proof context might
  have to be constructed by an ad-hoc \isa{init} operation.%
\end{isamarkuptext}%
\isamarkuptrue%
%
\isadelimmlref
%
\endisadelimmlref
%
\isatagmlref
%
\begin{isamarkuptext}%
\begin{mldecls}
  \indexmltype{Context.generic}\verb|type Context.generic| \\
  \indexml{Context.theory-of}\verb|Context.theory_of: Context.generic -> theory| \\
  \indexml{Context.proof-of}\verb|Context.proof_of: Context.generic -> Proof.context| \\
  \end{mldecls}

  \begin{description}

  \item \verb|Context.generic| is the direct sum of \verb|theory| and \verb|Proof.context|, with the datatype
  constructors \verb|Context.Theory| and \verb|Context.Proof|.

  \item \verb|Context.theory_of|~\isa{context} always produces a
  theory from the generic \isa{context}, using \verb|ProofContext.theory_of| as required.

  \item \verb|Context.proof_of|~\isa{context} always produces a
  proof context from the generic \isa{context}, using \verb|ProofContext.init| as required (note that this re-initializes the
  context data with each invocation).

  \end{description}%
\end{isamarkuptext}%
\isamarkuptrue%
%
\endisatagmlref
{\isafoldmlref}%
%
\isadelimmlref
%
\endisadelimmlref
%
\isamarkupsubsection{Context data \label{sec:context-data}%
}
\isamarkuptrue%
%
\begin{isamarkuptext}%
The main purpose of theory and proof contexts is to manage arbitrary
  data.  New data types can be declared incrementally at compile time.
  There are separate declaration mechanisms for any of the three kinds
  of contexts: theory, proof, generic.

  \paragraph{Theory data} may refer to destructive entities, which are
  maintained in direct correspondence to the linear evolution of
  theory values, including explicit copies.\footnote{Most existing
  instances of destructive theory data are merely historical relics
  (e.g.\ the destructive theorem storage, and destructive hints for
  the Simplifier and Classical rules).}  A theory data declaration
  needs to implement the following specification (depending on type
  \isa{T}):

  \medskip
  \begin{tabular}{ll}
  \isa{name{\isacharcolon}\ string} \\
  \isa{empty{\isacharcolon}\ T} & initial value \\
  \isa{copy{\isacharcolon}\ T\ {\isasymrightarrow}\ T} & refresh impure data \\
  \isa{extend{\isacharcolon}\ T\ {\isasymrightarrow}\ T} & re-initialize on import \\
  \isa{merge{\isacharcolon}\ T\ {\isasymtimes}\ T\ {\isasymrightarrow}\ T} & join on import \\
  \isa{print{\isacharcolon}\ T\ {\isasymrightarrow}\ unit} & diagnostic output \\
  \end{tabular}
  \medskip

  \noindent The \isa{name} acts as a comment for diagnostic
  messages; \isa{copy} is just the identity for pure data; \isa{extend} is acts like a unitary version of \isa{merge}, both
  should also include the functionality of \isa{copy} for impure
  data.

  \paragraph{Proof context data} is purely functional.  A declaration
  needs to implement the following specification:

  \medskip
  \begin{tabular}{ll}
  \isa{name{\isacharcolon}\ string} \\
  \isa{init{\isacharcolon}\ theory\ {\isasymrightarrow}\ T} & produce initial value \\
  \isa{print{\isacharcolon}\ T\ {\isasymrightarrow}\ unit} & diagnostic output \\
  \end{tabular}
  \medskip

  \noindent The \isa{init} operation is supposed to produce a pure
  value from the given background theory.  The remainder is analogous
  to theory data.

  \paragraph{Generic data} provides a hybrid interface for both theory
  and proof data.  The declaration is essentially the same as for
  (pure) theory data, without \isa{copy}, though.  The \isa{init} operation for proof contexts merely selects the current data
  value from the background theory.

  \bigskip In any case, a data declaration of type \isa{T} results
  in the following interface:

  \medskip
  \begin{tabular}{ll}
  \isa{init{\isacharcolon}\ theory\ {\isasymrightarrow}\ theory} \\
  \isa{get{\isacharcolon}\ context\ {\isasymrightarrow}\ T} \\
  \isa{put{\isacharcolon}\ T\ {\isasymrightarrow}\ context\ {\isasymrightarrow}\ context} \\
  \isa{map{\isacharcolon}\ {\isacharparenleft}T\ {\isasymrightarrow}\ T{\isacharparenright}\ {\isasymrightarrow}\ context\ {\isasymrightarrow}\ context} \\
  \isa{print{\isacharcolon}\ context\ {\isasymrightarrow}\ unit}
  \end{tabular}
  \medskip

  \noindent Here \isa{init} needs to be applied to the current
  theory context once, in order to register the initial setup.  The
  other operations provide access for the particular kind of context
  (theory, proof, or generic context).  Note that this is a safe
  interface: there is no other way to access the corresponding data
  slot of a context.  By keeping these operations private, a component
  may maintain abstract values authentically, without other components
  interfering.%
\end{isamarkuptext}%
\isamarkuptrue%
%
\isadelimmlref
%
\endisadelimmlref
%
\isatagmlref
%
\begin{isamarkuptext}%
\begin{mldecls}
  \indexmlfunctor{TheoryDataFun}\verb|functor TheoryDataFun| \\
  \indexmlfunctor{ProofDataFun}\verb|functor ProofDataFun| \\
  \indexmlfunctor{GenericDataFun}\verb|functor GenericDataFun| \\
  \end{mldecls}

  \begin{description}

  \item \verb|TheoryDataFun|\isa{{\isacharparenleft}spec{\isacharparenright}} declares data for
  type \verb|theory| according to the specification provided as
  argument structure.  The resulting structure provides data init and
  access operations as described above.

  \item \verb|ProofDataFun|\isa{{\isacharparenleft}spec{\isacharparenright}} is analogous to
  \verb|TheoryDataFun| for type \verb|Proof.context|.

  \item \verb|GenericDataFun|\isa{{\isacharparenleft}spec{\isacharparenright}} is analogous to
  \verb|TheoryDataFun| for type \verb|Context.generic|.

  \end{description}%
\end{isamarkuptext}%
\isamarkuptrue%
%
\endisatagmlref
{\isafoldmlref}%
%
\isadelimmlref
%
\endisadelimmlref
%
\isamarkupsection{Names%
}
\isamarkuptrue%
%
\begin{isamarkuptext}%
In principle, a name is just a string, but there are various
  convention for encoding additional structure.

  For example, the string ``\isa{Foo{\isachardot}bar{\isachardot}baz}'' is considered as a
  qualified name.  The most basic constituents of names may have their
  own structure, e.g.\ the string ``\verb,\,\verb,<alpha>,'' is
  considered as a single symbol (printed as ``\isa{{\isasymalpha}}'').%
\end{isamarkuptext}%
\isamarkuptrue%
%
\isamarkupsubsection{Strings of symbols%
}
\isamarkuptrue%
%
\begin{isamarkuptext}%
\glossary{Symbol}{The smallest unit of text in Isabelle, subsumes
  plain ASCII characters as well as an infinite collection of named
  symbols (for greek, math etc.).}

  A \emph{symbol} constitutes the smallest textual unit in Isabelle
  --- raw characters are normally not encountered.  Isabelle strings
  consist of a sequence of symbols, represented as a packed string or
  a list of symbols.  Each symbol is in itself a small string, which
  is of one of the following forms:

  \begin{enumerate}

  \item singleton ASCII character ``\isa{c}'' (character code
  0--127), for example ``\verb,a,'',

  \item regular symbol ``\verb,\,\verb,<,\isa{ident}\verb,>,'',
  for example ``\verb,\,\verb,<alpha>,'',

  \item control symbol ``\verb,\,\verb,<^,\isa{ident}\verb,>,'',
  for example ``\verb,\,\verb,<^bold>,'',

  \item raw symbol ``\verb,\,\verb,<^raw:,\isa{text}\verb,>,'' where
  \isa{text} is constists of printable characters excluding
  ``\verb,.,'' and ``\verb,>,'', for example
  ``\verb,\,\verb,<^raw:$\sum_{i = 1}^n$>,'',

  \item numbered raw control symbol ``\verb,\,\verb,<^raw,\isa{n}\verb,>, where \isa{n} consists of digits, for example
  ``\verb,\,\verb,<^raw42>,''.

  \end{enumerate}

  \noindent The \isa{ident} syntax for symbol names is \isa{letter\ {\isacharparenleft}letter\ {\isacharbar}\ digit{\isacharparenright}\isactrlsup {\isacharasterisk}}, where \isa{letter\ {\isacharequal}\ A{\isachardot}{\isachardot}Za{\isachardot}{\isachardot}z} and \isa{digit\ {\isacharequal}\ {\isadigit{0}}{\isachardot}{\isachardot}{\isadigit{9}}}.  There are infinitely many
  regular symbols and control symbols, but a fixed collection of
  standard symbols is treated specifically.  For example,
  ``\verb,\,\verb,<alpha>,'' is classified as a (non-ASCII) letter,
  which means it may occur within regular Isabelle identifier syntax.

  Note that the character set underlying Isabelle symbols is plain
  7-bit ASCII.  Since 8-bit characters are passed through
  transparently, Isabelle may process Unicode/UCS data (in UTF-8
  encoding) as well.  Unicode provides its own collection of
  mathematical symbols, but there is no built-in link to the ones of
  Isabelle.

  \medskip Output of Isabelle symbols depends on the print mode
  (\secref{FIXME}).  For example, the standard {\LaTeX} setup of the
  Isabelle document preparation system would present
  ``\verb,\,\verb,<alpha>,'' as \isa{{\isasymalpha}}, and
  ``\verb,\,\verb,<^bold>,\verb,\,\verb,<alpha>,'' as \isa{\isactrlbold {\isasymalpha}}.%
\end{isamarkuptext}%
\isamarkuptrue%
%
\isadelimmlref
%
\endisadelimmlref
%
\isatagmlref
%
\begin{isamarkuptext}%
\begin{mldecls}
  \indexmltype{Symbol.symbol}\verb|type Symbol.symbol| \\
  \indexml{Symbol.explode}\verb|Symbol.explode: string -> Symbol.symbol list| \\
  \indexml{Symbol.is-letter}\verb|Symbol.is_letter: Symbol.symbol -> bool| \\
  \indexml{Symbol.is-digit}\verb|Symbol.is_digit: Symbol.symbol -> bool| \\
  \indexml{Symbol.is-quasi}\verb|Symbol.is_quasi: Symbol.symbol -> bool| \\
  \indexml{Symbol.is-blank}\verb|Symbol.is_blank: Symbol.symbol -> bool| \\[1ex]
  \indexmltype{Symbol.sym}\verb|type Symbol.sym| \\
  \indexml{Symbol.decode}\verb|Symbol.decode: Symbol.symbol -> Symbol.sym| \\
  \end{mldecls}

  \begin{description}

  \item \verb|Symbol.symbol| represents Isabelle symbols.  This
  type is an alias for \verb|string|, but emphasizes the
  specific format encountered here.

  \item \verb|Symbol.explode|~\isa{str} produces a symbol list
  from the packed form that.  This function supercedes \verb|String.explode| for virtually all purposes of manipulating text in
  Isabelle!

  \item \verb|Symbol.is_letter|, \verb|Symbol.is_digit|, \verb|Symbol.is_quasi|, \verb|Symbol.is_blank| classify standard
  symbols according to fixed syntactic conventions of Isabelle, cf.\
  \cite{isabelle-isar-ref}.

  \item \verb|Symbol.sym| is a concrete datatype that represents
  the different kinds of symbols explicitly with constructors \verb|Symbol.Char|, \verb|Symbol.Sym|, \verb|Symbol.Ctrl|, or \verb|Symbol.Raw|.

  \item \verb|Symbol.decode| converts the string representation of a
  symbol into the datatype version.

  \end{description}%
\end{isamarkuptext}%
\isamarkuptrue%
%
\endisatagmlref
{\isafoldmlref}%
%
\isadelimmlref
%
\endisadelimmlref
%
\isamarkupsubsection{Basic names \label{sec:basic-names}%
}
\isamarkuptrue%
%
\begin{isamarkuptext}%
A \emph{basic name} essentially consists of a single Isabelle
  identifier.  There are conventions to mark separate classes of basic
  names, by attaching a suffix of underscores (\isa{{\isacharunderscore}}): one
  underscore means \emph{internal name}, two underscores means
  \emph{Skolem name}, three underscores means \emph{internal Skolem
  name}.

  For example, the basic name \isa{foo} has the internal version
  \isa{foo{\isacharunderscore}}, with Skolem versions \isa{foo{\isacharunderscore}{\isacharunderscore}} and \isa{foo{\isacharunderscore}{\isacharunderscore}{\isacharunderscore}}, respectively.

  Such special versions are required for bookkeeping of names that are
  apart from anything that may appear in the text given by the user.
  In particular, system generated variables in high-level Isar proof
  contexts are usually marked as internal, which prevents mysterious
  name references such as \isa{xaa} in the text.

  \medskip Basic manipulations of binding scopes requires names to be
  modified.  A \emph{name context} contains a collection of already
  used names, which is maintained by the \isa{declare} operation.

  The \isa{invents} operation derives a number of fresh names
  derived from a given starting point.  For example, three names
  derived from \isa{a} are \isa{a}, \isa{b}, \isa{c},
  provided there are no clashes with already used names.

  The \isa{variants} operation produces fresh names by
  incrementing given names as to base-26 numbers (with digits \isa{a{\isachardot}{\isachardot}z}).  For example, name \isa{foo} results in variants
  \isa{fooa}, \isa{foob}, \isa{fooc}, \dots, \isa{fooaa}, \isa{fooab}, \dots; each renaming step picks the next
  unused variant from this list.%
\end{isamarkuptext}%
\isamarkuptrue%
%
\isadelimmlref
%
\endisadelimmlref
%
\isatagmlref
%
\begin{isamarkuptext}%
\begin{mldecls}
  \indexml{Name.internal}\verb|Name.internal: string -> string| \\
  \indexml{Name.skolem}\verb|Name.skolem: string -> string| \\[1ex]
  \indexmltype{Name.context}\verb|type Name.context| \\
  \indexml{Name.context}\verb|Name.context: Name.context| \\
  \indexml{Name.declare}\verb|Name.declare: string -> Name.context -> Name.context| \\
  \indexml{Name.invents}\verb|Name.invents: Name.context -> string -> int -> string list| \\
  \indexml{Name.variants}\verb|Name.variants: string list -> Name.context -> string list * Name.context| \\
  \end{mldecls}

  \begin{description}

  \item \verb|Name.internal|~\isa{name} produces an internal name
  by adding one underscore.

  \item \verb|Name.skolem|~\isa{name} produces a Skolem name by
  adding two underscores.

  \item \verb|Name.context| represents the context of already used
  names; the initial value is \verb|Name.context|.

  \item \verb|Name.declare|~\isa{name} declares \isa{name} as
  being used.

  \item \verb|Name.invents|~\isa{context\ base\ n} produces \isa{n} fresh names derived from \isa{base}.

  \end{description}%
\end{isamarkuptext}%
\isamarkuptrue%
%
\endisatagmlref
{\isafoldmlref}%
%
\isadelimmlref
%
\endisadelimmlref
%
\isamarkupsubsection{Indexed names%
}
\isamarkuptrue%
%
\begin{isamarkuptext}%
An \emph{indexed name} (or \isa{indexname}) is a pair of a basic
  name with a natural number.  This representation allows efficient
  renaming by incrementing the second component only.  To rename two
  collections of indexnames apart from each other, first determine the
  maximum index \isa{maxidx} of the first collection, then
  increment all indexes of the second collection by \isa{maxidx\ {\isacharplus}\ {\isadigit{1}}}.  Note that the maximum index of an empty collection is \isa{{\isacharminus}{\isadigit{1}}}.

  Isabelle syntax observes the following rules for representing an
  indexname \isa{{\isacharparenleft}x{\isacharcomma}\ i{\isacharparenright}} as a packed string:

  \begin{itemize}

  \item \isa{{\isacharquery}x} if \isa{x} does not end with a digit and \isa{i\ {\isacharequal}\ {\isadigit{0}}},

  \item \isa{{\isacharquery}xi} if \isa{x} does not end with a digit,

  \item \isa{{\isacharquery}x{\isachardot}i} else.

  \end{itemize}

  Occasionally, basic names and indexed names are injected into the
  same pair type: the (improper) indexname \isa{{\isacharparenleft}x{\isacharcomma}\ {\isacharminus}{\isadigit{1}}{\isacharparenright}} is used
  to encode basic names.

  \medskip Indexnames may acquire arbitrary large index numbers over
  time.  Results are usually normalized towards \isa{{\isadigit{0}}} at certain
  checkpoints, such that the very end of a proof.  This works by
  producing variants of the corresponding basic names
  (\secref{sec:basic-names}).  For example, the collection \isa{{\isacharquery}x{\isachardot}{\isadigit{1}}{\isacharcomma}\ {\isacharquery}x{\isachardot}{\isadigit{7}}{\isacharcomma}\ {\isacharquery}x{\isachardot}{\isadigit{4}}{\isadigit{2}}} then becomes \isa{{\isacharquery}x{\isacharcomma}\ {\isacharquery}xa{\isacharcomma}\ {\isacharquery}xb}.%
\end{isamarkuptext}%
\isamarkuptrue%
%
\isadelimmlref
%
\endisadelimmlref
%
\isatagmlref
%
\begin{isamarkuptext}%
\begin{mldecls}
  \indexmltype{indexname}\verb|type indexname| \\
  \end{mldecls}

  \begin{description}

  \item \verb|indexname| represents indexed names.  This is an
  abbreviation for \verb|string * int|.  The second component is
  usually non-negative, except for situations where \isa{{\isacharparenleft}x{\isacharcomma}\ {\isacharminus}{\isadigit{1}}{\isacharparenright}}
  is used to embed plain names.

  \end{description}%
\end{isamarkuptext}%
\isamarkuptrue%
%
\endisatagmlref
{\isafoldmlref}%
%
\isadelimmlref
%
\endisadelimmlref
%
\isamarkupsubsection{Qualified names and name spaces%
}
\isamarkuptrue%
%
\begin{isamarkuptext}%
A \emph{qualified name} consists of a non-empty sequence of basic
  name components.  The packed representation a dot as separator, for
  example in ``\isa{A{\isachardot}b{\isachardot}c}''.  The last component is called
  \emph{base} name, the remaining prefix \emph{qualifier} (which may
  be empty).  The basic idea of qualified names is to encode a
  hierarchically structured name spaces by recording the access path
  as part of the name.  For example, \isa{A{\isachardot}b{\isachardot}c} may be understood
  as a local entity \isa{c} within a local structure \isa{b}
  of the enclosing structure \isa{A}.  Typically, name space
  hierarchies consist of 1--3 levels, but this need not be always so.

  The empty name is commonly used as an indication of unnamed
  entities, if this makes any sense.  The operations on qualified
  names are smart enough to pass through such improper names
  unchanged.

  \medskip A \isa{naming} policy tells how to turn a name
  specification into a fully qualified internal name (by the \isa{full} operation), and how to fully qualified names may be accessed
  externally.  For example, the default naming prefixes an implicit
  path from the context: \isa{x} is becomes \isa{path{\isachardot}x}
  internally; the standard accesses include \isa{x}, \isa{path{\isachardot}x}, and further partial \isa{path} specifications.
  Normally, the naming is implicit in the theory or proof context.
  There are separate versions of the corresponding operations for
  these context types.

  \medskip A \isa{name\ space} manages a collection of fully
  internalized names, together with a mapping between external names
  and internal names (in both directions).  The corresponding \isa{intern} and \isa{extern} operations are mostly used for
  parsing and printing only!  The \isa{declare} operation augments
  a name space according to a given naming policy.

  By general convention, there are separate name spaces for each kind
  of formal entity, such as logical constant, type, type class,
  theorem etc.  It is usually clear from the occurrence in concrete
  syntax (or from the scope) which kind of entity a name refers to.
  For example, the very same name \isa{c} may be used uniformly
  for a constant, type, type class, which are from separate syntactic
  categories.

  There are common schemes to name theorems systematically, according
  to the name of the main logical entity being involved, such as
  \isa{c{\isachardot}intro} and \isa{c{\isachardot}elim} for theorems related to
  constant \isa{c}.  This technique of mapping names from one
  space into another requires some care in order to avoid conflicts.
  In particular, theorem names derived from type or class names are
  better suffixed in addition to the usual qualification, e.g.\ \isa{c{\isacharunderscore}type{\isachardot}intro} and \isa{c{\isacharunderscore}class{\isachardot}intro} for theorems related to
  type \isa{c} and class \isa{c}, respectively.%
\end{isamarkuptext}%
\isamarkuptrue%
%
\isadelimmlref
%
\endisadelimmlref
%
\isatagmlref
%
\begin{isamarkuptext}%
\begin{mldecls}
  \indexml{NameSpace.base}\verb|NameSpace.base: string -> string| \\
  \indexml{NameSpace.drop-base}\verb|NameSpace.drop_base: string -> string| \\
  \indexml{NameSpace.append}\verb|NameSpace.append: string -> string -> string| \\
  \indexml{NameSpace.pack}\verb|NameSpace.pack: string list -> string| \\
  \indexml{NameSpace.unpack}\verb|NameSpace.unpack: string -> string list| \\[1ex]
  \indexmltype{NameSpace.naming}\verb|type NameSpace.naming| \\
  \indexml{NameSpace.default-naming}\verb|NameSpace.default_naming: NameSpace.naming| \\
  \indexml{NameSpace.add-path}\verb|NameSpace.add_path: string -> NameSpace.naming -> NameSpace.naming| \\
  \indexml{NameSpace.full}\verb|NameSpace.full: NameSpace.naming -> string -> string| \\[1ex]
  \indexmltype{NameSpace.T}\verb|type NameSpace.T| \\
  \indexml{NameSpace.empty}\verb|NameSpace.empty: NameSpace.T| \\
  \indexml{NameSpace.merge}\verb|NameSpace.merge: NameSpace.T * NameSpace.T -> NameSpace.T| \\
  \indexml{NameSpace.declare}\verb|NameSpace.declare: NameSpace.naming -> string -> NameSpace.T -> NameSpace.T| \\
  \indexml{NameSpace.intern}\verb|NameSpace.intern: NameSpace.T -> string -> string| \\
  \indexml{NameSpace.extern}\verb|NameSpace.extern: NameSpace.T -> string -> string| \\
  \end{mldecls}

  \begin{description}

  \item \verb|NameSpace.base|~\isa{name} returns the base name of a
  qualified name.

  \item \verb|NameSpace.drop_base|~\isa{name} returns the qualifier
  of a qualified name.

  \item \verb|NameSpace.append|~\isa{name\isactrlisub {\isadigit{1}}\ name\isactrlisub {\isadigit{2}}}
  appends two qualified names.

  \item \verb|NameSpace.pack|~\isa{name} and \isa{NameSpace{\isachardot}unpack}~\isa{names} convert between the packed
  string representation and explicit list form of qualified names.

  \item \verb|NameSpace.naming| represents the abstract concept of
  a naming policy.

  \item \verb|NameSpace.default_naming| is the default naming policy.
  In a theory context, this is usually augmented by a path prefix
  consisting of the theory name.

  \item \verb|NameSpace.add_path|~\isa{path\ naming} augments the
  naming policy by extending its path.

  \item \verb|NameSpace.full|\isa{naming\ name} turns a name
  specification (usually a basic name) into the fully qualified
  internal version, according to the given naming policy.

  \item \verb|NameSpace.T| represents name spaces.

  \item \verb|NameSpace.empty| and \verb|NameSpace.merge|~\isa{{\isacharparenleft}space\isactrlisub {\isadigit{1}}{\isacharcomma}\ space\isactrlisub {\isadigit{2}}{\isacharparenright}} provide basic operations for
  building name spaces in accordance to the usual theory data
  management (\secref{sec:context-data}).

  \item \verb|NameSpace.declare|~\isa{naming\ name\ space} enters a
  fully qualified name into the name space, with partial accesses
  being derived from the given policy.

  \item \verb|NameSpace.intern|~\isa{space\ name} internalizes a
  (partially qualified) external name.

  This operation is mostly for parsing.  Note that fully qualified
  names stemming from declarations are produced via \verb|NameSpace.full| (or derivatives for \verb|theory| or \verb|Proof.context|).

  \item \verb|NameSpace.extern|~\isa{space\ name} externalizes a
  (fully qualified) internal name.

  This operation is mostly for printing.  Note unqualified names are
  produced via \verb|NameSpace.base|.

  \end{description}%
\end{isamarkuptext}%
\isamarkuptrue%
%
\endisatagmlref
{\isafoldmlref}%
%
\isadelimmlref
%
\endisadelimmlref
%
\isadelimtheory
%
\endisadelimtheory
%
\isatagtheory
\isacommand{end}\isamarkupfalse%
%
\endisatagtheory
{\isafoldtheory}%
%
\isadelimtheory
%
\endisadelimtheory
\isanewline
\end{isabellebody}%
%%% Local Variables:
%%% mode: latex
%%% TeX-master: "root"
%%% End:

%
\begin{isabellebody}%
\def\isabellecontext{logic}%
%
\isadelimtheory
\isanewline
\isanewline
\isanewline
%
\endisadelimtheory
%
\isatagtheory
\isacommand{theory}\isamarkupfalse%
\ logic\ \isakeyword{imports}\ base\ \isakeyword{begin}%
\endisatagtheory
{\isafoldtheory}%
%
\isadelimtheory
%
\endisadelimtheory
%
\isamarkupchapter{Primitive logic \label{ch:logic}%
}
\isamarkuptrue%
%
\isamarkupsection{Names%
}
\isamarkuptrue%
%
\begin{isamarkuptext}%
FIXME%
\end{isamarkuptext}%
\isamarkuptrue%
%
\isamarkupsection{Types \label{sec:types}%
}
\isamarkuptrue%
%
\begin{isamarkuptext}%
\glossary{Type class}{FIXME}

  \glossary{Type arity}{FIXME}

  \glossary{Sort}{FIXME}

  FIXME classes and sorts


  \glossary{Type}{FIXME}

  \glossary{Type constructor}{FIXME}

  \glossary{Type variable}{FIXME}

  FIXME simple types%
\end{isamarkuptext}%
\isamarkuptrue%
%
\isamarkupsection{Terms \label{sec:terms}%
}
\isamarkuptrue%
%
\begin{isamarkuptext}%
\glossary{Term}{FIXME}

  FIXME de-Bruijn representation of lambda terms%
\end{isamarkuptext}%
\isamarkuptrue%
%
\begin{isamarkuptext}%
FIXME

\glossary{Schematic polymorphism}{FIXME}

\glossary{Type variable}{FIXME}%
\end{isamarkuptext}%
\isamarkuptrue%
%
\isamarkupsection{Theorems \label{sec:thms}%
}
\isamarkuptrue%
%
\begin{isamarkuptext}%
FIXME

\glossary{Proposition}{A \seeglossary{term} of \seeglossary{type}
\isa{prop}.  Internally, there is nothing special about
propositions apart from their type, but the concrete syntax enforces a
clear distinction.  Propositions are structured via implication \isa{A\ {\isasymLongrightarrow}\ B} or universal quantification \isa{{\isasymAnd}x{\isachardot}\ B\ x} --- anything
else is considered atomic.  The canonical form for propositions is
that of a \seeglossary{Hereditary Harrop Formula}.}

\glossary{Theorem}{A proven proposition within a certain theory and
proof context, formally \isa{{\isasymGamma}\ {\isasymturnstile}\isactrlsub {\isasymTheta}\ {\isasymphi}}; both contexts are
rarely spelled out explicitly.  Theorems are usually normalized
according to the \seeglossary{HHF} format.}

\glossary{Fact}{Sometimes used interchangably for
\seeglossary{theorem}.  Strictly speaking, a list of theorems,
essentially an extra-logical conjunction.  Facts emerge either as
local assumptions, or as results of local goal statements --- both may
be simultaneous, hence the list representation.}

\glossary{Schematic variable}{FIXME}

\glossary{Fixed variable}{A variable that is bound within a certain
proof context; an arbitrary-but-fixed entity within a portion of proof
text.}

\glossary{Free variable}{Synonymous for \seeglossary{fixed variable}.}

\glossary{Bound variable}{FIXME}

\glossary{Variable}{See \seeglossary{schematic variable},
\seeglossary{fixed variable}, \seeglossary{bound variable}, or
\seeglossary{type variable}.  The distinguishing feature of different
variables is their binding scope.}%
\end{isamarkuptext}%
\isamarkuptrue%
%
\isamarkupsubsection{Primitive inferences%
}
\isamarkuptrue%
%
\begin{isamarkuptext}%
FIXME%
\end{isamarkuptext}%
\isamarkuptrue%
%
\isamarkupsubsection{Higher-order resolution%
}
\isamarkuptrue%
%
\begin{isamarkuptext}%
FIXME

\glossary{Hereditary Harrop Formula}{The set of propositions in HHF
format is defined inductively as \isa{H\ {\isacharequal}\ {\isacharparenleft}{\isasymAnd}x\isactrlsup {\isacharasterisk}{\isachardot}\ H\isactrlsup {\isacharasterisk}\ {\isasymLongrightarrow}\ A{\isacharparenright}}, for variables \isa{x} and atomic propositions \isa{A}.
Any proposition may be put into HHF form by normalizing with the rule
\isa{{\isacharparenleft}A\ {\isasymLongrightarrow}\ {\isacharparenleft}{\isasymAnd}x{\isachardot}\ B\ x{\isacharparenright}{\isacharparenright}\ {\isasymequiv}\ {\isacharparenleft}{\isasymAnd}x{\isachardot}\ A\ {\isasymLongrightarrow}\ B\ x{\isacharparenright}}.  In Isabelle, the outermost
quantifier prefix is represented via \seeglossary{schematic
variables}, such that the top-level structure is merely that of a
\seeglossary{Horn Clause}}.

\glossary{HHF}{See \seeglossary{Hereditary Harrop Formula}.}%
\end{isamarkuptext}%
\isamarkuptrue%
%
\isamarkupsubsection{Equational reasoning%
}
\isamarkuptrue%
%
\begin{isamarkuptext}%
FIXME%
\end{isamarkuptext}%
\isamarkuptrue%
%
\isamarkupsection{Low-level specifications%
}
\isamarkuptrue%
%
\begin{isamarkuptext}%
FIXME

\glossary{Constant}{Essentially a \seeglossary{fixed variable} of the
theory context, but slightly more flexible since it may be used at
different type-instances, due to \seeglossary{schematic
polymorphism.}}

\glossary{Axiom}{FIXME}

\glossary{Primitive definition}{FIXME}%
\end{isamarkuptext}%
\isamarkuptrue%
%
\isadelimtheory
%
\endisadelimtheory
%
\isatagtheory
\isacommand{end}\isamarkupfalse%
%
\endisatagtheory
{\isafoldtheory}%
%
\isadelimtheory
%
\endisadelimtheory
\isanewline
\end{isabellebody}%
%%% Local Variables:
%%% mode: latex
%%% TeX-master: "root"
%%% End:

%% $Id$
\chapter{Tactics} \label{tactics}
\index{tactics|(} Tactics have type \mltydx{tactic}.  This is just an
abbreviation for functions from theorems to theorem sequences, where
the theorems represent states of a backward proof.  Tactics seldom
need to be coded from scratch, as functions; instead they are
expressed using basic tactics and tacticals.

This chapter only presents the primitive tactics.  Substantial proofs
require the power of automatic tools like simplification
(Chapter~\ref{chap:simplification}) and classical tableau reasoning
(Chapter~\ref{chap:classical}).

\section{Resolution and assumption tactics}
{\bf Resolution} is Isabelle's basic mechanism for refining a subgoal using
a rule.  {\bf Elim-resolution} is particularly suited for elimination
rules, while {\bf destruct-resolution} is particularly suited for
destruction rules.  The {\tt r}, {\tt e}, {\tt d} naming convention is
maintained for several different kinds of resolution tactics, as well as
the shortcuts in the subgoal module.

All the tactics in this section act on a subgoal designated by a positive
integer~$i$.  They fail (by returning the empty sequence) if~$i$ is out of
range.

\subsection{Resolution tactics}
\index{resolution!tactics}
\index{tactics!resolution|bold}
\begin{ttbox} 
resolve_tac  : thm list -> int -> tactic
eresolve_tac : thm list -> int -> tactic
dresolve_tac : thm list -> int -> tactic
forward_tac  : thm list -> int -> tactic 
\end{ttbox}
These perform resolution on a list of theorems, $thms$, representing a list
of object-rules.  When generating next states, they take each of the rules
in the order given.  Each rule may yield several next states, or none:
higher-order resolution may yield multiple resolvents.
\begin{ttdescription}
\item[\ttindexbold{resolve_tac} {\it thms} {\it i}] 
  refines the proof state using the rules, which should normally be
  introduction rules.  It resolves a rule's conclusion with
  subgoal~$i$ of the proof state.

\item[\ttindexbold{eresolve_tac} {\it thms} {\it i}] 
  \index{elim-resolution}
  performs elim-resolution with the rules, which should normally be
  elimination rules.  It resolves with a rule, proves its first premise by
  assumption, and finally \emph{deletes} that assumption from any new
  subgoals.  (To rotate a rule's premises,
  see \texttt{rotate_prems} in~\S\ref{MiscellaneousForwardRules}.)

\item[\ttindexbold{dresolve_tac} {\it thms} {\it i}] 
  \index{forward proof}\index{destruct-resolution}
  performs destruct-resolution with the rules, which normally should
  be destruction rules.  This replaces an assumption by the result of
  applying one of the rules.

\item[\ttindexbold{forward_tac}]\index{forward proof}
  is like {\tt dresolve_tac} except that the selected assumption is not
  deleted.  It applies a rule to an assumption, adding the result as a new
  assumption.
\end{ttdescription}

\subsection{Assumption tactics}
\index{tactics!assumption|bold}\index{assumptions!tactics for}
\begin{ttbox} 
assume_tac    : int -> tactic
eq_assume_tac : int -> tactic
\end{ttbox} 
\begin{ttdescription}
\item[\ttindexbold{assume_tac} {\it i}] 
attempts to solve subgoal~$i$ by assumption.

\item[\ttindexbold{eq_assume_tac}] 
is like {\tt assume_tac} but does not use unification.  It succeeds (with a
\emph{unique} next state) if one of the assumptions is identical to the
subgoal's conclusion.  Since it does not instantiate variables, it cannot
make other subgoals unprovable.  It is intended to be called from proof
strategies, not interactively.
\end{ttdescription}

\subsection{Matching tactics} \label{match_tac}
\index{tactics!matching}
\begin{ttbox} 
match_tac  : thm list -> int -> tactic
ematch_tac : thm list -> int -> tactic
dmatch_tac : thm list -> int -> tactic
\end{ttbox}
These are just like the resolution tactics except that they never
instantiate unknowns in the proof state.  Flexible subgoals are not updated
willy-nilly, but are left alone.  Matching --- strictly speaking --- means
treating the unknowns in the proof state as constants; these tactics merely
discard unifiers that would update the proof state.
\begin{ttdescription}
\item[\ttindexbold{match_tac} {\it thms} {\it i}] 
refines the proof state using the rules, matching a rule's
conclusion with subgoal~$i$ of the proof state.

\item[\ttindexbold{ematch_tac}] 
is like {\tt match_tac}, but performs elim-resolution.

\item[\ttindexbold{dmatch_tac}] 
is like {\tt match_tac}, but performs destruct-resolution.
\end{ttdescription}


\subsection{Resolution with instantiation} \label{res_inst_tac}
\index{tactics!instantiation}\index{instantiation}
\begin{ttbox} 
res_inst_tac  : (string*string)list -> thm -> int -> tactic
eres_inst_tac : (string*string)list -> thm -> int -> tactic
dres_inst_tac : (string*string)list -> thm -> int -> tactic
forw_inst_tac : (string*string)list -> thm -> int -> tactic
\end{ttbox}
These tactics are designed for applying rules such as substitution and
induction, which cause difficulties for higher-order unification.  The
tactics accept explicit instantiations for unknowns in the rule ---
typically, in the rule's conclusion.  Each instantiation is a pair
{\tt($v$,$e$)}, where $v$ is an unknown \emph{without} its leading
question mark!
\begin{itemize}
\item If $v$ is the type unknown {\tt'a}, then
the rule must contain a type unknown \verb$?'a$ of some
sort~$s$, and $e$ should be a type of sort $s$.

\item If $v$ is the unknown {\tt P}, then
the rule must contain an unknown \verb$?P$ of some type~$\tau$,
and $e$ should be a term of some type~$\sigma$ such that $\tau$ and
$\sigma$ are unifiable.  If the unification of $\tau$ and $\sigma$
instantiates any type unknowns in $\tau$, these instantiations
are recorded for application to the rule.
\end{itemize}
Types are instantiated before terms.  Because type instantiations are
inferred from term instantiations, explicit type instantiations are seldom
necessary --- if \verb$?t$ has type \verb$?'a$, then the instantiation list
\verb$[("'a","bool"),("t","True")]$ may be simplified to
\verb$[("t","True")]$.  Type unknowns in the proof state may cause
failure because the tactics cannot instantiate them.

The instantiation tactics act on a given subgoal.  Terms in the
instantiations are type-checked in the context of that subgoal --- in
particular, they may refer to that subgoal's parameters.  Any unknowns in
the terms receive subscripts and are lifted over the parameters; thus, you
may not refer to unknowns in the subgoal.

\begin{ttdescription}
\item[\ttindexbold{res_inst_tac} {\it insts} {\it thm} {\it i}]
instantiates the rule {\it thm} with the instantiations {\it insts}, as
described above, and then performs resolution on subgoal~$i$.  Resolution
typically causes further instantiations; you need not give explicit
instantiations for every unknown in the rule.

\item[\ttindexbold{eres_inst_tac}] 
is like {\tt res_inst_tac}, but performs elim-resolution.

\item[\ttindexbold{dres_inst_tac}] 
is like {\tt res_inst_tac}, but performs destruct-resolution.

\item[\ttindexbold{forw_inst_tac}] 
is like {\tt dres_inst_tac} except that the selected assumption is not
deleted.  It applies the instantiated rule to an assumption, adding the
result as a new assumption.
\end{ttdescription}


\section{Other basic tactics}
\subsection{Tactic shortcuts}
\index{shortcuts!for tactics}
\index{tactics!resolution}\index{tactics!assumption}
\index{tactics!meta-rewriting}
\begin{ttbox} 
rtac     :      thm -> int -> tactic
etac     :      thm -> int -> tactic
dtac     :      thm -> int -> tactic
atac     :             int -> tactic
ares_tac : thm list -> int -> tactic
rewtac   :      thm ->        tactic
\end{ttbox}
These abbreviate common uses of tactics.
\begin{ttdescription}
\item[\ttindexbold{rtac} {\it thm} {\it i}] 
abbreviates \hbox{\tt resolve_tac [{\it thm}] {\it i}}, doing resolution.

\item[\ttindexbold{etac} {\it thm} {\it i}] 
abbreviates \hbox{\tt eresolve_tac [{\it thm}] {\it i}}, doing elim-resolution.

\item[\ttindexbold{dtac} {\it thm} {\it i}] 
abbreviates \hbox{\tt dresolve_tac [{\it thm}] {\it i}}, doing
destruct-resolution.

\item[\ttindexbold{atac} {\it i}] 
abbreviates \hbox{\tt assume_tac {\it i}}, doing proof by assumption.

\item[\ttindexbold{ares_tac} {\it thms} {\it i}] 
tries proof by assumption and resolution; it abbreviates
\begin{ttbox}
assume_tac {\it i} ORELSE resolve_tac {\it thms} {\it i}
\end{ttbox}

\item[\ttindexbold{rewtac} {\it def}] 
abbreviates \hbox{\tt rewrite_goals_tac [{\it def}]}, unfolding a definition.
\end{ttdescription}


\subsection{Inserting premises and facts}\label{cut_facts_tac}
\index{tactics!for inserting facts}\index{assumptions!inserting}
\begin{ttbox} 
cut_facts_tac : thm list -> int -> tactic
cut_inst_tac  : (string*string)list -> thm -> int -> tactic
subgoal_tac   : string -> int -> tactic
subgoal_tacs  : string list -> int -> tactic
\end{ttbox}
These tactics add assumptions to a subgoal.
\begin{ttdescription}
\item[\ttindexbold{cut_facts_tac} {\it thms} {\it i}] 
  adds the {\it thms} as new assumptions to subgoal~$i$.  Once they have
  been inserted as assumptions, they become subject to tactics such as {\tt
    eresolve_tac} and {\tt rewrite_goals_tac}.  Only rules with no premises
  are inserted: Isabelle cannot use assumptions that contain $\Imp$
  or~$\Forall$.  Sometimes the theorems are premises of a rule being
  derived, returned by~{\tt goal}; instead of calling this tactic, you
  could state the goal with an outermost meta-quantifier.

\item[\ttindexbold{cut_inst_tac} {\it insts} {\it thm} {\it i}]
  instantiates the {\it thm} with the instantiations {\it insts}, as
  described in \S\ref{res_inst_tac}.  It adds the resulting theorem as a
  new assumption to subgoal~$i$. 

\item[\ttindexbold{subgoal_tac} {\it formula} {\it i}] 
adds the {\it formula} as a assumption to subgoal~$i$, and inserts the same
{\it formula} as a new subgoal, $i+1$.

\item[\ttindexbold{subgoals_tac} {\it formulae} {\it i}] 
  uses {\tt subgoal_tac} to add the members of the list of {\it
    formulae} as assumptions to subgoal~$i$. 
\end{ttdescription}


\subsection{``Putting off'' a subgoal}
\begin{ttbox} 
defer_tac : int -> tactic
\end{ttbox}
\begin{ttdescription}
\item[\ttindexbold{defer_tac} {\it i}] 
  moves subgoal~$i$ to the last position in the proof state.  It can be
  useful when correcting a proof script: if the tactic given for subgoal~$i$
  fails, calling {\tt defer_tac} instead will let you continue with the rest
  of the script.

  The tactic fails if subgoal~$i$ does not exist or if the proof state
  contains type unknowns. 
\end{ttdescription}


\subsection{Definitions and meta-level rewriting} \label{sec:rewrite_goals}
\index{tactics!meta-rewriting|bold}\index{meta-rewriting|bold}
\index{definitions}

Definitions in Isabelle have the form $t\equiv u$, where $t$ is typically a
constant or a constant applied to a list of variables, for example $\it
sqr(n)\equiv n\times n$.  Conditional definitions, $\phi\Imp t\equiv u$,
are also supported.  {\bf Unfolding} the definition ${t\equiv u}$ means using
it as a rewrite rule, replacing~$t$ by~$u$ throughout a theorem.  {\bf
Folding} $t\equiv u$ means replacing~$u$ by~$t$.  Rewriting continues until
no rewrites are applicable to any subterm.

There are rules for unfolding and folding definitions; Isabelle does not do
this automatically.  The corresponding tactics rewrite the proof state,
yielding a single next state.  See also the {\tt goalw} command, which is the
easiest way of handling definitions.
\begin{ttbox} 
rewrite_goals_tac : thm list -> tactic
rewrite_tac       : thm list -> tactic
fold_goals_tac    : thm list -> tactic
fold_tac          : thm list -> tactic
\end{ttbox}
\begin{ttdescription}
\item[\ttindexbold{rewrite_goals_tac} {\it defs}]  
unfolds the {\it defs} throughout the subgoals of the proof state, while
leaving the main goal unchanged.  Use \ttindex{SELECT_GOAL} to restrict it to a
particular subgoal.

\item[\ttindexbold{rewrite_tac} {\it defs}]  
unfolds the {\it defs} throughout the proof state, including the main goal
--- not normally desirable!

\item[\ttindexbold{fold_goals_tac} {\it defs}]  
folds the {\it defs} throughout the subgoals of the proof state, while
leaving the main goal unchanged.

\item[\ttindexbold{fold_tac} {\it defs}]  
folds the {\it defs} throughout the proof state.
\end{ttdescription}

\begin{warn}
  These tactics only cope with definitions expressed as meta-level
  equalities ($\equiv$).  More general equivalences are handled by the
  simplifier, provided that it is set up appropriately for your logic
  (see Chapter~\ref{chap:simplification}).
\end{warn}

\subsection{Theorems useful with tactics}
\index{theorems!of pure theory}
\begin{ttbox} 
asm_rl: thm 
cut_rl: thm 
\end{ttbox}
\begin{ttdescription}
\item[\tdx{asm_rl}] 
is $\psi\Imp\psi$.  Under elim-resolution it does proof by assumption, and
\hbox{\tt eresolve_tac (asm_rl::{\it thms}) {\it i}} is equivalent to
\begin{ttbox} 
assume_tac {\it i}  ORELSE  eresolve_tac {\it thms} {\it i}
\end{ttbox}

\item[\tdx{cut_rl}] 
is $\List{\psi\Imp\theta,\psi}\Imp\theta$.  It is useful for inserting
assumptions; it underlies {\tt forward_tac}, {\tt cut_facts_tac}
and {\tt subgoal_tac}.
\end{ttdescription}


\section{Obscure tactics}

\subsection{Renaming parameters in a goal} \index{parameters!renaming}
\begin{ttbox} 
rename_tac        : string -> int -> tactic
rename_last_tac   : string -> string list -> int -> tactic
Logic.set_rename_prefix : string -> unit
Logic.auto_rename       : bool ref      \hfill{\bf initially false}
\end{ttbox}
When creating a parameter, Isabelle chooses its name by matching variable
names via the object-rule.  Given the rule $(\forall I)$ formalized as
$\left(\Forall x. P(x)\right) \Imp \forall x.P(x)$, Isabelle will note that
the $\Forall$-bound variable in the premise has the same name as the
$\forall$-bound variable in the conclusion.  

Sometimes there is insufficient information and Isabelle chooses an
arbitrary name.  The renaming tactics let you override Isabelle's choice.
Because renaming parameters has no logical effect on the proof state, the
{\tt by} command prints the message {\tt Warning:\ same as previous
level}.

Alternatively, you can suppress the naming mechanism described above and
have Isabelle generate uniform names for parameters.  These names have the
form $p${\tt a}, $p${\tt b}, $p${\tt c},~\ldots, where $p$ is any desired
prefix.  They are ugly but predictable.

\begin{ttdescription}
\item[\ttindexbold{rename_tac} {\it str} {\it i}] 
interprets the string {\it str} as a series of blank-separated variable
names, and uses them to rename the parameters of subgoal~$i$.  The names
must be distinct.  If there are fewer names than parameters, then the
tactic renames the innermost parameters and may modify the remaining ones
to ensure that all the parameters are distinct.

\item[\ttindexbold{rename_last_tac} {\it prefix} {\it suffixes} {\it i}] 
generates a list of names by attaching each of the {\it suffixes\/} to the 
{\it prefix}.  It is intended for coding structural induction tactics,
where several of the new parameters should have related names.

\item[\ttindexbold{Logic.set_rename_prefix} {\it prefix};] 
sets the prefix for uniform renaming to~{\it prefix}.  The default prefix
is {\tt"k"}.

\item[set \ttindexbold{Logic.auto_rename};] 
makes Isabelle generate uniform names for parameters. 
\end{ttdescription}


\subsection{Manipulating assumptions}
\index{assumptions!rotating}
\begin{ttbox} 
thin_tac   : string -> int -> tactic
rotate_tac : int -> int -> tactic
\end{ttbox}
\begin{ttdescription}
\item[\ttindexbold{thin_tac} {\it formula} $i$]  
\index{assumptions!deleting}
deletes the specified assumption from subgoal $i$.  Often the assumption
can be abbreviated, replacing subformul{\ae} by unknowns; the first matching
assumption will be deleted.  Removing useless assumptions from a subgoal
increases its readability and can make search tactics run faster.

\item[\ttindexbold{rotate_tac} $n$ $i$]  
\index{assumptions!rotating}
rotates the assumptions of subgoal $i$ by $n$ positions: from right to left
if $n$ is positive, and from left to right if $n$ is negative.  This is 
sometimes necessary in connection with \ttindex{asm_full_simp_tac}, which 
processes assumptions from left to right.
\end{ttdescription}


\subsection{Tidying the proof state}
\index{duplicate subgoals!removing}
\index{parameters!removing unused}
\index{flex-flex constraints}
\begin{ttbox} 
distinct_subgoals_tac : tactic
prune_params_tac      : tactic
flexflex_tac          : tactic
\end{ttbox}
\begin{ttdescription}
\item[\ttindexbold{distinct_subgoals_tac}]  
  removes duplicate subgoals from a proof state.  (These arise especially
  in \ZF{}, where the subgoals are essentially type constraints.)

\item[\ttindexbold{prune_params_tac}]  
  removes unused parameters from all subgoals of the proof state.  It works
  by rewriting with the theorem $(\Forall x. V)\equiv V$.  This tactic can
  make the proof state more readable.  It is used with
  \ttindex{rule_by_tactic} to simplify the resulting theorem.

\item[\ttindexbold{flexflex_tac}]  
  removes all flex-flex pairs from the proof state by applying the trivial
  unifier.  This drastic step loses information, and should only be done as
  the last step of a proof.

  Flex-flex constraints arise from difficult cases of higher-order
  unification.  To prevent this, use \ttindex{res_inst_tac} to instantiate
  some variables in a rule~(\S\ref{res_inst_tac}).  Normally flex-flex
  constraints can be ignored; they often disappear as unknowns get
  instantiated.
\end{ttdescription}


\subsection{Composition: resolution without lifting}
\index{tactics!for composition}
\begin{ttbox}
compose_tac: (bool * thm * int) -> int -> tactic
\end{ttbox}
{\bf Composing} two rules means resolving them without prior lifting or
renaming of unknowns.  This low-level operation, which underlies the
resolution tactics, may occasionally be useful for special effects.
A typical application is \ttindex{res_inst_tac}, which lifts and instantiates a
rule, then passes the result to {\tt compose_tac}.
\begin{ttdescription}
\item[\ttindexbold{compose_tac} ($flag$, $rule$, $m$) $i$] 
refines subgoal~$i$ using $rule$, without lifting.  The $rule$ is taken to
have the form $\List{\psi@1; \ldots; \psi@m} \Imp \psi$, where $\psi$ need
not be atomic; thus $m$ determines the number of new subgoals.  If
$flag$ is {\tt true} then it performs elim-resolution --- it solves the
first premise of~$rule$ by assumption and deletes that assumption.
\end{ttdescription}


\section{*Managing lots of rules}
These operations are not intended for interactive use.  They are concerned
with the processing of large numbers of rules in automatic proof
strategies.  Higher-order resolution involving a long list of rules is
slow.  Filtering techniques can shorten the list of rules given to
resolution, and can also detect whether a subgoal is too flexible,
with too many rules applicable.

\subsection{Combined resolution and elim-resolution} \label{biresolve_tac}
\index{tactics!resolution}
\begin{ttbox} 
biresolve_tac   : (bool*thm)list -> int -> tactic
bimatch_tac     : (bool*thm)list -> int -> tactic
subgoals_of_brl : bool*thm -> int
lessb           : (bool*thm) * (bool*thm) -> bool
\end{ttbox}
{\bf Bi-resolution} takes a list of $\it (flag,rule)$ pairs.  For each
pair, it applies resolution if the flag is~{\tt false} and
elim-resolution if the flag is~{\tt true}.  A single tactic call handles a
mixture of introduction and elimination rules.

\begin{ttdescription}
\item[\ttindexbold{biresolve_tac} {\it brls} {\it i}] 
refines the proof state by resolution or elim-resolution on each rule, as
indicated by its flag.  It affects subgoal~$i$ of the proof state.

\item[\ttindexbold{bimatch_tac}] 
is like {\tt biresolve_tac}, but performs matching: unknowns in the
proof state are never updated (see~\S\ref{match_tac}).

\item[\ttindexbold{subgoals_of_brl}({\it flag},{\it rule})] 
returns the number of new subgoals that bi-res\-o\-lu\-tion would yield for the
pair (if applied to a suitable subgoal).  This is $n$ if the flag is
{\tt false} and $n-1$ if the flag is {\tt true}, where $n$ is the number
of premises of the rule.  Elim-resolution yields one fewer subgoal than
ordinary resolution because it solves the major premise by assumption.

\item[\ttindexbold{lessb} ({\it brl1},{\it brl2})] 
returns the result of 
\begin{ttbox}
subgoals_of_brl{\it brl1} < subgoals_of_brl{\it brl2}
\end{ttbox}
\end{ttdescription}
Note that \hbox{\tt sort lessb {\it brls}} sorts a list of $\it
(flag,rule)$ pairs by the number of new subgoals they will yield.  Thus,
those that yield the fewest subgoals should be tried first.


\subsection{Discrimination nets for fast resolution}\label{filt_resolve_tac}
\index{discrimination nets|bold}
\index{tactics!resolution}
\begin{ttbox} 
net_resolve_tac  : thm list -> int -> tactic
net_match_tac    : thm list -> int -> tactic
net_biresolve_tac: (bool*thm) list -> int -> tactic
net_bimatch_tac  : (bool*thm) list -> int -> tactic
filt_resolve_tac : thm list -> int -> int -> tactic
could_unify      : term*term->bool
filter_thms      : (term*term->bool) -> int*term*thm list -> thm list
\end{ttbox}
The module {\tt Net} implements a discrimination net data structure for
fast selection of rules \cite[Chapter 14]{charniak80}.  A term is
classified by the symbol list obtained by flattening it in preorder.
The flattening takes account of function applications, constants, and free
and bound variables; it identifies all unknowns and also regards
\index{lambda abs@$\lambda$-abstractions}
$\lambda$-abstractions as unknowns, since they could $\eta$-contract to
anything.  

A discrimination net serves as a polymorphic dictionary indexed by terms.
The module provides various functions for inserting and removing items from
nets.  It provides functions for returning all items whose term could match
or unify with a target term.  The matching and unification tests are
overly lax (due to the identifications mentioned above) but they serve as
useful filters.

A net can store introduction rules indexed by their conclusion, and
elimination rules indexed by their major premise.  Isabelle provides
several functions for `compiling' long lists of rules into fast
resolution tactics.  When supplied with a list of theorems, these functions
build a discrimination net; the net is used when the tactic is applied to a
goal.  To avoid repeatedly constructing the nets, use currying: bind the
resulting tactics to \ML{} identifiers.

\begin{ttdescription}
\item[\ttindexbold{net_resolve_tac} {\it thms}] 
builds a discrimination net to obtain the effect of a similar call to {\tt
resolve_tac}.

\item[\ttindexbold{net_match_tac} {\it thms}] 
builds a discrimination net to obtain the effect of a similar call to {\tt
match_tac}.

\item[\ttindexbold{net_biresolve_tac} {\it brls}] 
builds a discrimination net to obtain the effect of a similar call to {\tt
biresolve_tac}.

\item[\ttindexbold{net_bimatch_tac} {\it brls}] 
builds a discrimination net to obtain the effect of a similar call to {\tt
bimatch_tac}.

\item[\ttindexbold{filt_resolve_tac} {\it thms} {\it maxr} {\it i}] 
uses discrimination nets to extract the {\it thms} that are applicable to
subgoal~$i$.  If more than {\it maxr\/} theorems are applicable then the
tactic fails.  Otherwise it calls {\tt resolve_tac}.  

This tactic helps avoid runaway instantiation of unknowns, for example in
type inference.

\item[\ttindexbold{could_unify} ({\it t},{\it u})] 
returns {\tt false} if~$t$ and~$u$ are `obviously' non-unifiable, and
otherwise returns~{\tt true}.  It assumes all variables are distinct,
reporting that {\tt ?a=?a} may unify with {\tt 0=1}.

\item[\ttindexbold{filter_thms} $could\; (limit,prem,thms)$] 
returns the list of potentially resolvable rules (in {\it thms\/}) for the
subgoal {\it prem}, using the predicate {\it could\/} to compare the
conclusion of the subgoal with the conclusion of each rule.  The resulting list
is no longer than {\it limit}.
\end{ttdescription}


\section{Programming tools for proof strategies}
Do not consider using the primitives discussed in this section unless you
really need to code tactics from scratch.

\subsection{Operations on tactics}
\index{tactics!primitives for coding} A tactic maps theorems to sequences of
theorems.  The type constructor for sequences (lazy lists) is called
\mltydx{Seq.seq}.  To simplify the types of tactics and tacticals,
Isabelle defines a type abbreviation:
\begin{ttbox} 
type tactic = thm -> thm Seq.seq
\end{ttbox} 
The following operations provide means for coding tactics in a clean style.
\begin{ttbox} 
PRIMITIVE :                  (thm -> thm) -> tactic  
SUBGOAL   : ((term*int) -> tactic) -> int -> tactic
\end{ttbox} 
\begin{ttdescription}
\item[\ttindexbold{PRIMITIVE} $f$] packages the meta-rule~$f$ as a tactic that
  applies $f$ to the proof state and returns the result as a one-element
  sequence.  If $f$ raises an exception, then the tactic's result is the empty
  sequence.

\item[\ttindexbold{SUBGOAL} $f$ $i$] 
extracts subgoal~$i$ from the proof state as a term~$t$, and computes a
tactic by calling~$f(t,i)$.  It applies the resulting tactic to the same
state.  The tactic body is expressed using tactics and tacticals, but may
peek at a particular subgoal:
\begin{ttbox} 
SUBGOAL (fn (t,i) => {\it tactic-valued expression})
\end{ttbox} 
\end{ttdescription}


\subsection{Tracing}
\index{tactics!tracing}
\index{tracing!of tactics}
\begin{ttbox} 
pause_tac: tactic
print_tac: tactic
\end{ttbox}
These tactics print tracing information when they are applied to a proof
state.  Their output may be difficult to interpret.  Note that certain of
the searching tacticals, such as {\tt REPEAT}, have built-in tracing
options.
\begin{ttdescription}
\item[\ttindexbold{pause_tac}] 
prints {\footnotesize\tt** Press RETURN to continue:} and then reads a line
from the terminal.  If this line is blank then it returns the proof state
unchanged; otherwise it fails (which may terminate a repetition).

\item[\ttindexbold{print_tac}] 
returns the proof state unchanged, with the side effect of printing it at
the terminal.
\end{ttdescription}


\section{*Sequences}
\index{sequences (lazy lists)|bold}
The module {\tt Seq} declares a type of lazy lists.  It uses
Isabelle's type \mltydx{option} to represent the possible presence
(\ttindexbold{Some}) or absence (\ttindexbold{None}) of
a value:
\begin{ttbox}
datatype 'a option = None  |  Some of 'a;
\end{ttbox}
The {\tt Seq} structure is supposed to be accessed via fully qualified
names and should not be \texttt{open}.

\subsection{Basic operations on sequences}
\begin{ttbox} 
Seq.empty   : 'a seq
Seq.make    : (unit -> ('a * 'a seq) option) -> 'a seq
Seq.single  : 'a -> 'a seq
Seq.pull    : 'a seq -> ('a * 'a seq) option
\end{ttbox}
\begin{ttdescription}
\item[Seq.empty] is the empty sequence.

\item[\tt Seq.make (fn () => Some ($x$, $xq$))] constructs the
  sequence with head~$x$ and tail~$xq$, neither of which is evaluated.

\item[Seq.single $x$] 
constructs the sequence containing the single element~$x$.

\item[Seq.pull $xq$] returns {\tt None} if the sequence is empty and
  {\tt Some ($x$, $xq'$)} if the sequence has head~$x$ and tail~$xq'$.
  Warning: calling \hbox{Seq.pull $xq$} again will {\it recompute\/}
  the value of~$x$; it is not stored!
\end{ttdescription}


\subsection{Converting between sequences and lists}
\begin{ttbox} 
Seq.chop    : int * 'a seq -> 'a list * 'a seq
Seq.list_of : 'a seq -> 'a list
Seq.of_list : 'a list -> 'a seq
\end{ttbox}
\begin{ttdescription}
\item[Seq.chop ($n$, $xq$)] returns the first~$n$ elements of~$xq$ as a
  list, paired with the remaining elements of~$xq$.  If $xq$ has fewer
  than~$n$ elements, then so will the list.
  
\item[Seq.list_of $xq$] returns the elements of~$xq$, which must be
  finite, as a list.
  
\item[Seq.of_list $xs$] creates a sequence containing the elements
  of~$xs$.
\end{ttdescription}


\subsection{Combining sequences}
\begin{ttbox} 
Seq.append      : 'a seq * 'a seq -> 'a seq
Seq.interleave  : 'a seq * 'a seq -> 'a seq
Seq.flat        : 'a seq seq -> 'a seq
Seq.map         : ('a -> 'b) -> 'a seq -> 'b seq
Seq.filter      : ('a -> bool) -> 'a seq -> 'a seq
\end{ttbox} 
\begin{ttdescription}
\item[Seq.append ($xq$, $yq$)] concatenates $xq$ to $yq$.
  
\item[Seq.interleave ($xq$, $yq$)] joins $xq$ with $yq$ by
  interleaving their elements.  The result contains all the elements
  of the sequences, even if both are infinite.
  
\item[Seq.flat $xqq$] concatenates a sequence of sequences.
  
\item[Seq.map $f$ $xq$] applies $f$ to every element
  of~$xq=x@1,x@2,\ldots$, yielding the sequence $f(x@1),f(x@2),\ldots$.
  
\item[Seq.filter $p$ $xq$] returns the sequence consisting of all
  elements~$x$ of~$xq$ such that $p(x)$ is {\tt true}.
\end{ttdescription}

\index{tactics|)}


%%% Local Variables: 
%%% mode: latex
%%% TeX-master: "ref"
%%% End: 

%
\begin{isabellebody}%
\def\isabellecontext{proof}%
%
\isadelimtheory
\isanewline
\isanewline
\isanewline
%
\endisadelimtheory
%
\isatagtheory
\isacommand{theory}\isamarkupfalse%
\ {\isachardoublequoteopen}proof{\isachardoublequoteclose}\ \isakeyword{imports}\ base\ \isakeyword{begin}%
\endisatagtheory
{\isafoldtheory}%
%
\isadelimtheory
%
\endisadelimtheory
%
\isamarkupchapter{Structured proofs%
}
\isamarkuptrue%
%
\isamarkupsection{Variables and schematic polymorphism%
}
\isamarkuptrue%
%
\begin{isamarkuptext}%
FIXME%
\end{isamarkuptext}%
\isamarkuptrue%
%
\isadelimmlref
%
\endisadelimmlref
%
\isatagmlref
%
\begin{isamarkuptext}%
\begin{mldecls}
  \indexml{Variable.declare-term}\verb|Variable.declare_term: term -> Proof.context -> Proof.context| \\
  \indexml{Variable.add-fixes}\verb|Variable.add_fixes: string list -> Proof.context -> string list * Proof.context| \\
  \indexml{Variable.import}\verb|Variable.import: bool ->|\isasep\isanewline%
\verb|  thm list -> Proof.context -> ((ctyp list * cterm list) * thm list) * Proof.context| \\
  \indexml{Variable.export}\verb|Variable.export: Proof.context -> Proof.context -> thm list -> thm list| \\
  \indexml{Variable.trade}\verb|Variable.trade: Proof.context -> (thm list -> thm list) -> thm list -> thm list| \\
  \indexml{Variable.polymorphic}\verb|Variable.polymorphic: Proof.context -> term list -> term list| \\
  \end{mldecls}

  \begin{description}

  \item \verb|Variable.declare_term|~\isa{t\ ctxt} declares term
  \isa{t} to belong to the context.  This fixes free type
  variables, but not term variables.  Constraints for type and term
  variables are declared uniformly.

  \item \verb|Variable.add_fixes|~\isa{xs\ ctxt} fixes term
  variables \isa{xs} and returns the internal names of the
  resulting Skolem constants.  Note that term fixes refer to
  \emph{all} type instances that may occur in the future.

  \item \verb|Variable.invent_fixes| is similar to \verb|Variable.add_fixes|, but the given names merely act as hints for
  internal fixes produced here.

  \item \verb|Variable.import|~\isa{open\ ths\ ctxt} augments the
  context by new fixes for the schematic type and term variables
  occurring in \isa{ths}.  The \isa{open} flag indicates
  whether the fixed names should be accessible to the user, otherwise
  internal names are chosen.

  \item \verb|Variable.export|~\isa{inner\ outer\ ths} generalizes
  fixed type and term variables in \isa{ths} according to the
  difference of the \isa{inner} and \isa{outer} context.  Note
  that type variables occurring in term variables are still fixed.

  \verb|Variable.export| essentially reverses the effect of \verb|Variable.import| (up to renaming of schematic variables.

  \item \verb|Variable.trade| composes \verb|Variable.import| and \verb|Variable.export|, i.e.\ it provides a view on facts with all
  variables being fixed in the current context.

  \item \verb|Variable.polymorphic|~\isa{ctxt\ ts} generalizes type
  variables in \isa{ts} as far as possible, even those occurring
  in fixed term variables.  This operation essentially reverses the
  default policy of type-inference to introduce local polymorphism as
  fixed types.

  \end{description}%
\end{isamarkuptext}%
\isamarkuptrue%
%
\endisatagmlref
{\isafoldmlref}%
%
\isadelimmlref
%
\endisadelimmlref
%
\begin{isamarkuptext}%
FIXME%
\end{isamarkuptext}%
\isamarkuptrue%
%
\isamarkupsection{Assumptions%
}
\isamarkuptrue%
%
\begin{isamarkuptext}%
An \emph{assumption} is a proposition that it is postulated in the
  current context.  Local conclusions may use assumptions as
  additional facts, but this imposes implicit hypotheses that weaken
  the overall statement.

  Assumptions are restricted to fixed non-schematic statements, all
  generality needs to be expressed by explicit quantifiers.
  Nevertheless, the result will be in HHF normal form with outermost
  quantifiers stripped.  For example, by assuming \isa{{\isasymAnd}x\ {\isacharcolon}{\isacharcolon}\ {\isasymalpha}{\isachardot}\ P\ x} we get \isa{{\isasymAnd}x\ {\isacharcolon}{\isacharcolon}\ {\isasymalpha}{\isachardot}\ P\ x\ {\isasymturnstile}\ P\ {\isacharquery}x} for arbitrary \isa{{\isacharquery}x}
  of the fixed type \isa{{\isasymalpha}}.  Local derivations accumulate more
  and more explicit references to hypotheses: \isa{A\isactrlisub {\isadigit{1}}{\isacharcomma}\ {\isasymdots}{\isacharcomma}\ A\isactrlisub n\ {\isasymturnstile}\ B} where \isa{A\isactrlisub {\isadigit{1}}{\isacharcomma}\ {\isasymdots}{\isacharcomma}\ A\isactrlisub n} needs to
  be covered by the assumptions of the current context.

  \medskip The \isa{add{\isacharunderscore}assms} operation augments the context by
  local assumptions, which are parameterized by an arbitrary \isa{export} rule (see below).

  The \isa{export} operation moves facts from a (larger) inner
  context into a (smaller) outer context, by discharging the
  difference of the assumptions as specified by the associated export
  rules.  Note that the discharged portion is determined by the
  difference contexts, not the facts being exported!  There is a
  separate flag to indicate a goal context, where the result is meant
  to refine an enclosing sub-goal of a structured proof state (cf.\
  \secref{sec:isar-proof-state}).

  \medskip The most basic export rule discharges assumptions directly
  by means of the \isa{{\isasymLongrightarrow}} introduction rule:
  \[
  \infer[(\isa{{\isasymLongrightarrow}{\isacharunderscore}intro})]{\isa{{\isasymGamma}\ {\isacharbackslash}\ A\ {\isasymturnstile}\ A\ {\isasymLongrightarrow}\ B}}{\isa{{\isasymGamma}\ {\isasymturnstile}\ B}}
  \]

  The variant for goal refinements marks the newly introduced
  premises, which causes the builtin goal refinement scheme of Isar to
  enforce unification with local premises within the goal:
  \[
  \infer[(\isa{{\isacharhash}{\isasymLongrightarrow}{\isacharunderscore}intro})]{\isa{{\isasymGamma}\ {\isacharbackslash}\ A\ {\isasymturnstile}\ {\isacharhash}A\ {\isasymLongrightarrow}\ B}}{\isa{{\isasymGamma}\ {\isasymturnstile}\ B}}
  \]

  \medskip Alternative assumptions may perform arbitrary
  transformations on export, as long as a particular portion of
  hypotheses is removed from the given facts.  For example, a local
  definition works by fixing \isa{x} and assuming \isa{x\ {\isasymequiv}\ t},
  with the following export rule to reverse the effect:
  \[
  \infer{\isa{{\isasymGamma}\ {\isacharbackslash}\ x\ {\isasymequiv}\ t\ {\isasymturnstile}\ B\ t}}{\isa{{\isasymGamma}\ {\isasymturnstile}\ B\ x}}
  \]

  \medskip The general concept supports block-structured reasoning
  nicely, with arbitrary mechanisms for introducing local assumptions.
  The common reasoning pattern is as follows:

  \medskip
  \begin{tabular}{l}
  \isa{add{\isacharunderscore}assms\ e\isactrlisub {\isadigit{1}}\ A\isactrlisub {\isadigit{1}}} \\
  \isa{{\isasymdots}} \\
  \isa{add{\isacharunderscore}assms\ e\isactrlisub n\ A\isactrlisub n} \\
  \isa{export} \\
  \end{tabular}
  \medskip

  \noindent The final \isa{export} will turn any fact \isa{A\isactrlisub {\isadigit{1}}{\isacharcomma}\ {\isasymdots}{\isacharcomma}\ A\isactrlisub n\ {\isasymturnstile}\ B} into some \isa{{\isasymturnstile}\ B{\isacharprime}}, by
  applying the export rules \isa{e\isactrlisub {\isadigit{1}}{\isacharcomma}\ {\isasymdots}{\isacharcomma}\ e\isactrlisub n}
  inside-out.%
\end{isamarkuptext}%
\isamarkuptrue%
%
\isadelimmlref
%
\endisadelimmlref
%
\isatagmlref
%
\begin{isamarkuptext}%
\begin{mldecls}
  \indexmltype{Assumption.export}\verb|type Assumption.export| \\
  \indexml{Assumption.assume}\verb|Assumption.assume: cterm -> thm| \\
  \indexml{Assumption.add-assms}\verb|Assumption.add_assms: Assumption.export ->|\isasep\isanewline%
\verb|  cterm list -> Proof.context -> thm list * Proof.context| \\
  \indexml{Assumption.add-assumes}\verb|Assumption.add_assumes: |\isasep\isanewline%
\verb|  cterm list -> Proof.context -> thm list * Proof.context| \\
  \indexml{Assumption.export}\verb|Assumption.export: bool -> Proof.context -> Proof.context -> thm -> thm| \\
  \end{mldecls}

  \begin{description}

  \item \verb|Assumption.export| represents arbitrary export
  rules, which is any function of type \verb|bool -> cterm list -> thm -> thm|,
  where the \verb|bool| indicates goal mode, and the \verb|cterm list| the collection of assumptions to be discharged
  simultaneously.

  \item \verb|Assumption.assume|~\isa{A} turns proposition \isa{A} into a raw assumption \isa{A\ {\isasymturnstile}\ A{\isacharprime}}, where the conclusion
  \isa{A{\isacharprime}} is in HHF normal form.

  \item \verb|Assumption.add_assms|~\isa{e\ As} augments the context
  by assumptions \isa{As} with export rule \isa{e}.  The
  resulting facts are hypothetical theorems as produced by \verb|Assumption.assume|.

  \item \verb|Assumption.add_assumes|~\isa{As} is a special case of
  \verb|Assumption.add_assms| where the export rule performs \isa{{\isasymLongrightarrow}{\isacharunderscore}intro} or \isa{{\isacharhash}{\isasymLongrightarrow}{\isacharunderscore}intro}, depending on goal mode.

  \item \verb|Assumption.export|~\isa{is{\isacharunderscore}goal\ inner\ outer\ th}
  exports result \isa{th} from the the \isa{inner} context
  back into the \isa{outer} one; \isa{is{\isacharunderscore}goal\ {\isacharequal}\ true} means
  this is a goal context.  The result is in HHF normal form.  Note
  that \verb|ProofContext.export| combines \verb|Variable.export|
  and \verb|Assumption.export| in the canonical way.

  \end{description}%
\end{isamarkuptext}%
\isamarkuptrue%
%
\endisatagmlref
{\isafoldmlref}%
%
\isadelimmlref
%
\endisadelimmlref
%
\isamarkupsection{Conclusions%
}
\isamarkuptrue%
%
\begin{isamarkuptext}%
FIXME%
\end{isamarkuptext}%
\isamarkuptrue%
%
\isamarkupsection{Proof states \label{sec:isar-proof-state}%
}
\isamarkuptrue%
%
\begin{isamarkuptext}%
FIXME

\glossary{Proof state}{The whole configuration of a structured proof,
consisting of a \seeglossary{proof context} and an optional
\seeglossary{structured goal}.  Internally, an Isar proof state is
organized as a stack to accomodate block structure of proof texts.
For historical reasons, a low-level \seeglossary{tactical goal} is
occasionally called ``proof state'' as well.}

\glossary{Structured goal}{FIXME}

\glossary{Goal}{See \seeglossary{tactical goal} or \seeglossary{structured goal}. \norefpage}%
\end{isamarkuptext}%
\isamarkuptrue%
%
\isamarkupsection{Proof methods%
}
\isamarkuptrue%
%
\begin{isamarkuptext}%
FIXME%
\end{isamarkuptext}%
\isamarkuptrue%
%
\isamarkupsection{Attributes%
}
\isamarkuptrue%
%
\begin{isamarkuptext}%
FIXME ?!%
\end{isamarkuptext}%
\isamarkuptrue%
%
\isadelimtheory
%
\endisadelimtheory
%
\isatagtheory
\isacommand{end}\isamarkupfalse%
%
\endisatagtheory
{\isafoldtheory}%
%
\isadelimtheory
%
\endisadelimtheory
\isanewline
\end{isabellebody}%
%%% Local Variables:
%%% mode: latex
%%% TeX-master: "root"
%%% End:

%
\begin{isabellebody}%
\def\isabellecontext{locale}%
%
\isadelimtheory
\isanewline
\isanewline
\isanewline
%
\endisadelimtheory
%
\isatagtheory
\isacommand{theory}\isamarkupfalse%
\ {\isachardoublequoteopen}locale{\isachardoublequoteclose}\ \isakeyword{imports}\ base\ \isakeyword{begin}%
\endisatagtheory
{\isafoldtheory}%
%
\isadelimtheory
%
\endisadelimtheory
%
\isamarkupchapter{Structured specifications%
}
\isamarkuptrue%
%
\isamarkupsection{Specification elements%
}
\isamarkuptrue%
%
\begin{isamarkuptext}%
FIXME%
\end{isamarkuptext}%
\isamarkuptrue%
%
\isamarkupsection{Type-inference%
}
\isamarkuptrue%
%
\begin{isamarkuptext}%
FIXME%
\end{isamarkuptext}%
\isamarkuptrue%
%
\isamarkupsection{Local theories%
}
\isamarkuptrue%
%
\begin{isamarkuptext}%
FIXME

  \glossary{Local theory}{FIXME}%
\end{isamarkuptext}%
\isamarkuptrue%
%
\isadelimtheory
%
\endisadelimtheory
%
\isatagtheory
\isacommand{end}\isamarkupfalse%
%
\endisatagtheory
{\isafoldtheory}%
%
\isadelimtheory
%
\endisadelimtheory
\isanewline
\end{isabellebody}%
%%% Local Variables:
%%% mode: latex
%%% TeX-master: "root"
%%% End:

%
\begin{isabellebody}%
\def\isabellecontext{integration}%
%
\isadelimtheory
\isanewline
\isanewline
\isanewline
%
\endisadelimtheory
%
\isatagtheory
\isacommand{theory}\isamarkupfalse%
\ integration\ \isakeyword{imports}\ base\ \isakeyword{begin}%
\endisatagtheory
{\isafoldtheory}%
%
\isadelimtheory
%
\endisadelimtheory
%
\isamarkupchapter{System integration%
}
\isamarkuptrue%
%
\isamarkupsection{Isar toplevel \label{sec:isar-toplevel}%
}
\isamarkuptrue%
%
\begin{isamarkuptext}%
The Isar toplevel may be considered the centeral hub of the
  Isabelle/Isar system, where all key components and sub-systems are
  integrated into a single read-eval-print loop of Isar commands.
  Here we even incorporate the existing {\ML} toplevel of the compiler
  and run-time system (cf.\ \secref{sec:ML-toplevel}).

  Isabelle/Isar departs from original ``LCF system architecture''
  where {\ML} was really The Meta Language for defining theories and
  conducting proofs.  Instead, {\ML} merely serves as the
  implementation language for the system (and user extensions), while
  our specific Isar toplevel supports particular notions of
  incremental theory and proof development more directly.  This
  includes the graph structure of theories and the block structure of
  proofs, support for unlimited undo, facilities for tracing,
  debugging, timing, profiling.

  \medskip The toplevel maintains an implicit state, which is
  transformed by a sequence of transitions -- either interactively or
  in batch-mode.  In interactive mode, Isar state transitions are
  encapsulated as safe transactions, such that both failure and undo
  are handled conveniently without destroying the underlying draft
  theory (cf.~\secref{sec:context-theory}).  In batch mode,
  transitions operate in a strictly linear (destructive) fashion, such
  that error conditions abort the present attempt to construct a
  theory altogether.

  The toplevel state is a disjoint sum of empty \isa{toplevel}, or
  \isa{theory}, or \isa{proof}.  On entering the main Isar loop we
  start with an empty toplevel.  A theory is commenced by giving a
  \isa{{\isasymTHEORY}} header; within a theory we may issue theory
  commands such as \isa{{\isasymDEFINITION}}, or state a \isa{{\isasymTHEOREM}} to be proven.  Now we are within a proof state, with a
  rich collection of Isar proof commands for structured proof
  composition, or unstructured proof scripts.  When the proof is
  concluded we get back to the theory, which is then updated by
  storing the resulting fact.  Further theory declarations or theorem
  statements with proofs may follow, until we eventually conclude the
  theory development by issuing \isa{{\isasymEND}}.  The resulting theory
  is then stored within the theory database and we are back to the
  empty toplevel.

  In addition to these proper state transformations, there are also
  some diagnostic commands for peeking at the toplevel state without
  modifying it (e.g.\ \isakeyword{thm}, \isakeyword{term},
  \isakeyword{print-cases}).%
\end{isamarkuptext}%
\isamarkuptrue%
%
\isadelimmlref
%
\endisadelimmlref
%
\isatagmlref
%
\begin{isamarkuptext}%
\begin{mldecls}
  \indexmltype{Toplevel.state}\verb|type Toplevel.state| \\
  \indexml{Toplevel.UNDEF}\verb|Toplevel.UNDEF: exn| \\
  \indexml{Toplevel.is-toplevel}\verb|Toplevel.is_toplevel: Toplevel.state -> bool| \\
  \indexml{Toplevel.theory-of}\verb|Toplevel.theory_of: Toplevel.state -> theory| \\
  \indexml{Toplevel.proof-of}\verb|Toplevel.proof_of: Toplevel.state -> Proof.state| \\
  \indexml{Toplevel.debug}\verb|Toplevel.debug: bool ref| \\
  \indexml{Toplevel.timing}\verb|Toplevel.timing: bool ref| \\
  \indexml{Toplevel.profiling}\verb|Toplevel.profiling: int ref| \\
  \end{mldecls}

  \begin{description}

  \item \verb|Toplevel.state| represents Isar toplevel states,
  which are normally only manipulated through the toplevel transition
  concept (\secref{sec:toplevel-transition}).  Also note that a
  toplevel state is subject to the same linerarity restrictions as a
  theory context (cf.~\secref{sec:context-theory}).

  \item \verb|Toplevel.UNDEF| is raised for undefined toplevel
  operations: \verb|Toplevel.state| is a sum type, many operations
  work only partially for certain cases.

  \item \verb|Toplevel.is_toplevel| checks for an empty toplevel state.

  \item \verb|Toplevel.theory_of| gets the theory of a theory or proof
  (!), otherwise raises \verb|Toplevel.UNDEF|.

  \item \verb|Toplevel.proof_of| gets the Isar proof state if
  available, otherwise raises \verb|Toplevel.UNDEF|.

  \item \verb|set Toplevel.debug| makes the toplevel print further
  details about internal error conditions, exceptions being raised
  etc.

  \item \verb|set Toplevel.timing| makes the toplevel print timing
  information for each Isar command being executed.

  \item \verb|Toplevel.profiling| controls low-level profiling of the
  underlying {\ML} runtime system.\footnote{For Poly/ML, 1 means time
  and 2 space profiling.}

  \end{description}%
\end{isamarkuptext}%
\isamarkuptrue%
%
\endisatagmlref
{\isafoldmlref}%
%
\isadelimmlref
%
\endisadelimmlref
%
\isamarkupsubsection{Toplevel transitions%
}
\isamarkuptrue%
%
\begin{isamarkuptext}%
An Isar toplevel transition consists of a partial
  function on the toplevel state, with additional information for
  diagnostics and error reporting: there are fields for command name,
  source position, optional source text, as well as flags for
  interactive-only commands (which issue a warning in batch-mode),
  printing of result state, etc.

  The operational part is represented as a sequential union of a list
  of partial functions, which are tried in turn until the first one
  succeeds (i.e.\ does not raise \verb|Toplevel.UNDEF|).  For example,
  a single Isar command like \isacommand{qed} consists of the union of
  some function \verb|Proof.state -> Proof.state| for proofs
  within proofs, plus \verb|Proof.state -> theory| for proofs at
  the outer theory level.

  Toplevel transitions are composed via transition transformers.
  Internally, Isar commands are put together from an empty transition
  extended by name and source position (and optional source text).  It
  is then left to the individual command parser to turn the given
  syntax body into a suitable transition transformer that adjoin
  actual operations on a theory or proof state etc.%
\end{isamarkuptext}%
\isamarkuptrue%
%
\isadelimmlref
%
\endisadelimmlref
%
\isatagmlref
%
\begin{isamarkuptext}%
\begin{mldecls}
  \indexml{Toplevel.print}\verb|Toplevel.print: Toplevel.transition -> Toplevel.transition| \\
  \indexml{Toplevel.no-timing}\verb|Toplevel.no_timing: Toplevel.transition -> Toplevel.transition| \\
  \indexml{Toplevel.keep}\verb|Toplevel.keep: (Toplevel.state -> unit) ->|\isasep\isanewline%
\verb|  Toplevel.transition -> Toplevel.transition| \\
  \indexml{Toplevel.theory}\verb|Toplevel.theory: (theory -> theory) ->|\isasep\isanewline%
\verb|  Toplevel.transition -> Toplevel.transition| \\
  \indexml{Toplevel.theory-to-proof}\verb|Toplevel.theory_to_proof: (theory -> Proof.state) ->|\isasep\isanewline%
\verb|  Toplevel.transition -> Toplevel.transition| \\
  \indexml{Toplevel.proof}\verb|Toplevel.proof: (Proof.state -> Proof.state) ->|\isasep\isanewline%
\verb|  Toplevel.transition -> Toplevel.transition| \\
  \indexml{Toplevel.proofs}\verb|Toplevel.proofs: (Proof.state -> Proof.state Seq.seq) ->|\isasep\isanewline%
\verb|  Toplevel.transition -> Toplevel.transition| \\
  \indexml{Toplevel.proof-to-theory}\verb|Toplevel.proof_to_theory: (Proof.state -> theory) ->|\isasep\isanewline%
\verb|  Toplevel.transition -> Toplevel.transition| \\
  \end{mldecls}

  \begin{description}

  \item \verb|Toplevel.print| sets the print flag, which causes the
  resulting state of the transition to be echoed in interactive mode.

  \item \verb|Toplevel.no_timing| indicates that the transition should
  never show timing information, e.g.\ because it is merely a
  diagnostic command.

  \item \verb|Toplevel.keep| adjoins a diagnostic function.

  \item \verb|Toplevel.theory| adjoins a theory transformer.

  \item \verb|Toplevel.theory_to_proof| adjoins a global goal function,
  which turns a theory into a proof state.  The theory may be changed
  before entering the proof; the generic Isar goal setup includes an
  argument that specifies how to apply the proven result to the
  theory, when the proof is finished.

  \item \verb|Toplevel.proof| adjoins a deterministic proof command,
  with a singleton result state.

  \item \verb|Toplevel.proofs| adjoins a general proof command, with
  zero or more result states (represented as a lazy list).

  \item \verb|Toplevel.proof_to_theory| adjoins a concluding proof
  command, that returns the resulting theory, after storing the
  resulting facts etc.

  \end{description}%
\end{isamarkuptext}%
\isamarkuptrue%
%
\endisatagmlref
{\isafoldmlref}%
%
\isadelimmlref
%
\endisadelimmlref
%
\isamarkupsubsection{Toplevel control%
}
\isamarkuptrue%
%
\begin{isamarkuptext}%
Apart from regular toplevel transactions there are a few
  special control commands that modify the behavior the toplevel
  itself, and only make sense in interactive mode.  Under normal
  circumstances, the user encounters these only implicitly as part of
  the protocol between the Isabelle/Isar system and a user-interface
  such as ProofGeneral.

  \begin{description}

  \item \isacommand{undo} follows the three-level hierarchy of empty
  toplevel vs.\ theory vs.\ proof: undo within a proof reverts to the
  previous proof context, undo after a proof reverts to the theory
  before the initial goal statement, undo of a theory command reverts
  to the previous theory value, undo of a theory header discontinues
  the current theory development and removes it from the theory
  database (\secref{sec:theory-database}).

  \item \isacommand{kill} aborts the current level of development:
  kill in a proof context reverts to the theory before the initial
  goal statement, kill in a theory context aborts the current theory
  development, removing it from the database.

  \item \isacommand{exit} drops out of the Isar toplevel into the
  underlying {\ML} toplevel (\secref{sec:ML-toplevel}).  The Isar
  toplevel state is preserved and may be continued later.

  \item \isacommand{quit} terminates the Isabelle/Isar process without
  saving.

  \end{description}%
\end{isamarkuptext}%
\isamarkuptrue%
%
\isamarkupsection{ML toplevel \label{sec:ML-toplevel}%
}
\isamarkuptrue%
%
\begin{isamarkuptext}%
The {\ML} toplevel provides a read-compile-eval-print loop for
  {\ML} values, types, structures, and functors.  {\ML} declarations
  operate on the global system state, which consists of the compiler
  environment plus the values of {\ML} reference variables.  There is
  no clean way to undo {\ML} declarations, except for reverting to a
  previously saved state of the whole Isabelle process.  {\ML} input
  is either read interactively from a TTY, or from a string (usually
  within a theory text), or from a source file (usually associated
  with a theory).

  Whenever the {\ML} toplevel is active, the current Isabelle theory
  context is passed as an internal reference variable.  Thus {\ML}
  code may access the theory context during compilation, it may even
  change the value of a theory being under construction --- following
  the usual linearity restrictions (cf.~\secref{sec:context-theory}).%
\end{isamarkuptext}%
\isamarkuptrue%
%
\isadelimmlref
%
\endisadelimmlref
%
\isatagmlref
%
\begin{isamarkuptext}%
\begin{mldecls}
  \indexml{context}\verb|context: theory -> unit| \\
  \indexml{the-context}\verb|the_context: unit -> theory| \\
  \indexml{Context.$>$$>$ }\verb|Context.>> : (theory -> theory) -> unit| \\
  \end{mldecls}

  \begin{description}

  \item \verb|context|~\isa{thy} sets the {\ML} theory context to
  \isa{thy}.  This is usually performed automatically by the system,
  when dropping out of the interactive Isar toplevel into {\ML}, or
  when Isar invokes {\ML} to process code from a string or a file.

  \item \verb|the_context ()| refers to the theory context of the
  {\ML} toplevel --- at compile time!  {\ML} code needs to take care
  to refer to \verb|the_context ()| correctly, recall that evaluation
  of a function body is delayed until actual runtime.  Moreover,
  persistent {\ML} toplevel bindings to an unfinished theory should be
  avoided: code should either project out the desired information
  immediately, or produce an explicit \verb|theory_ref| (cf.\
  \secref{sec:context-theory}).

  \item \verb|Context.>>|~\isa{f} applies theory transformation
  \isa{f} to the current theory of the {\ML} toplevel.  In order to
  work as expected, the theory should be still under construction, and
  the Isar language element that invoked the {\ML} compiler in the
  first place shoule be ready to accept the changed theory value
  (e.g.\ \isakeyword{ML-setup}, but not plain \isakeyword{ML}).
  Otherwise the theory may get destroyed!

  \end{description}

  It is very important to note that the above functions are really
  restricted to the compile time, even though the {\ML} compiler is
  invoked at runtime!  The majority of {\ML} code uses explicit
  functional arguments of a theory or proof context, as required.
  Thus it may get run in an arbitrary context later on.

  \bigskip

  \begin{mldecls}
  \indexml{Isar.main}\verb|Isar.main: unit -> unit| \\
  \indexml{Isar.loop}\verb|Isar.loop: unit -> unit| \\
  \indexml{Isar.state}\verb|Isar.state: unit -> Toplevel.state| \\
  \indexml{Isar.context}\verb|Isar.context: unit -> Proof.context| \\
  \indexml{Isar.exn}\verb|Isar.exn: unit -> (exn * string) option| \\
  \end{mldecls}

  \begin{description}

  \item \verb|Isar.main ()| invokes the Isar toplevel from {\ML},
  initializing the state to empty toplevel state.

  \item \verb|Isar.loop ()| continues the Isar toplevel with the
  current state, after dropping out of the Isar toplevel loop.

  \item \verb|Isar.state ()| and \verb|Isar.exn ()| get current
  toplevel state and optional error condition, respectively.  This
  only works after dropping out of the Isar toplevel loop.

  \item \verb|Isar.context ()| produces the proof context from \verb|Isar.state ()|.

  \end{description}%
\end{isamarkuptext}%
\isamarkuptrue%
%
\endisatagmlref
{\isafoldmlref}%
%
\isadelimmlref
%
\endisadelimmlref
%
\isamarkupsection{Theory database%
}
\isamarkuptrue%
%
\begin{isamarkuptext}%
The theory database maintains a collection of theories,
  together with some administrative information about the original
  sources, which are held in an external store (i.e.\ a collection of
  directories within the regular file system of the underlying
  platform).

  The theory database is organized as a directed acyclic graph, with
  entries referenced by theory name.  Although some external
  interfaces allow to include a directory specification, this is only
  a hint to the underlying theory loader mechanism: the internal
  theory name space is flat.

  Theory \isa{A} is associated with the main theory file \isa{A}\verb,.thy,, which needs to be accessible through the theory
  loader path.  A number of optional {\ML} source files may be
  associated with each theory, by declaring these dependencies in the
  theory header as \isa{{\isasymUSES}}, and loading them consecutively
  within the theory context.  The system keeps track of incoming {\ML}
  sources and associates them with the current theory.  The special
  theory {\ML} file \isa{A}\verb,.ML, is loaded after a theory has
  been concluded, in order to support legacy proof {\ML} proof
  scripts.

  The basic internal actions of the theory database are \isa{update}, \isa{outdate}, and \isa{remove}:

  \begin{itemize}

  \item \isa{update\ A} introduces a link of \isa{A} with a
  \isa{theory} value of the same name; it asserts that the theory
  sources are consistent with that value.

  \item \isa{outdate\ A} invalidates the link of a theory database
  entry to its sources, but retains the present theory value.

  \item \isa{remove\ A} removes entry \isa{A} from the theory
  database.
  
  \end{itemize}

  These actions are propagated to sub- or super-graphs of a theory
  entry in the usual way, in order to preserve global consistency of
  the state of all loaded theories with the sources of the external
  store.  This implies causal dependencies of certain actions: \isa{update} or \isa{outdate} of an entry will \isa{outdate}
  all descendants; \isa{remove} will \isa{remove} all
  descendants.

  \medskip There are separate user-level interfaces to operate on the
  theory database directly or indirectly.  The primitive actions then
  just happen automatically while working with the system.  In
  particular, processing a theory header \isa{{\isasymTHEORY}\ A\ {\isasymIMPORTS}\ B\isactrlsub {\isadigit{1}}\ {\isasymdots}\ B\isactrlsub n\ {\isasymBEGIN}} ensure that the
  sub-graph of the collective imports \isa{B\isactrlsub {\isadigit{1}}\ {\isasymdots}\ B\isactrlsub n}
  is up-to-date.  Earlier theories are reloaded as required, with
  \isa{update} actions proceeding in topological order according to
  theory dependencies.  There may be also a wave of implied \isa{outdate} actions for derived theory nodes until a stable situation
  is achieved eventually.%
\end{isamarkuptext}%
\isamarkuptrue%
%
\isadelimmlref
%
\endisadelimmlref
%
\isatagmlref
%
\begin{isamarkuptext}%
\begin{mldecls}
  \indexml{theory}\verb|theory: string -> theory| \\
  \indexml{use-thy}\verb|use_thy: string -> unit| \\
  \indexml{update-thy}\verb|update_thy: string -> unit| \\
  \indexml{use-thy-only}\verb|use_thy_only: string -> unit| \\
  \indexml{update-thy-only}\verb|update_thy_only: string -> unit| \\
  \indexml{touch-thy}\verb|touch_thy: string -> unit| \\
  \indexml{remove-thy}\verb|remove_thy: string -> unit| \\[1ex]
  \indexml{ThyInfo.begin-theory}\verb|ThyInfo.begin_theory|\verb|: ... -> bool -> theory| \\
  \indexml{ThyInfo.end-theory}\verb|ThyInfo.end_theory: theory -> theory| \\
  \indexml{ThyInfo.register-theory}\verb|ThyInfo.register_theory: theory -> unit| \\[1ex]
  \verb|datatype action = Update |\verb,|,\verb| Outdate |\verb,|,\verb| Remove| \\
  \indexml{ThyInfo.add-hook}\verb|ThyInfo.add_hook: (ThyInfo.action -> string -> unit) -> unit| \\
  \end{mldecls}

  \begin{description}

  \item \verb|theory|~\isa{A} retrieves the theory value presently
  associated with \isa{A}.  The result is not necessarily
  up-to-date!

  \item \verb|use_thy|~\isa{A} loads theory \isa{A} if it is absent
  or out-of-date.  It ensures that all parent theories are available
  as well, but does not reload them if older versions are already
  present.

  \item \verb|update_thy| is similar to \verb|use_thy|, but ensures that
  the \isa{A} and all of its ancestors are fully up-to-date.

  \item \verb|use_thy_only|~\isa{A} is like \verb|use_thy|~\isa{A},
  but refrains from loading the attached \isa{A}\verb,.ML, file.
  This is occasionally useful in replaying legacy {\ML} proof scripts
  by hand.
  
  \item \verb|update_thy_only| is analogous to \verb|use_thy_only|, but
  proceeds like \verb|update_thy| for ancestors.

  \item \verb|touch_thy|~\isa{A} performs \isa{outdate} action on
  theory \isa{A} and all of its descendants.

  \item \verb|remove_thy|~\isa{A} removes \isa{A} and all of its
  descendants from the theory database.

  \item \verb|ThyInfo.begin_theory| is the basic operation behind a
  \isa{{\isasymTHEORY}} header declaration.  The boolean argument
  indicates the strictness of treating ancestors: for \verb|true| (as
  in interactive mode) like \verb|update_thy|, and for \verb|false| (as
  in batch mode) like \verb|use_thy|.  This is {\ML} functions is
  normally not invoked directly.

  \item \verb|ThyInfo.end_theory| concludes the loading of a theory
  proper; an attached theory {\ML} file may be still loaded later on.
  This is {\ML} functions is normally not invoked directly.

  \item \verb|ThyInfo.register_theory|~{text thy} registers an existing
  theory value with the theory loader database.  There is no
  management of associated sources; this is mainly for bootstrapping.

  \item \verb|ThyInfo.add_hook|~\isa{f} registers function \isa{f} as a hook for theory database actions.  The function will be
  invoked with the action and theory name being involved; thus derived
  actions may be performed in associated system components, e.g.\
  maintaining the state of an editor for theory sources.

  The kind and order of actions occurring in practice depends both on
  user interactions and the internal process of resolving theory
  imports.  Hooks should not rely on a particular policy here!  Any
  exceptions raised by the hook are ignored by the theory database.

  \end{description}%
\end{isamarkuptext}%
\isamarkuptrue%
%
\endisatagmlref
{\isafoldmlref}%
%
\isadelimmlref
%
\endisadelimmlref
%
\isadelimtheory
%
\endisadelimtheory
%
\isatagtheory
\isacommand{end}\isamarkupfalse%
%
\endisatagtheory
{\isafoldtheory}%
%
\isadelimtheory
%
\endisadelimtheory
\isanewline
\end{isabellebody}%
%%% Local Variables:
%%% mode: latex
%%% TeX-master: "root"
%%% End:


\appendix
%
\begin{isabellebody}%
\def\isabellecontext{ML}%
%
\isadelimtheory
%
\endisadelimtheory
%
\isatagtheory
\isacommand{theory}\isamarkupfalse%
\ {\isachardoublequoteopen}ML{\isachardoublequoteclose}\isanewline
\isakeyword{imports}\ Base\isanewline
\isakeyword{begin}%
\endisatagtheory
{\isafoldtheory}%
%
\isadelimtheory
%
\endisadelimtheory
%
\isamarkupchapter{Isabelle/ML%
}
\isamarkuptrue%
%
\begin{isamarkuptext}%
Isabelle/ML is best understood as a certain culture based on
  Standard ML.  Thus it is not a new programming language, but a
  certain way to use SML at an advanced level within the Isabelle
  environment.  This covers a variety of aspects that are geared
  towards an efficient and robust platform for applications of formal
  logic with fully foundational proof construction --- according to
  the well-known \emph{LCF principle}.  There is specific
  infrastructure with library modules to address the needs of this
  difficult task.  For example, the raw parallel programming model of
  Poly/ML is presented as considerably more abstract concept of
  \emph{future values}, which is then used to augment the inference
  kernel, proof interpreter, and theory loader accordingly.

  The main aspects of Isabelle/ML are introduced below.  These
  first-hand explanations should help to understand how proper
  Isabelle/ML is to be read and written, and to get access to the
  wealth of experience that is expressed in the source text and its
  history of changes.\footnote{See
  \url{http://isabelle.in.tum.de/repos/isabelle} for the full
  Mercurial history.  There are symbolic tags to refer to official
  Isabelle releases, as opposed to arbitrary \emph{tip} versions that
  merely reflect snapshots that are never really up-to-date.}%
\end{isamarkuptext}%
\isamarkuptrue%
%
\isamarkupsection{Style and orthography%
}
\isamarkuptrue%
%
\begin{isamarkuptext}%
The sources of Isabelle/Isar are optimized for
  \emph{readability} and \emph{maintainability}.  The main purpose is
  to tell an informed reader what is really going on and how things
  really work.  This is a non-trivial aim, but it is supported by a
  certain style of writing Isabelle/ML that has emerged from long
  years of system development.\footnote{See also the interesting style
  guide for OCaml
  \url{http://caml.inria.fr/resources/doc/guides/guidelines.en.html}
  which shares many of our means and ends.}

  The main principle behind any coding style is \emph{consistency}.
  For a single author of a small program this merely means ``choose
  your style and stick to it''.  A complex project like Isabelle, with
  long years of development and different contributors, requires more
  standardization.  A coding style that is changed every few years or
  with every new contributor is no style at all, because consistency
  is quickly lost.  Global consistency is hard to achieve, though.
  Nonetheless, one should always strive at least for local consistency
  of modules and sub-systems, without deviating from some general
  principles how to write Isabelle/ML.

  In a sense, good coding style is like an \emph{orthography} for the
  sources: it helps to read quickly over the text and see through the
  main points, without getting distracted by accidental presentation
  of free-style code.%
\end{isamarkuptext}%
\isamarkuptrue%
%
\isamarkupsubsection{Header and sectioning%
}
\isamarkuptrue%
%
\begin{isamarkuptext}%
Isabelle source files have a certain standardized header
  format (with precise spacing) that follows ancient traditions
  reaching back to the earliest versions of the system by Larry
  Paulson.  See \hyperlink{file.~~/src/Pure/thm.ML}{\mbox{\isa{\isatt{{\isachartilde}{\isachartilde}{\isacharslash}src{\isacharslash}Pure{\isacharslash}thm{\isachardot}ML}}}}, for example.

  The header includes at least \verb|Title| and \verb|Author| entries, followed by a prose description of the purpose of
  the module.  The latter can range from a single line to several
  paragraphs of explanations.

  The rest of the file is divided into sections, subsections,
  subsubsections, paragraphs etc.\ using a simple layout via ML
  comments as follows.

\begin{verbatim}
(*** section ***)

(** subsection **)

(* subsubsection *)

(*short paragraph*)

(*
  long paragraph,
  with more text
*)
\end{verbatim}

  As in regular typography, there is some extra space \emph{before}
  section headings that are adjacent to plain text (not other headings
  as in the example above).

  \medskip The precise wording of the prose text given in these
  headings is chosen carefully to introduce the main theme of the
  subsequent formal ML text.%
\end{isamarkuptext}%
\isamarkuptrue%
%
\isamarkupsubsection{Naming conventions%
}
\isamarkuptrue%
%
\begin{isamarkuptext}%
Since ML is the primary medium to express the meaning of the
  source text, naming of ML entities requires special care.

  \paragraph{Notation.}  A name consists of 1--3 \emph{words} (rarely
  4, but not more) that are separated by underscore.  There are three
  variants concerning upper or lower case letters, which are used for
  certain ML categories as follows:

  \medskip
  \begin{tabular}{lll}
  variant & example & ML categories \\\hline
  lower-case & \verb|foo_bar| & values, types, record fields \\
  capitalized & \verb|Foo_Bar| & datatype constructors, structures, functors \\
  upper-case & \verb|FOO_BAR| & special values, exception constructors, signatures \\
  \end{tabular}
  \medskip

  For historical reasons, many capitalized names omit underscores,
  e.g.\ old-style \verb|FooBar| instead of \verb|Foo_Bar|.
  Genuine mixed-case names are \emph{not} used, bacause clear division
  of words is essential for readability.\footnote{Camel-case was
  invented to workaround the lack of underscore in some early
  non-ASCII character sets.  Later it became habitual in some language
  communities that are now strong in numbers.}

  A single (capital) character does not count as ``word'' in this
  respect: some Isabelle/ML names are suffixed by extra markers like
  this: \verb|foo_barT|.

  Name variants are produced by adding 1--3 primes, e.g.\ \verb|foo'|, \verb|foo''|, or \verb|foo'''|, but not \verb|foo''''| or more.  Decimal digits scale better to larger numbers,
  e.g.\ \verb|foo0|, \verb|foo1|, \verb|foo42|.

  \paragraph{Scopes.}  Apart from very basic library modules, ML
  structures are not ``opened'', but names are referenced with
  explicit qualification, as in \verb|Syntax.string_of_term| for
  example.  When devising names for structures and their components it
  is important aim at eye-catching compositions of both parts, because
  this is how they are seen in the sources and documentation.  For the
  same reasons, aliases of well-known library functions should be
  avoided.

  Local names of function abstraction or case/let bindings are
  typically shorter, sometimes using only rudiments of ``words'',
  while still avoiding cryptic shorthands.  An auxiliary function
  called \verb|helper|, \verb|aux|, or \verb|f| is
  considered bad style.

  Example:

  \begin{verbatim}
  (* RIGHT *)

  fun print_foo ctxt foo =
    let
      fun print t = ... Syntax.string_of_term ctxt t ...
    in ... end;


  (* RIGHT *)

  fun print_foo ctxt foo =
    let
      val string_of_term = Syntax.string_of_term ctxt;
      fun print t = ... string_of_term t ...
    in ... end;


  (* WRONG *)

  val string_of_term = Syntax.string_of_term;

  fun print_foo ctxt foo =
    let
      fun aux t = ... string_of_term ctxt t ...
    in ... end;

  \end{verbatim}


  \paragraph{Specific conventions.} Here are some specific name forms
  that occur frequently in the sources.

  \begin{itemize}

  \item A function that maps \verb|foo| to \verb|bar| is
  called \verb|foo_to_bar| or \verb|bar_of_foo| (never
  \verb|foo2bar|, \verb|bar_from_foo|, \verb|bar_for_foo|, or \verb|bar4foo|).

  \item The name component \verb|legacy| means that the operation
  is about to be discontinued soon.

  \item The name component \verb|old| means that this is historic
  material that might disappear at some later stage.

  \item The name component \verb|global| means that this works
  with the background theory instead of the regular local context
  (\secref{sec:context}), sometimes for historical reasons, sometimes
  due a genuine lack of locality of the concept involved, sometimes as
  a fall-back for the lack of a proper context in the application
  code.  Whenever there is a non-global variant available, the
  application should be migrated to use it with a proper local
  context.

  \item Variables of the main context types of the Isabelle/Isar
  framework (\secref{sec:context} and \chref{ch:local-theory}) have
  firm naming conventions as follows:

  \begin{itemize}

  \item theories are called \verb|thy|, rarely \verb|theory|
  (never \verb|thry|)

  \item proof contexts are called \verb|ctxt|, rarely \verb|context| (never \verb|ctx|)

  \item generic contexts are called \verb|context|, rarely
  \verb|ctxt|

  \item local theories are called \verb|lthy|, except for local
  theories that are treated as proof context (which is a semantic
  super-type)

  \end{itemize}

  Variations with primed or decimal numbers are always possible, as
  well as sematic prefixes like \verb|foo_thy| or \verb|bar_ctxt|, but the base conventions above need to be preserved.
  This allows to visualize the their data flow via plain regular
  expressions in the editor.

  \item The main logical entities (\secref{ch:logic}) have established
  naming convention as follows:

  \begin{itemize}

  \item sorts are called \verb|S|

  \item types are called \verb|T|, \verb|U|, or \verb|ty| (never \verb|t|)

  \item terms are called \verb|t|, \verb|u|, or \verb|tm| (never \verb|trm|)

  \item certified types are called \verb|cT|, rarely \verb|T|, with variants as for types

  \item certified terms are called \verb|ct|, rarely \verb|t|, with variants as for terms

  \item theorems are called \verb|th|, or \verb|thm|

  \end{itemize}

  Proper semantic names override these conventions completely.  For
  example, the left-hand side of an equation (as a term) can be called
  \verb|lhs| (not \verb|lhs_tm|).  Or a term that is known
  to be a variable can be called \verb|v| or \verb|x|.

  \end{itemize}%
\end{isamarkuptext}%
\isamarkuptrue%
%
\isamarkupsubsection{General source layout%
}
\isamarkuptrue%
%
\begin{isamarkuptext}%
The general Isabelle/ML source layout imitates regular
  type-setting to some extent, augmented by the requirements for
  deeply nested expressions that are commonplace in functional
  programming.

  \paragraph{Line length} is 80 characters according to ancient
  standards, but we allow as much as 100 characters (not
  more).\footnote{Readability requires to keep the beginning of a line
  in view while watching its end.  Modern wide-screen displays do not
  change the way how the human brain works.  Sources also need to be
  printable on plain paper with reasonable font-size.} The extra 20
  characters acknowledge the space requirements due to qualified
  library references in Isabelle/ML.

  \paragraph{White-space} is used to emphasize the structure of
  expressions, following mostly standard conventions for mathematical
  typesetting, as can be seen in plain {\TeX} or {\LaTeX}.  This
  defines positioning of spaces for parentheses, punctuation, and
  infixes as illustrated here:

  \begin{verbatim}
  val x = y + z * (a + b);
  val pair = (a, b);
  val record = {foo = 1, bar = 2};
  \end{verbatim}

  Lines are normally broken \emph{after} an infix operator or
  punctuation character.  For example:

  \begin{verbatim}
  val x =
    a +
    b +
    c;

  val tuple =
   (a,
    b,
    c);
  \end{verbatim}

  Some special infixes (e.g.\ \verb||\verb,|,\verb|>|) work better at the
  start of the line, but punctuation is always at the end.

  Function application follows the tradition of \isa{{\isasymlambda}}-calculus,
  not informal mathematics.  For example: \verb|f a b| for a
  curried function, or \verb|g (a, b)| for a tupled function.
  Note that the space between \verb|g| and the pair \verb|(a, b)| follows the important principle of
  \emph{compositionality}: the layout of \verb|g p| does not
  change when \verb|p| is refined to the concrete pair
  \verb|(a, b)|.

  \paragraph{Indentation} uses plain spaces, never hard
  tabulators.\footnote{Tabulators were invented to move the carriage
  of a type-writer to certain predefined positions.  In software they
  could be used as a primitive run-length compression of consecutive
  spaces, but the precise result would depend on non-standardized
  editor configuration.}

  Each level of nesting is indented by 2 spaces, sometimes 1, very
  rarely 4, never 8 or any other odd number.

  Indentation follows a simple logical format that only depends on the
  nesting depth, not the accidental length of the text that initiates
  a level of nesting.  Example:

  \begin{verbatim}
  (* RIGHT *)

  if b then
    expr1_part1
    expr1_part2
  else
    expr2_part1
    expr2_part2


  (* WRONG *)

  if b then expr1_part1
            expr1_part2
  else expr2_part1
       expr2_part2
  \end{verbatim}

  The second form has many problems: it assumes a fixed-width font
  when viewing the sources, it uses more space on the line and thus
  makes it hard to observe its strict length limit (working against
  \emph{readability}), it requires extra editing to adapt the layout
  to changes of the initial text (working against
  \emph{maintainability}) etc.

  \medskip For similar reasons, any kind of two-dimensional or tabular
  layouts, ASCII-art with lines or boxes of asterisks etc.\ should be
  avoided.

  \paragraph{Complex expressions} that consist of multi-clausal
  function definitions, \verb|handle|, \verb|case|,
  \verb|let| (and combinations) require special attention.  The
  syntax of Standard ML is quite ambitious and admits a lot of
  variance that can distort the meaning of the text.

  Clauses of \verb|fun|, \verb|fn|, \verb|handle|,
  \verb|case| get extra indentation to indicate the nesting
  clearly.  Example:

  \begin{verbatim}
  (* RIGHT *)

  fun foo p1 =
        expr1
    | foo p2 =
        expr2


  (* WRONG *)

  fun foo p1 =
    expr1
    | foo p2 =
    expr2
  \end{verbatim}

  Body expressions consisting of \verb|case| or \verb|let|
  require care to maintain compositionality, to prevent loss of
  logical indentation where it is especially important to see the
  structure of the text.  Example:

  \begin{verbatim}
  (* RIGHT *)

  fun foo p1 =
        (case e of
          q1 => ...
        | q2 => ...)
    | foo p2 =
        let
          ...
        in
          ...
        end


  (* WRONG *)

  fun foo p1 = case e of
      q1 => ...
    | q2 => ...
    | foo p2 =
    let
      ...
    in
      ...
    end
  \end{verbatim}

  Extra parentheses around \verb|case| expressions are optional,
  but help to analyse the nesting based on character matching in the
  editor.

  \medskip There are two main exceptions to the overall principle of
  compositionality in the layout of complex expressions.

  \begin{enumerate}

  \item \verb|if| expressions are iterated as if there would be
  a multi-branch conditional in SML, e.g.

  \begin{verbatim}
  (* RIGHT *)

  if b1 then e1
  else if b2 then e2
  else e3
  \end{verbatim}

  \item \verb|fn| abstractions are often layed-out as if they
  would lack any structure by themselves.  This traditional form is
  motivated by the possibility to shift function arguments back and
  forth wrt.\ additional combinators.  Example:

  \begin{verbatim}
  (* RIGHT *)

  fun foo x y = fold (fn z =>
    expr)
  \end{verbatim}

  Here the visual appearance is that of three arguments \verb|x|,
  \verb|y|, \verb|z|.

  \end{enumerate}

  Such weakly structured layout should be use with great care.  Here
  is a counter-example involving \verb|let| expressions:

  \begin{verbatim}
  (* WRONG *)

  fun foo x = let
      val y = ...
    in ... end

  fun foo x =
  let
    val y = ...
  in ... end
  \end{verbatim}

  \medskip In general the source layout is meant to emphasize the
  structure of complex language expressions, not to pretend that SML
  had a completely different syntax (say that of Haskell or Java).%
\end{isamarkuptext}%
\isamarkuptrue%
%
\isamarkupsection{SML embedded into Isabelle/Isar%
}
\isamarkuptrue%
%
\begin{isamarkuptext}%
ML and Isar are intertwined via an open-ended bootstrap
  process that provides more and more programming facilities and
  logical content in an alternating manner.  Bootstrapping starts from
  the raw environment of existing implementations of Standard ML
  (mainly Poly/ML, but also SML/NJ).

  Isabelle/Pure marks the point where the original ML toplevel is
  superseded by the Isar toplevel that maintains a uniform context for
  arbitrary ML values (see also \secref{sec:context}).  This formal
  environment holds ML compiler bindings, logical entities, and many
  other things.  Raw SML is never encountered again after the initial
  bootstrap of Isabelle/Pure.

  Object-logics like Isabelle/HOL are built within the
  Isabelle/ML/Isar environment by introducing suitable theories with
  associated ML modules, either inlined or as separate files.  Thus
  Isabelle/HOL is defined as a regular user-space application within
  the Isabelle framework.  Further add-on tools can be implemented in
  ML within the Isar context in the same manner: ML is part of the
  standard repertoire of Isabelle, and there is no distinction between
  ``user'' and ``developer'' in this respect.%
\end{isamarkuptext}%
\isamarkuptrue%
%
\isamarkupsubsection{Isar ML commands%
}
\isamarkuptrue%
%
\begin{isamarkuptext}%
The primary Isar source language provides facilities to ``open
  a window'' to the underlying ML compiler.  Especially see the Isar
  commands \indexref{}{command}{use}\hyperlink{command.use}{\mbox{\isa{\isacommand{use}}}} and \indexref{}{command}{ML}\hyperlink{command.ML}{\mbox{\isa{\isacommand{ML}}}}: both work the
  same way, only the source text is provided via a file vs.\ inlined,
  respectively.  Apart from embedding ML into the main theory
  definition like that, there are many more commands that refer to ML
  source, such as \indexref{}{command}{setup}\hyperlink{command.setup}{\mbox{\isa{\isacommand{setup}}}} or \indexref{}{command}{declaration}\hyperlink{command.declaration}{\mbox{\isa{\isacommand{declaration}}}}.
  Even more fine-grained embedding of ML into Isar is encountered in
  the proof method \indexref{}{method}{tactic}\hyperlink{method.tactic}{\mbox{\isa{tactic}}}, which refines the pending
  goal state via a given expression of type \verb|tactic|.%
\end{isamarkuptext}%
\isamarkuptrue%
%
\isadelimmlex
%
\endisadelimmlex
%
\isatagmlex
%
\begin{isamarkuptext}%
The following artificial example demonstrates some ML
  toplevel declarations within the implicit Isar theory context.  This
  is regular functional programming without referring to logical
  entities yet.%
\end{isamarkuptext}%
\isamarkuptrue%
%
\endisatagmlex
{\isafoldmlex}%
%
\isadelimmlex
%
\endisadelimmlex
%
\isadelimML
%
\endisadelimML
%
\isatagML
\isacommand{ML}\isamarkupfalse%
\ {\isacharverbatimopen}\isanewline
\ \ fun\ factorial\ {\isadigit{0}}\ {\isacharequal}\ {\isadigit{1}}\isanewline
\ \ \ \ {\isacharbar}\ factorial\ n\ {\isacharequal}\ n\ {\isacharasterisk}\ factorial\ {\isacharparenleft}n\ {\isacharminus}\ {\isadigit{1}}{\isacharparenright}\isanewline
{\isacharverbatimclose}%
\endisatagML
{\isafoldML}%
%
\isadelimML
%
\endisadelimML
%
\begin{isamarkuptext}%
Here the ML environment is already managed by Isabelle, i.e.\
  the \verb|factorial| function is not yet accessible in the preceding
  paragraph, nor in a different theory that is independent from the
  current one in the import hierarchy.

  Removing the above ML declaration from the source text will remove
  any trace of this definition as expected.  The Isabelle/ML toplevel
  environment is managed in a \emph{stateless} way: unlike the raw ML
  toplevel there are no global side-effects involved
  here.\footnote{Such a stateless compilation environment is also a
  prerequisite for robust parallel compilation within independent
  nodes of the implicit theory development graph.}

  \medskip The next example shows how to embed ML into Isar proofs, using
 \indexref{}{command}{ML\_prf}\hyperlink{command.ML-prf}{\mbox{\isa{\isacommand{ML{\isacharunderscore}prf}}}} instead of Instead of \indexref{}{command}{ML}\hyperlink{command.ML}{\mbox{\isa{\isacommand{ML}}}}.
  As illustrated below, the effect on the ML environment is local to
  the whole proof body, ignoring the block structure.%
\end{isamarkuptext}%
\isamarkuptrue%
\isacommand{example{\isacharunderscore}proof}\isamarkupfalse%
\isanewline
%
\isadelimML
\ \ %
\endisadelimML
%
\isatagML
\isacommand{ML{\isacharunderscore}prf}\isamarkupfalse%
\ {\isacharverbatimopen}\ val\ a\ {\isacharequal}\ {\isadigit{1}}\ {\isacharverbatimclose}\isanewline
\ \ \isacommand{{\isacharbraceleft}}\isamarkupfalse%
\isanewline
\ \ \ \ \isacommand{ML{\isacharunderscore}prf}\isamarkupfalse%
\ {\isacharverbatimopen}\ val\ b\ {\isacharequal}\ a\ {\isacharplus}\ {\isadigit{1}}\ {\isacharverbatimclose}\isanewline
\ \ \isacommand{{\isacharbraceright}}\isamarkupfalse%
\ %
\isamarkupcmt{Isar block structure ignored by ML environment%
}
\isanewline
\ \ \isacommand{ML{\isacharunderscore}prf}\isamarkupfalse%
\ {\isacharverbatimopen}\ val\ c\ {\isacharequal}\ b\ {\isacharplus}\ {\isadigit{1}}\ {\isacharverbatimclose}\isanewline
\isacommand{qed}\isamarkupfalse%
%
\endisatagML
{\isafoldML}%
%
\isadelimML
%
\endisadelimML
%
\begin{isamarkuptext}%
By side-stepping the normal scoping rules for Isar proof
  blocks, embedded ML code can refer to the different contexts and
  manipulate corresponding entities, e.g.\ export a fact from a block
  context.

  \medskip Two further ML commands are useful in certain situations:
  \indexref{}{command}{ML\_val}\hyperlink{command.ML-val}{\mbox{\isa{\isacommand{ML{\isacharunderscore}val}}}} and \indexref{}{command}{ML\_command}\hyperlink{command.ML-command}{\mbox{\isa{\isacommand{ML{\isacharunderscore}command}}}} are
  \emph{diagnostic} in the sense that there is no effect on the
  underlying environment, and can thus used anywhere (even outside a
  theory).  The examples below produce long strings of digits by
  invoking \verb|factorial|: \hyperlink{command.ML-val}{\mbox{\isa{\isacommand{ML{\isacharunderscore}val}}}} already takes care of
  printing the ML toplevel result, but \hyperlink{command.ML-command}{\mbox{\isa{\isacommand{ML{\isacharunderscore}command}}}} is silent
  so we produce an explicit output message.%
\end{isamarkuptext}%
\isamarkuptrue%
%
\isadelimML
%
\endisadelimML
%
\isatagML
\isacommand{ML{\isacharunderscore}val}\isamarkupfalse%
\ {\isacharverbatimopen}\ factorial\ {\isadigit{1}}{\isadigit{0}}{\isadigit{0}}\ {\isacharverbatimclose}\isanewline
\isacommand{ML{\isacharunderscore}command}\isamarkupfalse%
\ {\isacharverbatimopen}\ writeln\ {\isacharparenleft}string{\isacharunderscore}of{\isacharunderscore}int\ {\isacharparenleft}factorial\ {\isadigit{1}}{\isadigit{0}}{\isadigit{0}}{\isacharparenright}{\isacharparenright}\ {\isacharverbatimclose}%
\endisatagML
{\isafoldML}%
%
\isadelimML
%
\endisadelimML
\isanewline
\isanewline
\isacommand{example{\isacharunderscore}proof}\isamarkupfalse%
\isanewline
%
\isadelimML
\ \ %
\endisadelimML
%
\isatagML
\isacommand{ML{\isacharunderscore}val}\isamarkupfalse%
\ {\isacharverbatimopen}\ factorial\ {\isadigit{1}}{\isadigit{0}}{\isadigit{0}}\ {\isacharverbatimclose}\ \ \isanewline
\ \ \isacommand{ML{\isacharunderscore}command}\isamarkupfalse%
\ {\isacharverbatimopen}\ writeln\ {\isacharparenleft}string{\isacharunderscore}of{\isacharunderscore}int\ {\isacharparenleft}factorial\ {\isadigit{1}}{\isadigit{0}}{\isadigit{0}}{\isacharparenright}{\isacharparenright}\ {\isacharverbatimclose}%
\endisatagML
{\isafoldML}%
%
\isadelimML
\isanewline
%
\endisadelimML
%
\isadelimproof
%
\endisadelimproof
%
\isatagproof
\isacommand{qed}\isamarkupfalse%
%
\endisatagproof
{\isafoldproof}%
%
\isadelimproof
%
\endisadelimproof
%
\isamarkupsubsection{Compile-time context%
}
\isamarkuptrue%
%
\begin{isamarkuptext}%
Whenever the ML compiler is invoked within Isabelle/Isar, the
  formal context is passed as a thread-local reference variable.  Thus
  ML code may access the theory context during compilation, by reading
  or writing the (local) theory under construction.  Note that such
  direct access to the compile-time context is rare.  In practice it
  is typically done via some derived ML functions instead.%
\end{isamarkuptext}%
\isamarkuptrue%
%
\isadelimmlref
%
\endisadelimmlref
%
\isatagmlref
%
\begin{isamarkuptext}%
\begin{mldecls}
  \indexdef{}{ML}{ML\_Context.the\_generic\_context}\verb|ML_Context.the_generic_context: unit -> Context.generic| \\
  \indexdef{}{ML}{Context.$>$$>$}\verb|Context.>> : (Context.generic -> Context.generic) -> unit| \\
  \indexdef{}{ML}{bind\_thms}\verb|bind_thms: string * thm list -> unit| \\
  \indexdef{}{ML}{bind\_thm}\verb|bind_thm: string * thm -> unit| \\
  \end{mldecls}

  \begin{description}

  \item \verb|ML_Context.the_generic_context ()| refers to the theory
  context of the ML toplevel --- at compile time.  ML code needs to
  take care to refer to \verb|ML_Context.the_generic_context ()|
  correctly.  Recall that evaluation of a function body is delayed
  until actual run-time.

  \item \verb|Context.>>|~\isa{f} applies context transformation
  \isa{f} to the implicit context of the ML toplevel.

  \item \verb|bind_thms|~\isa{{\isacharparenleft}name{\isacharcomma}\ thms{\isacharparenright}} stores a list of
  theorems produced in ML both in the (global) theory context and the
  ML toplevel, associating it with the provided name.  Theorems are
  put into a global ``standard'' format before being stored.

  \item \verb|bind_thm| is similar to \verb|bind_thms| but refers to a
  singleton fact.

  \end{description}

  It is important to note that the above functions are really
  restricted to the compile time, even though the ML compiler is
  invoked at run-time.  The majority of ML code either uses static
  antiquotations (\secref{sec:ML-antiq}) or refers to the theory or
  proof context at run-time, by explicit functional abstraction.%
\end{isamarkuptext}%
\isamarkuptrue%
%
\endisatagmlref
{\isafoldmlref}%
%
\isadelimmlref
%
\endisadelimmlref
%
\isamarkupsubsection{Antiquotations \label{sec:ML-antiq}%
}
\isamarkuptrue%
%
\begin{isamarkuptext}%
A very important consequence of embedding SML into Isar is the
  concept of \emph{ML antiquotation}.  The standard token language of
  ML is augmented by special syntactic entities of the following form:

  \indexouternonterm{antiquote}
  \begin{rail}
  antiquote: atsign lbrace nameref args rbrace | lbracesym | rbracesym
  ;
  \end{rail}

  Here \hyperlink{syntax.nameref}{\mbox{\isa{nameref}}} and \hyperlink{syntax.args}{\mbox{\isa{args}}} are regular outer syntax
  categories \cite{isabelle-isar-ref}.  Attributes and proof methods
  use similar syntax.

  \medskip A regular antiquotation \isa{{\isacharat}{\isacharbraceleft}name\ args{\isacharbraceright}} processes
  its arguments by the usual means of the Isar source language, and
  produces corresponding ML source text, either as literal
  \emph{inline} text (e.g. \isa{{\isacharat}{\isacharbraceleft}term\ t{\isacharbraceright}}) or abstract
  \emph{value} (e.g. \isa{{\isacharat}{\isacharbraceleft}thm\ th{\isacharbraceright}}).  This pre-compilation
  scheme allows to refer to formal entities in a robust manner, with
  proper static scoping and with some degree of logical checking of
  small portions of the code.

  Special antiquotations like \isa{{\isacharat}{\isacharbraceleft}let\ {\isasymdots}{\isacharbraceright}} or \isa{{\isacharat}{\isacharbraceleft}note\ {\isasymdots}{\isacharbraceright}} augment the compilation context without generating code.  The
  non-ASCII braces \isa{{\isasymlbrace}} and \isa{{\isasymrbrace}} allow to delimit the
  effect by introducing local blocks within the pre-compilation
  environment.

  \medskip See also \cite{Wenzel-Chaieb:2007b} for a broader
  perspective on Isabelle/ML antiquotations.%
\end{isamarkuptext}%
\isamarkuptrue%
%
\isadelimmlantiq
%
\endisadelimmlantiq
%
\isatagmlantiq
%
\begin{isamarkuptext}%
\begin{matharray}{rcl}
  \indexdef{}{ML antiquotation}{let}\hypertarget{ML antiquotation.let}{\hyperlink{ML antiquotation.let}{\mbox{\isa{let}}}} & : & \isa{ML{\isacharunderscore}antiquotation} \\
  \indexdef{}{ML antiquotation}{note}\hypertarget{ML antiquotation.note}{\hyperlink{ML antiquotation.note}{\mbox{\isa{note}}}} & : & \isa{ML{\isacharunderscore}antiquotation} \\
  \end{matharray}

  \begin{rail}
  'let' ((term + 'and') '=' term + 'and')
  ;

  'note' (thmdef? thmrefs + 'and')
  ;
  \end{rail}

  \begin{description}

  \item \isa{{\isacharat}{\isacharbraceleft}let\ p\ {\isacharequal}\ t{\isacharbraceright}} binds schematic variables in the
  pattern \isa{p} by higher-order matching against the term \isa{t}.  This is analogous to the regular \indexref{}{command}{let}\hyperlink{command.let}{\mbox{\isa{\isacommand{let}}}} command
  in the Isar proof language.  The pre-compilation environment is
  augmented by auxiliary term bindings, without emitting ML source.

  \item \isa{{\isacharat}{\isacharbraceleft}note\ a\ {\isacharequal}\ b\isactrlsub {\isadigit{1}}\ {\isasymdots}\ b\isactrlsub n{\isacharbraceright}} recalls existing facts \isa{b\isactrlsub {\isadigit{1}}{\isacharcomma}\ {\isasymdots}{\isacharcomma}\ b\isactrlsub n}, binding the result as \isa{a}.  This is analogous to
  the regular \indexref{}{command}{note}\hyperlink{command.note}{\mbox{\isa{\isacommand{note}}}} command in the Isar proof language.
  The pre-compilation environment is augmented by auxiliary fact
  bindings, without emitting ML source.

  \end{description}%
\end{isamarkuptext}%
\isamarkuptrue%
%
\endisatagmlantiq
{\isafoldmlantiq}%
%
\isadelimmlantiq
%
\endisadelimmlantiq
%
\isadelimmlex
%
\endisadelimmlex
%
\isatagmlex
%
\begin{isamarkuptext}%
The following artificial example gives some impression
  about the antiquotation elements introduced so far, together with
  the important \isa{{\isacharat}{\isacharbraceleft}thm{\isacharbraceright}} antiquotation defined later.%
\end{isamarkuptext}%
\isamarkuptrue%
%
\endisatagmlex
{\isafoldmlex}%
%
\isadelimmlex
%
\endisadelimmlex
%
\isadelimML
%
\endisadelimML
%
\isatagML
\isacommand{ML}\isamarkupfalse%
\ {\isacharverbatimopen}\isanewline
\ \ {\isaantiqopen}\isanewline
\ \ \ \ %
\isaantiq
let\ {\isacharquery}t\ {\isacharequal}\ my{\isacharunderscore}term%
\endisaantiq
\isanewline
\ \ \ \ %
\isaantiq
note\ my{\isacharunderscore}refl\ {\isacharequal}\ reflexive\ {\isacharbrackleft}of\ {\isacharquery}t{\isacharbrackright}%
\endisaantiq
\isanewline
\ \ \ \ fun\ foo\ th\ {\isacharequal}\ Thm{\isachardot}transitive\ th\ %
\isaantiq
thm\ my{\isacharunderscore}refl%
\endisaantiq
\isanewline
\ \ {\isaantiqclose}\isanewline
{\isacharverbatimclose}%
\endisatagML
{\isafoldML}%
%
\isadelimML
%
\endisadelimML
%
\begin{isamarkuptext}%
The extra block delimiters do not affect the compiled code
  itself, i.e.\ function \verb|foo| is available in the present context
  of this paragraph.%
\end{isamarkuptext}%
\isamarkuptrue%
%
\isamarkupsection{Canonical argument order \label{sec:canonical-argument-order}%
}
\isamarkuptrue%
%
\begin{isamarkuptext}%
Standard ML is a language in the tradition of \isa{{\isasymlambda}}-calculus and \emph{higher-order functional programming},
  similar to OCaml, Haskell, or Isabelle/Pure and HOL as logical
  languages.  Getting acquainted with the native style of representing
  functions in that setting can save a lot of extra boiler-plate of
  redundant shuffling of arguments, auxiliary abstractions etc.

  Functions are usually \emph{curried}: the idea of turning arguments
  of type \isa{{\isasymtau}\isactrlsub i} (for \isa{i\ {\isasymin}\ {\isacharbraceleft}{\isadigit{1}}{\isacharcomma}\ {\isasymdots}\ n{\isacharbraceright}}) into a result of
  type \isa{{\isasymtau}} is represented by the iterated function space
  \isa{{\isasymtau}\isactrlsub {\isadigit{1}}\ {\isasymrightarrow}\ {\isasymdots}\ {\isasymrightarrow}\ {\isasymtau}\isactrlsub n\ {\isasymrightarrow}\ {\isasymtau}}.  This is isomorphic to the well-known
  encoding via tuples \isa{{\isasymtau}\isactrlsub {\isadigit{1}}\ {\isasymtimes}\ {\isasymdots}\ {\isasymtimes}\ {\isasymtau}\isactrlsub n\ {\isasymrightarrow}\ {\isasymtau}}, but the curried
  version fits more smoothly into the basic calculus.\footnote{The
  difference is even more significant in higher-order logic, because
  the redundant tuple structure needs to be accommodated by formal
  reasoning.}

  Currying gives some flexiblity due to \emph{partial application}.  A
  function \isa{f{\isacharcolon}\ {\isasymtau}\isactrlsub {\isadigit{1}}\ {\isasymrightarrow}\ {\isasymtau}\isactrlbsub {\isadigit{2}}\isactrlesub \ {\isasymrightarrow}\ {\isasymtau}} can be applied to \isa{x{\isacharcolon}\ {\isasymtau}\isactrlsub {\isadigit{1}}}
  and the remaining \isa{{\isacharparenleft}f\ x{\isacharparenright}{\isacharcolon}\ {\isasymtau}\isactrlsub {\isadigit{2}}\ {\isasymrightarrow}\ {\isasymtau}} passed to another function
  etc.  How well this works in practice depends on the order of
  arguments.  In the worst case, arguments are arranged erratically,
  and using a function in a certain situation always requires some
  glue code.  Thus we would get exponentially many oppurtunities to
  decorate the code with meaningless permutations of arguments.

  This can be avoided by \emph{canonical argument order}, which
  observes certain standard patterns and minimizes adhoc permutations
  in their application.  In Isabelle/ML, large portions text can be
  written without ever using \isa{swap{\isacharcolon}\ {\isasymalpha}\ {\isasymtimes}\ {\isasymbeta}\ {\isasymrightarrow}\ {\isasymbeta}\ {\isasymtimes}\ {\isasymalpha}}, or the
  combinator \isa{C{\isacharcolon}\ {\isacharparenleft}{\isasymalpha}\ {\isasymrightarrow}\ {\isasymbeta}\ {\isasymrightarrow}\ {\isasymgamma}{\isacharparenright}\ {\isasymrightarrow}\ {\isacharparenleft}{\isasymbeta}\ {\isasymrightarrow}\ {\isasymalpha}\ {\isasymrightarrow}\ {\isasymgamma}{\isacharparenright}} that is not even
  defined in our library.

  \medskip The basic idea is that arguments that vary less are moved
  further to the left than those that vary more.  Two particularly
  important categories of functions are \emph{selectors} and
  \emph{updates}.

  The subsequent scheme is based on a hypothetical set-like container
  of type \isa{{\isasymbeta}} that manages elements of type \isa{{\isasymalpha}}.  Both
  the names and types of the associated operations are canonical for
  Isabelle/ML.

  \medskip
  \begin{tabular}{ll}
  kind & canonical name and type \\\hline
  selector & \isa{member{\isacharcolon}\ {\isasymbeta}\ {\isasymrightarrow}\ {\isasymalpha}\ {\isasymrightarrow}\ bool} \\
  update & \isa{insert{\isacharcolon}\ {\isasymalpha}\ {\isasymrightarrow}\ {\isasymbeta}\ {\isasymrightarrow}\ {\isasymbeta}} \\
  \end{tabular}
  \medskip

  Given a container \isa{B{\isacharcolon}\ {\isasymbeta}}, the partially applied \isa{member\ B} is a predicate over elements \isa{{\isasymalpha}\ {\isasymrightarrow}\ bool}, and
  thus represents the intended denotation directly.  It is customary
  to pass the abstract predicate to further operations, not the
  concrete container.  The argument order makes it easy to use other
  combinators: \isa{forall\ {\isacharparenleft}member\ B{\isacharparenright}\ list} will check a list of
  elements for membership in \isa{B} etc. Often the explicit
  \isa{list} is pointless and can be contracted to \isa{forall\ {\isacharparenleft}member\ B{\isacharparenright}} to get directly a predicate again.

  In contrast, an update operation varies the container, so it moves
  to the right: \isa{insert\ a} is a function \isa{{\isasymbeta}\ {\isasymrightarrow}\ {\isasymbeta}} to
  insert a value \isa{a}.  These can be composed naturally as
  \isa{insert\ c\ {\isasymcirc}\ insert\ b\ {\isasymcirc}\ insert\ a}.  The slightly awkward
  inversion if the composition order is due to conventional
  mathematical notation, which can be easily amended as explained
  below.%
\end{isamarkuptext}%
\isamarkuptrue%
%
\isamarkupsubsection{Forward application and composition%
}
\isamarkuptrue%
%
\begin{isamarkuptext}%
Regular function application and infix notation works best for
  relatively deeply structured expressions, e.g.\ \isa{h\ {\isacharparenleft}f\ x\ y\ {\isacharplus}\ g\ z{\isacharparenright}}.  The important special case of \emph{linear transformation}
  applies a cascade of functions \isa{f\isactrlsub n\ {\isacharparenleft}{\isasymdots}\ {\isacharparenleft}f\isactrlsub {\isadigit{1}}\ x{\isacharparenright}{\isacharparenright}}.  This
  becomes hard to read and maintain if the functions are themselves
  given as complex expressions.  The notation can be significantly
  improved by introducing \emph{forward} versions of application and
  composition as follows:

  \medskip
  \begin{tabular}{lll}
  \isa{x\ {\isacharbar}{\isachargreater}\ f} & \isa{{\isasymequiv}} & \isa{f\ x} \\
  \isa{f\ {\isacharhash}{\isachargreater}\ g} & \isa{{\isasymequiv}} & \isa{x\ {\isacharbar}{\isachargreater}\ f\ {\isacharbar}{\isachargreater}\ g} \\
  \end{tabular}
  \medskip

  This enables to write conveniently \isa{x\ {\isacharbar}{\isachargreater}\ f\isactrlsub {\isadigit{1}}\ {\isacharbar}{\isachargreater}\ {\isasymdots}\ {\isacharbar}{\isachargreater}\ f\isactrlsub n} or
  \isa{f\isactrlsub {\isadigit{1}}\ {\isacharhash}{\isachargreater}\ {\isasymdots}\ {\isacharhash}{\isachargreater}\ f\isactrlsub n} for its functional abstraction over \isa{x}.

  \medskip There is an additional set of combinators to accommodate
  multiple results (via pairs) that are passed on as multiple
  arguments (via currying).

  \medskip
  \begin{tabular}{lll}
  \isa{{\isacharparenleft}x{\isacharcomma}\ y{\isacharparenright}\ {\isacharbar}{\isacharminus}{\isachargreater}\ f} & \isa{{\isasymequiv}} & \isa{f\ x\ y} \\
  \isa{f\ {\isacharhash}{\isacharminus}{\isachargreater}\ g} & \isa{{\isasymequiv}} & \isa{x\ {\isacharbar}{\isachargreater}\ f\ {\isacharbar}{\isacharminus}{\isachargreater}\ g} \\
  \end{tabular}
  \medskip%
\end{isamarkuptext}%
\isamarkuptrue%
%
\isadelimmlref
%
\endisadelimmlref
%
\isatagmlref
%
\begin{isamarkuptext}%
\begin{mldecls}
  \indexdef{}{ML}{$\mid$$>$}\verb|op |\verb,|,\verb|>  : 'a * ('a -> 'b) -> 'b| \\
  \indexdef{}{ML}{$\mid$-$>$}\verb|op |\verb,|,\verb|->  : ('c * 'a) * ('c -> 'a -> 'b) -> 'b| \\
  \indexdef{}{ML}{\#$>$}\verb|op #>  : ('a -> 'b) * ('b -> 'c) -> 'a -> 'c| \\
  \indexdef{}{ML}{\#-$>$}\verb|op #->  : ('a -> 'c * 'b) * ('c -> 'b -> 'd) -> 'a -> 'd| \\
  \end{mldecls}

  %FIXME description!?%
\end{isamarkuptext}%
\isamarkuptrue%
%
\endisatagmlref
{\isafoldmlref}%
%
\isadelimmlref
%
\endisadelimmlref
%
\isamarkupsubsection{Canonical iteration%
}
\isamarkuptrue%
%
\begin{isamarkuptext}%
As explained above, a function \isa{f{\isacharcolon}\ {\isasymalpha}\ {\isasymrightarrow}\ {\isasymbeta}\ {\isasymrightarrow}\ {\isasymbeta}} can be
  understood as update on a configuration of type \isa{{\isasymbeta}},
  parametrized by arguments of type \isa{{\isasymalpha}}.  Given \isa{a{\isacharcolon}\ {\isasymalpha}}
  the partial application \isa{{\isacharparenleft}f\ a{\isacharparenright}{\isacharcolon}\ {\isasymbeta}\ {\isasymrightarrow}\ {\isasymbeta}} operates
  homogeneously on \isa{{\isasymbeta}}.  This can be iterated naturally over a
  list of parameters \isa{{\isacharbrackleft}a\isactrlsub {\isadigit{1}}{\isacharcomma}\ {\isasymdots}{\isacharcomma}\ a\isactrlsub n{\isacharbrackright}} as \isa{f\ a\isactrlsub {\isadigit{1}}\ {\isacharhash}{\isachargreater}\ {\isasymdots}\ {\isacharhash}{\isachargreater}\ f\ a\isactrlbsub n\isactrlesub \isactrlbsub \isactrlesub }.  The latter expression is again a function \isa{{\isasymbeta}\ {\isasymrightarrow}\ {\isasymbeta}}.
  It can be applied to an initial configuration \isa{b{\isacharcolon}\ {\isasymbeta}} to
  start the iteration over the given list of arguments: each \isa{a} in \isa{a\isactrlsub {\isadigit{1}}{\isacharcomma}\ {\isasymdots}{\isacharcomma}\ a\isactrlsub n} is applied consecutively by updating a
  cumulative configuration.

  The \isa{fold} combinator in Isabelle/ML lifts a function \isa{f} as above to its iterated version over a list of arguments.
  Lifting can be repeated, e.g.\ \isa{{\isacharparenleft}fold\ {\isasymcirc}\ fold{\isacharparenright}\ f} iterates
  over a list of lists as expected.

  The variant \isa{fold{\isacharunderscore}rev} works inside-out over the list of
  arguments, such that \isa{fold{\isacharunderscore}rev\ f\ {\isasymequiv}\ fold\ f\ {\isasymcirc}\ rev} holds.

  The \isa{fold{\isacharunderscore}map} combinator essentially performs \isa{fold} and \isa{map} simultaneously: each application of \isa{f} produces an updated configuration together with a side-result;
  the iteration collects all such side-results as a separate list.%
\end{isamarkuptext}%
\isamarkuptrue%
%
\isadelimmlref
%
\endisadelimmlref
%
\isatagmlref
%
\begin{isamarkuptext}%
\begin{mldecls}
  \indexdef{}{ML}{fold}\verb|fold: ('a -> 'b -> 'b) -> 'a list -> 'b -> 'b| \\
  \indexdef{}{ML}{fold\_rev}\verb|fold_rev: ('a -> 'b -> 'b) -> 'a list -> 'b -> 'b| \\
  \indexdef{}{ML}{fold\_map}\verb|fold_map: ('a -> 'b -> 'c * 'b) -> 'a list -> 'b -> 'c list * 'b| \\
  \end{mldecls}

  \begin{description}

  \item \verb|fold|~\isa{f} lifts the parametrized update function
  \isa{f} to a list of parameters.

  \item \verb|fold_rev|~\isa{f} is similar to \verb|fold|~\isa{f}, but works inside-out.

  \item \verb|fold_map|~\isa{f} lifts the parametrized update
  function \isa{f} (with side-result) to a list of parameters and
  cumulative side-results.

  \end{description}

  \begin{warn}
  The literature on functional programming provides a multitude of
  combinators called \isa{foldl}, \isa{foldr} etc.  SML97
  provides its own variations as \verb|List.foldl| and \verb|List.foldr|, while the classic Isabelle library also has the
  historic \verb|Library.foldl| and \verb|Library.foldr|.  To avoid
  further confusion, all of this should be ignored, and \verb|fold| (or
  \verb|fold_rev|) used exclusively.
  \end{warn}%
\end{isamarkuptext}%
\isamarkuptrue%
%
\endisatagmlref
{\isafoldmlref}%
%
\isadelimmlref
%
\endisadelimmlref
%
\isadelimmlex
%
\endisadelimmlex
%
\isatagmlex
%
\begin{isamarkuptext}%
The following example shows how to fill a text buffer
  incrementally by adding strings, either individually or from a given
  list.%
\end{isamarkuptext}%
\isamarkuptrue%
%
\endisatagmlex
{\isafoldmlex}%
%
\isadelimmlex
%
\endisadelimmlex
%
\isadelimML
%
\endisadelimML
%
\isatagML
\isacommand{ML}\isamarkupfalse%
\ {\isacharverbatimopen}\isanewline
\ \ val\ s\ {\isacharequal}\isanewline
\ \ \ \ Buffer{\isachardot}empty\isanewline
\ \ \ \ {\isacharbar}{\isachargreater}\ Buffer{\isachardot}add\ {\isachardoublequote}digits{\isacharcolon}\ {\isachardoublequote}\isanewline
\ \ \ \ {\isacharbar}{\isachargreater}\ fold\ {\isacharparenleft}Buffer{\isachardot}add\ o\ string{\isacharunderscore}of{\isacharunderscore}int{\isacharparenright}\ {\isacharparenleft}{\isadigit{0}}\ upto\ {\isadigit{9}}{\isacharparenright}\isanewline
\ \ \ \ {\isacharbar}{\isachargreater}\ Buffer{\isachardot}content{\isacharsemicolon}\isanewline
\isanewline
\ \ %
\isaantiq
assert%
\endisaantiq
\ {\isacharparenleft}s\ {\isacharequal}\ {\isachardoublequote}digits{\isacharcolon}\ {\isadigit{0}}{\isadigit{1}}{\isadigit{2}}{\isadigit{3}}{\isadigit{4}}{\isadigit{5}}{\isadigit{6}}{\isadigit{7}}{\isadigit{8}}{\isadigit{9}}{\isachardoublequote}{\isacharparenright}{\isacharsemicolon}\isanewline
{\isacharverbatimclose}%
\endisatagML
{\isafoldML}%
%
\isadelimML
%
\endisadelimML
%
\begin{isamarkuptext}%
Note how \verb|fold (Buffer.add o string_of_int)| above saves
  an extra \verb|map| over the given list.  This kind of peephole
  optimization reduces both the code size and the tree structures in
  memory (``deforestation''), but requires some practice to read and
  write it fluently.

  \medskip The next example elaborates the idea of canonical
  iteration, demonstrating fast accumulation of tree content using a
  text buffer.%
\end{isamarkuptext}%
\isamarkuptrue%
%
\isadelimML
%
\endisadelimML
%
\isatagML
\isacommand{ML}\isamarkupfalse%
\ {\isacharverbatimopen}\isanewline
\ \ datatype\ tree\ {\isacharequal}\ Text\ of\ string\ {\isacharbar}\ Elem\ of\ string\ {\isacharasterisk}\ tree\ list{\isacharsemicolon}\isanewline
\isanewline
\ \ fun\ slow{\isacharunderscore}content\ {\isacharparenleft}Text\ txt{\isacharparenright}\ {\isacharequal}\ txt\isanewline
\ \ \ \ {\isacharbar}\ slow{\isacharunderscore}content\ {\isacharparenleft}Elem\ {\isacharparenleft}name{\isacharcomma}\ ts{\isacharparenright}{\isacharparenright}\ {\isacharequal}\isanewline
\ \ \ \ \ \ \ \ {\isachardoublequote}{\isacharless}{\isachardoublequote}\ {\isacharcircum}\ name\ {\isacharcircum}\ {\isachardoublequote}{\isachargreater}{\isachardoublequote}\ {\isacharcircum}\isanewline
\ \ \ \ \ \ \ \ implode\ {\isacharparenleft}map\ slow{\isacharunderscore}content\ ts{\isacharparenright}\ {\isacharcircum}\isanewline
\ \ \ \ \ \ \ \ {\isachardoublequote}{\isacharless}{\isacharslash}{\isachardoublequote}\ {\isacharcircum}\ name\ {\isacharcircum}\ {\isachardoublequote}{\isachargreater}{\isachardoublequote}\isanewline
\isanewline
\ \ fun\ add{\isacharunderscore}content\ {\isacharparenleft}Text\ txt{\isacharparenright}\ {\isacharequal}\ Buffer{\isachardot}add\ txt\isanewline
\ \ \ \ {\isacharbar}\ add{\isacharunderscore}content\ {\isacharparenleft}Elem\ {\isacharparenleft}name{\isacharcomma}\ ts{\isacharparenright}{\isacharparenright}\ {\isacharequal}\isanewline
\ \ \ \ \ \ \ \ Buffer{\isachardot}add\ {\isacharparenleft}{\isachardoublequote}{\isacharless}{\isachardoublequote}\ {\isacharcircum}\ name\ {\isacharcircum}\ {\isachardoublequote}{\isachargreater}{\isachardoublequote}{\isacharparenright}\ {\isacharhash}{\isachargreater}\isanewline
\ \ \ \ \ \ \ \ fold\ add{\isacharunderscore}content\ ts\ {\isacharhash}{\isachargreater}\isanewline
\ \ \ \ \ \ \ \ Buffer{\isachardot}add\ {\isacharparenleft}{\isachardoublequote}{\isacharless}{\isacharslash}{\isachardoublequote}\ {\isacharcircum}\ name\ {\isacharcircum}\ {\isachardoublequote}{\isachargreater}{\isachardoublequote}{\isacharparenright}{\isacharsemicolon}\isanewline
\isanewline
\ \ fun\ fast{\isacharunderscore}content\ tree\ {\isacharequal}\isanewline
\ \ \ \ Buffer{\isachardot}empty\ {\isacharbar}{\isachargreater}\ add{\isacharunderscore}content\ tree\ {\isacharbar}{\isachargreater}\ Buffer{\isachardot}content{\isacharsemicolon}\isanewline
{\isacharverbatimclose}%
\endisatagML
{\isafoldML}%
%
\isadelimML
%
\endisadelimML
%
\begin{isamarkuptext}%
The slow part of \verb|slow_content| is the \verb|implode| of
  the recursive results, because it copies previously produced strings
  again.

  The incremental \verb|add_content| avoids this by operating on a
  buffer that is passed through in a linear fashion.  Using \verb|#>| and contraction over the actual buffer argument saves some
  additional boiler-plate.  Of course, the two \verb|Buffer.add|
  invocations with concatenated strings could have been split into
  smaller parts, but this would have obfuscated the source without
  making a big difference in allocations.  Here we have done some
  peephole-optimization for the sake of readability.

  Another benefit of \verb|add_content| is its ``open'' form as a
  function on buffers that can be continued in further linear
  transformations, folding etc.  Thus it is more compositional than
  the naive \verb|slow_content|.  As realistic example, compare the
  old-style \verb|Term.maxidx_of_term: term -> int| with the newer
  \verb|Term.maxidx_term: term -> int -> int| in Isabelle/Pure.

  Note that \verb|fast_content| above is only defined as example.  In
  many practical situations, it is customary to provide the
  incremental \verb|add_content| only and leave the initialization and
  termination to the concrete application by the user.%
\end{isamarkuptext}%
\isamarkuptrue%
%
\isamarkupsection{Message output channels \label{sec:message-channels}%
}
\isamarkuptrue%
%
\begin{isamarkuptext}%
Isabelle provides output channels for different kinds of
  messages: regular output, high-volume tracing information, warnings,
  and errors.

  Depending on the user interface involved, these messages may appear
  in different text styles or colours.  The standard output for
  terminal sessions prefixes each line of warnings by \verb|###| and errors by \verb|***|, but leaves anything else
  unchanged.

  Messages are associated with the transaction context of the running
  Isar command.  This enables the front-end to manage commands and
  resulting messages together.  For example, after deleting a command
  from a given theory document version, the corresponding message
  output can be retracted from the display.%
\end{isamarkuptext}%
\isamarkuptrue%
%
\isadelimmlref
%
\endisadelimmlref
%
\isatagmlref
%
\begin{isamarkuptext}%
\begin{mldecls}
  \indexdef{}{ML}{writeln}\verb|writeln: string -> unit| \\
  \indexdef{}{ML}{tracing}\verb|tracing: string -> unit| \\
  \indexdef{}{ML}{warning}\verb|warning: string -> unit| \\
  \indexdef{}{ML}{error}\verb|error: string -> 'a| \\
  \end{mldecls}

  \begin{description}

  \item \verb|writeln|~\isa{text} outputs \isa{text} as regular
  message.  This is the primary message output operation of Isabelle
  and should be used by default.

  \item \verb|tracing|~\isa{text} outputs \isa{text} as special
  tracing message, indicating potential high-volume output to the
  front-end (hundreds or thousands of messages issued by a single
  command).  The idea is to allow the user-interface to downgrade the
  quality of message display to achieve higher throughput.

  Note that the user might have to take special actions to see tracing
  output, e.g.\ switch to a different output window.  So this channel
  should not be used for regular output.

  \item \verb|warning|~\isa{text} outputs \isa{text} as
  warning, which typically means some extra emphasis on the front-end
  side (color highlighting, icons, etc.).

  \item \verb|error|~\isa{text} raises exception \verb|ERROR|~\isa{text} and thus lets the Isar toplevel print \isa{text} on the
  error channel, which typically means some extra emphasis on the
  front-end side (color highlighting, icons, etc.).

  This assumes that the exception is not handled before the command
  terminates.  Handling exception \verb|ERROR|~\isa{text} is a
  perfectly legal alternative: it means that the error is absorbed
  without any message output.

  \begin{warn}
  The actual error channel is accessed via \verb|Output.error_msg|, but
  the interaction protocol of Proof~General \emph{crashes} if that
  function is used in regular ML code: error output and toplevel
  command failure always need to coincide.
  \end{warn}

  \end{description}

  \begin{warn}
  Regular Isabelle/ML code should output messages exclusively by the
  official channels.  Using raw I/O on \emph{stdout} or \emph{stderr}
  instead (e.g.\ via \verb|TextIO.output|) is apt to cause problems in
  the presence of parallel and asynchronous processing of Isabelle
  theories.  Such raw output might be displayed by the front-end in
  some system console log, with a low chance that the user will ever
  see it.  Moreover, as a genuine side-effect on global process
  channels, there is no proper way to retract output when Isar command
  transactions are reset by the system.
  \end{warn}

  \begin{warn}
  The message channels should be used in a message-oriented manner.
  This means that multi-line output that logically belongs together is
  issued by a \emph{single} invocation of \verb|writeln| etc.\ with the
  functional concatenation of all message constituents.
  \end{warn}%
\end{isamarkuptext}%
\isamarkuptrue%
%
\endisatagmlref
{\isafoldmlref}%
%
\isadelimmlref
%
\endisadelimmlref
%
\isadelimmlex
%
\endisadelimmlex
%
\isatagmlex
%
\begin{isamarkuptext}%
The following example demonstrates a multi-line
  warning.  Note that in some situations the user sees only the first
  line, so the most important point should be made first.%
\end{isamarkuptext}%
\isamarkuptrue%
%
\endisatagmlex
{\isafoldmlex}%
%
\isadelimmlex
%
\endisadelimmlex
%
\isadelimML
%
\endisadelimML
%
\isatagML
\isacommand{ML{\isacharunderscore}command}\isamarkupfalse%
\ {\isacharverbatimopen}\isanewline
\ \ warning\ {\isacharparenleft}cat{\isacharunderscore}lines\isanewline
\ \ \ {\isacharbrackleft}{\isachardoublequote}Beware\ the\ Jabberwock{\isacharcomma}\ my\ son{\isacharbang}{\isachardoublequote}{\isacharcomma}\isanewline
\ \ \ \ {\isachardoublequote}The\ jaws\ that\ bite{\isacharcomma}\ the\ claws\ that\ catch{\isacharbang}{\isachardoublequote}{\isacharcomma}\isanewline
\ \ \ \ {\isachardoublequote}Beware\ the\ Jubjub\ Bird{\isacharcomma}\ and\ shun{\isachardoublequote}{\isacharcomma}\isanewline
\ \ \ \ {\isachardoublequote}The\ frumious\ Bandersnatch{\isacharbang}{\isachardoublequote}{\isacharbrackright}{\isacharparenright}{\isacharsemicolon}\isanewline
{\isacharverbatimclose}%
\endisatagML
{\isafoldML}%
%
\isadelimML
%
\endisadelimML
%
\isamarkupsection{Exceptions \label{sec:exceptions}%
}
\isamarkuptrue%
%
\begin{isamarkuptext}%
The Standard ML semantics of strict functional evaluation
  together with exceptions is rather well defined, but some delicate
  points need to be observed to avoid that ML programs go wrong
  despite static type-checking.  Exceptions in Isabelle/ML are
  subsequently categorized as follows.

  \paragraph{Regular user errors.}  These are meant to provide
  informative feedback about malformed input etc.

  The \emph{error} function raises the corresponding \emph{ERROR}
  exception, with a plain text message as argument.  \emph{ERROR}
  exceptions can be handled internally, in order to be ignored, turned
  into other exceptions, or cascaded by appending messages.  If the
  corresponding Isabelle/Isar command terminates with an \emph{ERROR}
  exception state, the toplevel will print the result on the error
  channel (see \secref{sec:message-channels}).

  It is considered bad style to refer to internal function names or
  values in ML source notation in user error messages.

  Grammatical correctness of error messages can be improved by
  \emph{omitting} final punctuation: messages are often concatenated
  or put into a larger context (e.g.\ augmented with source position).
  By not insisting in the final word at the origin of the error, the
  system can perform its administrative tasks more easily and
  robustly.

  \paragraph{Program failures.}  There is a handful of standard
  exceptions that indicate general failure situations, or failures of
  core operations on logical entities (types, terms, theorems,
  theories, see \chref{ch:logic}).

  These exceptions indicate a genuine breakdown of the program, so the
  main purpose is to determine quickly what has happened where.
  Traditionally, the (short) exception message would include the name
  of an ML function, although this is no longer necessary, because the
  ML runtime system prints a detailed source position of the
  corresponding \verb|raise| keyword.

  \medskip User modules can always introduce their own custom
  exceptions locally, e.g.\ to organize internal failures robustly
  without overlapping with existing exceptions.  Exceptions that are
  exposed in module signatures require extra care, though, and should
  \emph{not} be introduced by default.  Surprise by users of a module
  can be often minimized by using plain user errors instead.

  \paragraph{Interrupts.}  These indicate arbitrary system events:
  both the ML runtime system and the Isabelle/ML infrastructure signal
  various exceptional situations by raising the special
  \emph{Interrupt} exception in user code.

  This is the one and only way that physical events can intrude an
  Isabelle/ML program.  Such an interrupt can mean out-of-memory,
  stack overflow, timeout, internal signaling of threads, or the user
  producing a console interrupt manually etc.  An Isabelle/ML program
  that intercepts interrupts becomes dependent on physical effects of
  the environment.  Even worse, exception handling patterns that are
  too general by accident, e.g.\ by mispelled exception constructors,
  will cover interrupts unintentionally and thus render the program
  semantics ill-defined.

  Note that the Interrupt exception dates back to the original SML90
  language definition.  It was excluded from the SML97 version to
  avoid its malign impact on ML program semantics, but without
  providing a viable alternative.  Isabelle/ML recovers physical
  interruptibility (which an indispensable tool to implement managed
  evaluation of command transactions), but requires user code to be
  strictly transparent wrt.\ interrupts.

  \begin{warn}
  Isabelle/ML user code needs to terminate promptly on interruption,
  without guessing at its meaning to the system infrastructure.
  Temporary handling of interrupts for cleanup of global resources
  etc.\ needs to be followed immediately by re-raising of the original
  exception.
  \end{warn}%
\end{isamarkuptext}%
\isamarkuptrue%
%
\isadelimmlref
%
\endisadelimmlref
%
\isatagmlref
%
\begin{isamarkuptext}%
\begin{mldecls}
  \indexdef{}{ML}{try}\verb|try: ('a -> 'b) -> 'a -> 'b option| \\
  \indexdef{}{ML}{can}\verb|can: ('a -> 'b) -> 'a -> bool| \\
  \indexdef{}{ML}{ERROR}\verb|ERROR: string -> exn| \\
  \indexdef{}{ML}{Fail}\verb|Fail: string -> exn| \\
  \indexdef{}{ML}{Exn.is\_interrupt}\verb|Exn.is_interrupt: exn -> bool| \\
  \indexdef{}{ML}{reraise}\verb|reraise: exn -> 'a| \\
  \indexdef{}{ML}{exception\_trace}\verb|exception_trace: (unit -> 'a) -> 'a| \\
  \end{mldecls}

  \begin{description}

  \item \verb|try|~\isa{f\ x} makes the partiality of evaluating
  \isa{f\ x} explicit via the option datatype.  Interrupts are
  \emph{not} handled here, i.e.\ this form serves as safe replacement
  for the \emph{unsafe} version \verb|(SOME|~\isa{f\ x}~\verb|handle _ => NONE)| that is occasionally seen in
  books about SML.

  \item \verb|can| is similar to \verb|try| with more abstract result.

  \item \verb|ERROR|~\isa{msg} represents user errors; this
  exception is normally raised indirectly via the \verb|error| function
  (see \secref{sec:message-channels}).

  \item \verb|Fail|~\isa{msg} represents general program failures.

  \item \verb|Exn.is_interrupt| identifies interrupts robustly, without
  mentioning concrete exception constructors in user code.  Handled
  interrupts need to be re-raised promptly!

  \item \verb|reraise|~\isa{exn} raises exception \isa{exn}
  while preserving its implicit position information (if possible,
  depending on the ML platform).

  \item \verb|exception_trace|~\verb|(fn () =>|~\isa{e}\verb|)| evaluates expression \isa{e} while printing
  a full trace of its stack of nested exceptions (if possible,
  depending on the ML platform).\footnote{In versions of Poly/ML the
  trace will appear on raw stdout of the Isabelle process.}

  Inserting \verb|exception_trace| into ML code temporarily is useful
  for debugging, but not suitable for production code.

  \end{description}%
\end{isamarkuptext}%
\isamarkuptrue%
%
\endisatagmlref
{\isafoldmlref}%
%
\isadelimmlref
%
\endisadelimmlref
%
\isadelimmlantiq
%
\endisadelimmlantiq
%
\isatagmlantiq
%
\begin{isamarkuptext}%
\begin{matharray}{rcl}
  \indexdef{}{ML antiquotation}{assert}\hypertarget{ML antiquotation.assert}{\hyperlink{ML antiquotation.assert}{\mbox{\isa{assert}}}} & : & \isa{ML{\isacharunderscore}antiquotation} \\
  \end{matharray}

  \begin{description}

  \item \isa{{\isacharat}{\isacharbraceleft}assert{\isacharbraceright}} inlines a function
  \verb|bool -> unit| that raises \verb|Fail| if the argument is
  \verb|false|.  Due to inlining the source position of failed
  assertions is included in the error output.

  \end{description}%
\end{isamarkuptext}%
\isamarkuptrue%
%
\endisatagmlantiq
{\isafoldmlantiq}%
%
\isadelimmlantiq
%
\endisadelimmlantiq
%
\isamarkupsection{Basic data types%
}
\isamarkuptrue%
%
\begin{isamarkuptext}%
The basis library proposal of SML97 needs to be treated with
  caution.  Many of its operations simply do not fit with important
  Isabelle/ML conventions (like ``canonical argument order'', see
  \secref{sec:canonical-argument-order}), others cause problems with
  the parallel evaluation model of Isabelle/ML (such as \verb|TextIO.print| or \verb|OS.Process.system|).

  Subsequently we give a brief overview of important operations on
  basic ML data types.%
\end{isamarkuptext}%
\isamarkuptrue%
%
\isamarkupsubsection{Characters%
}
\isamarkuptrue%
%
\isadelimmlref
%
\endisadelimmlref
%
\isatagmlref
%
\begin{isamarkuptext}%
\begin{mldecls}
  \indexdef{}{ML type}{char}\verb|type char| \\
  \end{mldecls}

  \begin{description}

  \item Type \verb|char| is \emph{not} used.  The smallest textual
  unit in Isabelle is represented as a ``symbol'' (see
  \secref{sec:symbols}).

  \end{description}%
\end{isamarkuptext}%
\isamarkuptrue%
%
\endisatagmlref
{\isafoldmlref}%
%
\isadelimmlref
%
\endisadelimmlref
%
\isamarkupsubsection{Integers%
}
\isamarkuptrue%
%
\isadelimmlref
%
\endisadelimmlref
%
\isatagmlref
%
\begin{isamarkuptext}%
\begin{mldecls}
  \indexdef{}{ML type}{int}\verb|type int| \\
  \end{mldecls}

  \begin{description}

  \item Type \verb|int| represents regular mathematical integers,
  which are \emph{unbounded}.  Overflow never happens in
  practice.\footnote{The size limit for integer bit patterns in memory
  is 64\,MB for 32-bit Poly/ML, and much higher for 64-bit systems.}
  This works uniformly for all supported ML platforms (Poly/ML and
  SML/NJ).

  Literal integers in ML text are forced to be of this one true
  integer type --- overloading of SML97 is disabled.

  Structure \verb|IntInf| of SML97 is obsolete and superseded by
  \verb|Int|.  Structure \verb|Integer| in \hyperlink{file.~~/src/Pure/General/integer.ML}{\mbox{\isa{\isatt{{\isachartilde}{\isachartilde}{\isacharslash}src{\isacharslash}Pure{\isacharslash}General{\isacharslash}integer{\isachardot}ML}}}} provides some additional
  operations.

  \end{description}%
\end{isamarkuptext}%
\isamarkuptrue%
%
\endisatagmlref
{\isafoldmlref}%
%
\isadelimmlref
%
\endisadelimmlref
%
\isamarkupsubsection{Options%
}
\isamarkuptrue%
%
\isadelimmlref
%
\endisadelimmlref
%
\isatagmlref
%
\begin{isamarkuptext}%
\begin{mldecls}
  \indexdef{}{ML}{Option.map}\verb|Option.map: ('a -> 'b) -> 'a option -> 'b option| \\
  \indexdef{}{ML}{is\_some}\verb|is_some: 'a option -> bool| \\
  \indexdef{}{ML}{is\_none}\verb|is_none: 'a option -> bool| \\
  \indexdef{}{ML}{the}\verb|the: 'a option -> 'a| \\
  \indexdef{}{ML}{these}\verb|these: 'a list option -> 'a list| \\
  \indexdef{}{ML}{the\_list}\verb|the_list: 'a option -> 'a list| \\
  \indexdef{}{ML}{the\_default}\verb|the_default: 'a -> 'a option -> 'a| \\
  \end{mldecls}%
\end{isamarkuptext}%
\isamarkuptrue%
%
\endisatagmlref
{\isafoldmlref}%
%
\isadelimmlref
%
\endisadelimmlref
%
\begin{isamarkuptext}%
Apart from \verb|Option.map| most operations defined in
  structure \verb|Option| are alien to Isabelle/ML.  The
  operations shown above are defined in \hyperlink{file.~~/src/Pure/General/basics.ML}{\mbox{\isa{\isatt{{\isachartilde}{\isachartilde}{\isacharslash}src{\isacharslash}Pure{\isacharslash}General{\isacharslash}basics{\isachardot}ML}}}}, among others.%
\end{isamarkuptext}%
\isamarkuptrue%
%
\isamarkupsubsection{Lists%
}
\isamarkuptrue%
%
\begin{isamarkuptext}%
Lists are ubiquitous in ML as simple and light-weight
  ``collections'' for many everyday programming tasks.  Isabelle/ML
  provides important additions and improvements over operations that
  are predefined in the SML97 library.%
\end{isamarkuptext}%
\isamarkuptrue%
%
\isadelimmlref
%
\endisadelimmlref
%
\isatagmlref
%
\begin{isamarkuptext}%
\begin{mldecls}
  \indexdef{}{ML}{cons}\verb|cons: 'a -> 'a list -> 'a list| \\
  \indexdef{}{ML}{member}\verb|member: ('b * 'a -> bool) -> 'a list -> 'b -> bool| \\
  \indexdef{}{ML}{insert}\verb|insert: ('a * 'a -> bool) -> 'a -> 'a list -> 'a list| \\
  \indexdef{}{ML}{remove}\verb|remove: ('b * 'a -> bool) -> 'b -> 'a list -> 'a list| \\
  \indexdef{}{ML}{update}\verb|update: ('a * 'a -> bool) -> 'a -> 'a list -> 'a list| \\
  \end{mldecls}

  \begin{description}

  \item \verb|cons|~\isa{x\ xs} evaluates to \isa{x\ {\isacharcolon}{\isacharcolon}\ xs}.

  Tupled infix operators are a historical accident in Standard ML.
  The curried \verb|cons| amends this, but it should be only used when
  partial application is required.

  \item \verb|member|, \verb|insert|, \verb|remove|, \verb|update| treat
  lists as a set-like container that maintains the order of elements.
  See \hyperlink{file.~~/src/Pure/library.ML}{\mbox{\isa{\isatt{{\isachartilde}{\isachartilde}{\isacharslash}src{\isacharslash}Pure{\isacharslash}library{\isachardot}ML}}}} for the full specifications
  (written in ML).  There are some further derived operations like
  \verb|union| or \verb|inter|.

  Note that \verb|insert| is conservative about elements that are
  already a \verb|member| of the list, while \verb|update| ensures that
  the latest entry is always put in front.  The latter discipline is
  often more appropriate in declarations of context data
  (\secref{sec:context-data}) that are issued by the user in Isar
  source: more recent declarations normally take precedence over
  earlier ones.

  \end{description}%
\end{isamarkuptext}%
\isamarkuptrue%
%
\endisatagmlref
{\isafoldmlref}%
%
\isadelimmlref
%
\endisadelimmlref
%
\isadelimmlex
%
\endisadelimmlex
%
\isatagmlex
%
\begin{isamarkuptext}%
Using canonical \verb|fold| together with \verb|cons|, or
  similar standard operations, alternates the orientation of data.
  The is quite natural and should not be altered forcible by inserting
  extra applications of \verb|rev|.  The alternative \verb|fold_rev| can
  be used in the few situations, where alternation should be
  prevented.%
\end{isamarkuptext}%
\isamarkuptrue%
%
\endisatagmlex
{\isafoldmlex}%
%
\isadelimmlex
%
\endisadelimmlex
%
\isadelimML
%
\endisadelimML
%
\isatagML
\isacommand{ML}\isamarkupfalse%
\ {\isacharverbatimopen}\isanewline
\ \ val\ items\ {\isacharequal}\ {\isacharbrackleft}{\isadigit{1}}{\isacharcomma}\ {\isadigit{2}}{\isacharcomma}\ {\isadigit{3}}{\isacharcomma}\ {\isadigit{4}}{\isacharcomma}\ {\isadigit{5}}{\isacharcomma}\ {\isadigit{6}}{\isacharcomma}\ {\isadigit{7}}{\isacharcomma}\ {\isadigit{8}}{\isacharcomma}\ {\isadigit{9}}{\isacharcomma}\ {\isadigit{1}}{\isadigit{0}}{\isacharbrackright}{\isacharsemicolon}\isanewline
\isanewline
\ \ val\ list{\isadigit{1}}\ {\isacharequal}\ fold\ cons\ items\ {\isacharbrackleft}{\isacharbrackright}{\isacharsemicolon}\isanewline
\ \ %
\isaantiq
assert%
\endisaantiq
\ {\isacharparenleft}list{\isadigit{1}}\ {\isacharequal}\ rev\ items{\isacharparenright}{\isacharsemicolon}\isanewline
\isanewline
\ \ val\ list{\isadigit{2}}\ {\isacharequal}\ fold{\isacharunderscore}rev\ cons\ items\ {\isacharbrackleft}{\isacharbrackright}{\isacharsemicolon}\isanewline
\ \ %
\isaantiq
assert%
\endisaantiq
\ {\isacharparenleft}list{\isadigit{2}}\ {\isacharequal}\ items{\isacharparenright}{\isacharsemicolon}\isanewline
{\isacharverbatimclose}%
\endisatagML
{\isafoldML}%
%
\isadelimML
%
\endisadelimML
%
\begin{isamarkuptext}%
The subsequent example demonstrates how to \emph{merge} two
  lists in a natural way.%
\end{isamarkuptext}%
\isamarkuptrue%
%
\isadelimML
%
\endisadelimML
%
\isatagML
\isacommand{ML}\isamarkupfalse%
\ {\isacharverbatimopen}\isanewline
\ \ fun\ merge{\isacharunderscore}lists\ eq\ {\isacharparenleft}xs{\isacharcomma}\ ys{\isacharparenright}\ {\isacharequal}\ fold{\isacharunderscore}rev\ {\isacharparenleft}insert\ eq{\isacharparenright}\ ys\ xs{\isacharsemicolon}\isanewline
{\isacharverbatimclose}%
\endisatagML
{\isafoldML}%
%
\isadelimML
%
\endisadelimML
%
\begin{isamarkuptext}%
Here the first list is treated conservatively: only the new
  elements from the second list are inserted.  The inside-out order of
  insertion via \verb|fold_rev| attempts to preserve the order of
  elements in the result.

  This way of merging lists is typical for context data
  (\secref{sec:context-data}).  See also \verb|merge| as defined in
  \hyperlink{file.~~/src/Pure/library.ML}{\mbox{\isa{\isatt{{\isachartilde}{\isachartilde}{\isacharslash}src{\isacharslash}Pure{\isacharslash}library{\isachardot}ML}}}}.%
\end{isamarkuptext}%
\isamarkuptrue%
%
\isamarkupsubsection{Association lists%
}
\isamarkuptrue%
%
\begin{isamarkuptext}%
The operations for association lists interpret a concrete list
  of pairs as a finite function from keys to values.  Redundant
  representations with multiple occurrences of the same key are
  implicitly normalized: lookup and update only take the first
  occurrence into account.%
\end{isamarkuptext}%
\isamarkuptrue%
%
\begin{isamarkuptext}%
\begin{mldecls}
  \indexdef{}{ML}{AList.lookup}\verb|AList.lookup: ('a * 'b -> bool) -> ('b * 'c) list -> 'a -> 'c option| \\
  \indexdef{}{ML}{AList.defined}\verb|AList.defined: ('a * 'b -> bool) -> ('b * 'c) list -> 'a -> bool| \\
  \indexdef{}{ML}{AList.update}\verb|AList.update: ('a * 'a -> bool) -> 'a * 'b -> ('a * 'b) list -> ('a * 'b) list| \\
  \end{mldecls}

  \begin{description}

  \item \verb|AList.lookup|, \verb|AList.defined|, \verb|AList.update|
  implement the main ``framework operations'' for mappings in
  Isabelle/ML, following standard conventions for their names and
  types.

  Note that a function called \isa{lookup} is obliged to express its
  partiality via an explicit option element.  There is no choice to
  raise an exception, without changing the name to something like
  \isa{the{\isacharunderscore}element} or \isa{get}.

  The \isa{defined} operation is essentially a contraction of \verb|is_some| and \isa{lookup}, but this is sufficiently frequent to
  justify its independent existence.  This also gives the
  implementation some opportunity for peep-hole optimization.

  \end{description}

  Association lists are adequate as simple and light-weight
  implementation of finite mappings in many practical situations.  A
  more heavy-duty table structure is defined in \hyperlink{file.~~/src/Pure/General/table.ML}{\mbox{\isa{\isatt{{\isachartilde}{\isachartilde}{\isacharslash}src{\isacharslash}Pure{\isacharslash}General{\isacharslash}table{\isachardot}ML}}}}; that version scales easily to
  thousands or millions of elements.%
\end{isamarkuptext}%
\isamarkuptrue%
%
\isamarkupsubsection{Unsynchronized references%
}
\isamarkuptrue%
%
\isadelimmlref
%
\endisadelimmlref
%
\isatagmlref
%
\begin{isamarkuptext}%
\begin{mldecls}
  \indexdef{}{ML type}{Unsynchronized.ref}\verb|type 'a Unsynchronized.ref| \\
  \indexdef{}{ML}{Unsynchronized.ref}\verb|Unsynchronized.ref: 'a -> 'a Unsynchronized.ref| \\
  \indexdef{}{ML}{!}\verb|! : 'a Unsynchronized.ref -> 'a| \\
  \indexdef{}{ML}{:=}\verb|op := : 'a Unsynchronized.ref * 'a -> unit| \\
  \end{mldecls}%
\end{isamarkuptext}%
\isamarkuptrue%
%
\endisatagmlref
{\isafoldmlref}%
%
\isadelimmlref
%
\endisadelimmlref
%
\begin{isamarkuptext}%
Due to ubiquitous parallelism in Isabelle/ML (see also
  \secref{sec:multi-threading}), the mutable reference cells of
  Standard ML are notorious for causing problems.  In a highly
  parallel system, both correctness \emph{and} performance are easily
  degraded when using mutable data.

  The unwieldy name of \verb|Unsynchronized.ref| for the constructor
  for references in Isabelle/ML emphasizes the inconveniences caused by
  mutability.  Existing operations \verb|!|  and \verb|op :=| are
  unchanged, but should be used with special precautions, say in a
  strictly local situation that is guaranteed to be restricted to
  sequential evaluation --- now and in the future.%
\end{isamarkuptext}%
\isamarkuptrue%
%
\isamarkupsection{Thread-safe programming \label{sec:multi-threading}%
}
\isamarkuptrue%
%
\begin{isamarkuptext}%
Multi-threaded execution has become an everyday reality in
  Isabelle since Poly/ML 5.2.1 and Isabelle2008.  Isabelle/ML provides
  implicit and explicit parallelism by default, and there is no way
  for user-space tools to ``opt out''.  ML programs that are purely
  functional, output messages only via the official channels
  (\secref{sec:message-channels}), and do not intercept interrupts
  (\secref{sec:exceptions}) can participate in the multi-threaded
  environment immediately without further ado.

  More ambitious tools with more fine-grained interaction with the
  environment need to observe the principles explained below.%
\end{isamarkuptext}%
\isamarkuptrue%
%
\isamarkupsubsection{Multi-threading with shared memory%
}
\isamarkuptrue%
%
\begin{isamarkuptext}%
Multiple threads help to organize advanced operations of the
  system, such as real-time conditions on command transactions,
  sub-components with explicit communication, general asynchronous
  interaction etc.  Moreover, parallel evaluation is a prerequisite to
  make adequate use of the CPU resources that are available on
  multi-core systems.\footnote{Multi-core computing does not mean that
  there are ``spare cycles'' to be wasted.  It means that the
  continued exponential speedup of CPU performance due to ``Moore's
  Law'' follows different rules: clock frequency has reached its peak
  around 2005, and applications need to be parallelized in order to
  avoid a perceived loss of performance.  See also
  \cite{Sutter:2005}.}

  Isabelle/Isar exploits the inherent structure of theories and proofs
  to support \emph{implicit parallelism} to a large extent.  LCF-style
  theorem provides almost ideal conditions for that, see also
  \cite{Wenzel:2009}.  This means, significant parts of theory and
  proof checking is parallelized by default.  A maximum speedup-factor
  of 3.0 on 4 cores and 5.0 on 8 cores can be
  expected.\footnote{Further scalability is limited due to garbage
  collection, which is still sequential in Poly/ML 5.2/5.3/5.4.  It
  helps to provide initial heap space generously, using the
  \texttt{-H} option.  Initial heap size needs to be scaled-up
  together with the number of CPU cores: approximately 1--2\,GB per
  core..}

  \medskip ML threads lack the memory protection of separate
  processes, and operate concurrently on shared heap memory.  This has
  the advantage that results of independent computations are directly
  available to other threads: abstract values can be passed without
  copying or awkward serialization that is typically required for
  separate processes.

  To make shared-memory multi-threading work robustly and efficiently,
  some programming guidelines need to be observed.  While the ML
  system is responsible to maintain basic integrity of the
  representation of ML values in memory, the application programmer
  needs to ensure that multi-threaded execution does not break the
  intended semantics.

  \begin{warn}
  To participate in implicit parallelism, tools need to be
  thread-safe.  A single ill-behaved tool can affect the stability and
  performance of the whole system.
  \end{warn}

  Apart from observing the principles of thread-safeness passively,
  advanced tools may also exploit parallelism actively, e.g.\ by using
  ``future values'' (\secref{sec:futures}) or the more basic library
  functions for parallel list operations (\secref{sec:parlist}).

  \begin{warn}
  Parallel computing resources are managed centrally by the
  Isabelle/ML infrastructure.  User programs must not fork their own
  ML threads to perform computations.
  \end{warn}%
\end{isamarkuptext}%
\isamarkuptrue%
%
\isamarkupsubsection{Critical shared resources%
}
\isamarkuptrue%
%
\begin{isamarkuptext}%
Thread-safeness is mainly concerned about concurrent
  read/write access to shared resources, which are outside the purely
  functional world of ML.  This covers the following in particular.

  \begin{itemize}

  \item Global references (or arrays), i.e.\ mutable memory cells that
  persist over several invocations of associated
  operations.\footnote{This is independent of the visibility of such
  mutable values in the toplevel scope.}

  \item Global state of the running Isabelle/ML process, i.e.\ raw I/O
  channels, environment variables, current working directory.

  \item Writable resources in the file-system that are shared among
  different threads or external processes.

  \end{itemize}

  Isabelle/ML provides various mechanisms to avoid critical shared
  resources in most situations.  As last resort there are some
  mechanisms for explicit synchronization.  The following guidelines
  help to make Isabelle/ML programs work smoothly in a concurrent
  environment.

  \begin{itemize}

  \item Avoid global references altogether.  Isabelle/Isar maintains a
  uniform context that incorporates arbitrary data declared by user
  programs (\secref{sec:context-data}).  This context is passed as
  plain value and user tools can get/map their own data in a purely
  functional manner.  Configuration options within the context
  (\secref{sec:config-options}) provide simple drop-in replacements
  for historic reference variables.

  \item Keep components with local state information re-entrant.
  Instead of poking initial values into (private) global references, a
  new state record can be created on each invocation, and passed
  through any auxiliary functions of the component.  The state record
  may well contain mutable references, without requiring any special
  synchronizations, as long as each invocation gets its own copy.

  \item Avoid raw output on \isa{stdout} or \isa{stderr}.  The
  Poly/ML library is thread-safe for each individual output operation,
  but the ordering of parallel invocations is arbitrary.  This means
  raw output will appear on some system console with unpredictable
  interleaving of atomic chunks.

  Note that this does not affect regular message output channels
  (\secref{sec:message-channels}).  An official message is associated
  with the command transaction from where it originates, independently
  of other transactions.  This means each running Isar command has
  effectively its own set of message channels, and interleaving can
  only happen when commands use parallelism internally (and only at
  message boundaries).

  \item Treat environment variables and the current working directory
  of the running process as strictly read-only.

  \item Restrict writing to the file-system to unique temporary files.
  Isabelle already provides a temporary directory that is unique for
  the running process, and there is a centralized source of unique
  serial numbers in Isabelle/ML.  Thus temporary files that are passed
  to to some external process will be always disjoint, and thus
  thread-safe.

  \end{itemize}%
\end{isamarkuptext}%
\isamarkuptrue%
%
\isadelimmlref
%
\endisadelimmlref
%
\isatagmlref
%
\begin{isamarkuptext}%
\begin{mldecls}
  \indexdef{}{ML}{File.tmp\_path}\verb|File.tmp_path: Path.T -> Path.T| \\
  \indexdef{}{ML}{serial\_string}\verb|serial_string: unit -> string| \\
  \end{mldecls}

  \begin{description}

  \item \verb|File.tmp_path|~\isa{path} relocates the base
  component of \isa{path} into the unique temporary directory of
  the running Isabelle/ML process.

  \item \verb|serial_string|~\isa{{\isacharparenleft}{\isacharparenright}} creates a new serial number
  that is unique over the runtime of the Isabelle/ML process.

  \end{description}%
\end{isamarkuptext}%
\isamarkuptrue%
%
\endisatagmlref
{\isafoldmlref}%
%
\isadelimmlref
%
\endisadelimmlref
%
\isadelimmlex
%
\endisadelimmlex
%
\isatagmlex
%
\begin{isamarkuptext}%
The following example shows how to create unique
  temporary file names.%
\end{isamarkuptext}%
\isamarkuptrue%
%
\endisatagmlex
{\isafoldmlex}%
%
\isadelimmlex
%
\endisadelimmlex
%
\isadelimML
%
\endisadelimML
%
\isatagML
\isacommand{ML}\isamarkupfalse%
\ {\isacharverbatimopen}\isanewline
\ \ val\ tmp{\isadigit{1}}\ {\isacharequal}\ File{\isachardot}tmp{\isacharunderscore}path\ {\isacharparenleft}Path{\isachardot}basic\ {\isacharparenleft}{\isachardoublequote}foo{\isachardoublequote}\ {\isacharcircum}\ serial{\isacharunderscore}string\ {\isacharparenleft}{\isacharparenright}{\isacharparenright}{\isacharparenright}{\isacharsemicolon}\isanewline
\ \ val\ tmp{\isadigit{2}}\ {\isacharequal}\ File{\isachardot}tmp{\isacharunderscore}path\ {\isacharparenleft}Path{\isachardot}basic\ {\isacharparenleft}{\isachardoublequote}foo{\isachardoublequote}\ {\isacharcircum}\ serial{\isacharunderscore}string\ {\isacharparenleft}{\isacharparenright}{\isacharparenright}{\isacharparenright}{\isacharsemicolon}\isanewline
\ \ %
\isaantiq
assert%
\endisaantiq
\ {\isacharparenleft}tmp{\isadigit{1}}\ {\isacharless}{\isachargreater}\ tmp{\isadigit{2}}{\isacharparenright}{\isacharsemicolon}\isanewline
{\isacharverbatimclose}%
\endisatagML
{\isafoldML}%
%
\isadelimML
%
\endisadelimML
%
\isamarkupsubsection{Explicit synchronization%
}
\isamarkuptrue%
%
\begin{isamarkuptext}%
Isabelle/ML also provides some explicit synchronization
  mechanisms, for the rare situations where mutable shared resources
  are really required.  These are based on the synchronizations
  primitives of Poly/ML, which have been adapted to the specific
  assumptions of the concurrent Isabelle/ML environment.  User code
  must not use the Poly/ML primitives directly!

  \medskip The most basic synchronization concept is a single
  \emph{critical section} (also called ``monitor'' in the literature).
  A thread that enters the critical section prevents all other threads
  from doing the same.  A thread that is already within the critical
  section may re-enter it in an idempotent manner.

  Such centralized locking is convenient, because it prevents
  deadlocks by construction.

  \medskip More fine-grained locking works via \emph{synchronized
  variables}.  An explicit state component is associated with
  mechanisms for locking and signaling.  There are operations to
  await a condition, change the state, and signal the change to all
  other waiting threads.

  Here the synchronized access to the state variable is \emph{not}
  re-entrant: direct or indirect nesting within the same thread will
  cause a deadlock!%
\end{isamarkuptext}%
\isamarkuptrue%
%
\isadelimmlref
%
\endisadelimmlref
%
\isatagmlref
%
\begin{isamarkuptext}%
\begin{mldecls}
  \indexdef{}{ML}{NAMED\_CRITICAL}\verb|NAMED_CRITICAL: string -> (unit -> 'a) -> 'a| \\
  \indexdef{}{ML}{CRITICAL}\verb|CRITICAL: (unit -> 'a) -> 'a| \\
  \end{mldecls}
  \begin{mldecls}
  \indexdef{}{ML type}{Synchronized.var}\verb|type 'a Synchronized.var| \\
  \indexdef{}{ML}{Synchronized.var}\verb|Synchronized.var: string -> 'a -> 'a Synchronized.var| \\
  \indexdef{}{ML}{Synchronized.guarded\_access}\verb|Synchronized.guarded_access: 'a Synchronized.var ->|\isasep\isanewline%
\verb|  ('a -> ('b * 'a) option) -> 'b| \\
  \end{mldecls}

  \begin{description}

  \item \verb|NAMED_CRITICAL|~\isa{name\ e} evaluates \isa{e\ {\isacharparenleft}{\isacharparenright}}
  within the central critical section of Isabelle/ML.  No other thread
  may do so at the same time, but non-critical parallel execution will
  continue.  The \isa{name} argument is used for tracing and might
  help to spot sources of congestion.

  Entering the critical section without contention is very fast, and
  several basic system operations do so frequently.  Each thread
  should stay within the critical section quickly only very briefly,
  otherwise parallel performance may degrade.

  \item \verb|CRITICAL| is the same as \verb|NAMED_CRITICAL| with empty
  name argument.

  \item Type \verb|'a Synchronized.var| represents synchronized
  variables with state of type \verb|'a|.

  \item \verb|Synchronized.var|~\isa{name\ x} creates a synchronized
  variable that is initialized with value \isa{x}.  The \isa{name} is used for tracing.

  \item \verb|Synchronized.guarded_access|~\isa{var\ f} lets the
  function \isa{f} operate within a critical section on the state
  \isa{x} as follows: if \isa{f\ x} produces \verb|NONE|, it
  continues to wait on the internal condition variable, expecting that
  some other thread will eventually change the content in a suitable
  manner; if \isa{f\ x} produces \verb|SOME|~\isa{{\isacharparenleft}y{\isacharcomma}\ x{\isacharprime}{\isacharparenright}} it is
  satisfied and assigns the new state value \isa{x{\isacharprime}}, broadcasts a
  signal to all waiting threads on the associated condition variable,
  and returns the result \isa{y}.

  \end{description}

  There are some further variants of the \verb|Synchronized.guarded_access| combinator, see \hyperlink{file.~~/src/Pure/Concurrent/synchronized.ML}{\mbox{\isa{\isatt{{\isachartilde}{\isachartilde}{\isacharslash}src{\isacharslash}Pure{\isacharslash}Concurrent{\isacharslash}synchronized{\isachardot}ML}}}} for details.%
\end{isamarkuptext}%
\isamarkuptrue%
%
\endisatagmlref
{\isafoldmlref}%
%
\isadelimmlref
%
\endisadelimmlref
%
\isadelimmlex
%
\endisadelimmlex
%
\isatagmlex
%
\begin{isamarkuptext}%
The following example implements a counter that produces
  positive integers that are unique over the runtime of the Isabelle
  process:%
\end{isamarkuptext}%
\isamarkuptrue%
%
\endisatagmlex
{\isafoldmlex}%
%
\isadelimmlex
%
\endisadelimmlex
%
\isadelimML
%
\endisadelimML
%
\isatagML
\isacommand{ML}\isamarkupfalse%
\ {\isacharverbatimopen}\isanewline
\ \ local\isanewline
\ \ \ \ val\ counter\ {\isacharequal}\ Synchronized{\isachardot}var\ {\isachardoublequote}counter{\isachardoublequote}\ {\isadigit{0}}{\isacharsemicolon}\isanewline
\ \ in\isanewline
\ \ \ \ fun\ next\ {\isacharparenleft}{\isacharparenright}\ {\isacharequal}\isanewline
\ \ \ \ \ \ Synchronized{\isachardot}guarded{\isacharunderscore}access\ counter\isanewline
\ \ \ \ \ \ \ \ {\isacharparenleft}fn\ i\ {\isacharequal}{\isachargreater}\isanewline
\ \ \ \ \ \ \ \ \ \ let\ val\ j\ {\isacharequal}\ i\ {\isacharplus}\ {\isadigit{1}}\isanewline
\ \ \ \ \ \ \ \ \ \ in\ SOME\ {\isacharparenleft}j{\isacharcomma}\ j{\isacharparenright}\ end{\isacharparenright}{\isacharsemicolon}\isanewline
\ \ end{\isacharsemicolon}\isanewline
{\isacharverbatimclose}\isanewline
\isanewline
\isacommand{ML}\isamarkupfalse%
\ {\isacharverbatimopen}\isanewline
\ \ val\ a\ {\isacharequal}\ next\ {\isacharparenleft}{\isacharparenright}{\isacharsemicolon}\isanewline
\ \ val\ b\ {\isacharequal}\ next\ {\isacharparenleft}{\isacharparenright}{\isacharsemicolon}\isanewline
\ \ %
\isaantiq
assert%
\endisaantiq
\ {\isacharparenleft}a\ {\isacharless}{\isachargreater}\ b{\isacharparenright}{\isacharsemicolon}\isanewline
{\isacharverbatimclose}%
\endisatagML
{\isafoldML}%
%
\isadelimML
%
\endisadelimML
%
\begin{isamarkuptext}%
\medskip See \hyperlink{file.~~/src/Pure/Concurrent/mailbox.ML}{\mbox{\isa{\isatt{{\isachartilde}{\isachartilde}{\isacharslash}src{\isacharslash}Pure{\isacharslash}Concurrent{\isacharslash}mailbox{\isachardot}ML}}}} how
  to implement a mailbox as synchronized variable over a purely
  functional queue.%
\end{isamarkuptext}%
\isamarkuptrue%
%
\isadelimtheory
%
\endisadelimtheory
%
\isatagtheory
\isacommand{end}\isamarkupfalse%
%
\endisatagtheory
{\isafoldtheory}%
%
\isadelimtheory
%
\endisadelimtheory
\end{isabellebody}%
%%% Local Variables:
%%% mode: latex
%%% TeX-master: "root"
%%% End:


\begingroup
\tocentry{\bibname}
\bibliographystyle{plain} \small\raggedright\frenchspacing
\bibliography{../manual}
\endgroup

\tocentry{\glossaryname}
\printglossary

\tocentry{\indexname}
\printindex

\end{document}


%%% Local Variables: 
%%% mode: latex
%%% TeX-master: t
%%% End: 
